\documentclass[11pt]{article}

\usepackage[T1]{fontenc}
\usepackage[utf8]{inputenc}
\usepackage{lmodern}
\usepackage[margin=1in]{geometry}
\usepackage{amsmath,amssymb,amsthm,mathtools}
\usepackage{enumitem}
\usepackage{booktabs}
\usepackage{array}
\usepackage{longtable}
\usepackage{xcolor}
\usepackage{microtype}
\usepackage{etoolbox}
\usepackage{url}
\usepackage[colorlinks=true,linkcolor=blue!50!black,citecolor=green!35!black,urlcolor=blue!60!black]{hyperref}
\usepackage[capitalize,nameinlink]{cleveref}

% ---------------------------------------------------------------------------
% Theorem environments
% ---------------------------------------------------------------------------
\theoremstyle{plain}
\newtheorem{theorem}{Theorem}[section]
\newtheorem{lemma}[theorem]{Lemma}
\newtheorem{corollary}[theorem]{Corollary}
\newtheorem{proposition}[theorem]{Proposition}
\newtheorem{conjecture}[theorem]{Conjecture}
\newtheorem{question}[theorem]{Question}
\newtheorem{problem}[theorem]{Problem}
\theoremstyle{definition}
\newtheorem{definition}[theorem]{Definition}
\newtheorem{remark}[theorem]{Remark}
\newtheorem{example}[theorem]{Example}

% The Lean proof is the proof; the text gives a faithful outline of it.
\newenvironment{outline}[1][Outline of the formal proof]{\begin{proof}[#1]}{\end{proof}}
% Results that are not formalized get complete proofs, labelled as such.
\newenvironment{informalproof}{\begin{proof}[Proof (not formalized)]}{\end{proof}}

% ---------------------------------------------------------------------------
% Lean references
%   \lean{name}            a declaration, relative to the namespace DirectedReciprocity
%                          (declarations of earlier papers are written with their full names)
%   \formal{file}{names}   the formal counterpart of a numbered result: a file relative to
%                          lean/DirectedReciprocity/ and a comma-separated list of declarations
%   \partformal{file}{names}{gap}
%                          a numbered result whose formal counterpart covers all but the
%                          part described in {gap}; the gap is proved or cited in the text
%   \notformalized         a numbered result that has no formal counterpart
%   \cited{where}          a result of an earlier paper used without proof: where it is proved,
%                          and whether it is Lean-verified there
% ---------------------------------------------------------------------------
\DeclareUrlCommand\leanname{\def\UrlFont{\ttfamily\upshape}\def\UrlBreaks{\do\_\do\.\do\/\do\-}}
\renewcommand{\UrlBreaks}{\do\_\do\.\do\/\do\-}
\emergencystretch=2.5em
\newcommand{\lean}[1]{\relax\ifmmode\text{\leanname{#1}}\else\leanname{#1}\fi}
\newcommand{\leanfile}[1]{\leanname{#1}}
\newcommand{\leanlist}[1]{\begingroup\def\leansep{}%
  \renewcommand*{\do}[1]{\leansep\leanname{##1}\def\leansep{, }}\docsvlist{#1}\endgroup}
\newcommand{\formal}[2]{\par\noindent{\small\upshape\raggedright\textsf{Lean:} \leanlist{#2}\unskip
  \nobreak\hfil\penalty50\hskip1em\hbox{}\nobreak\hfill[\leanfile{#1}]\parfillskip=0pt\par}%
  \smallskip\noindent\ignorespaces}
\newcommand{\partformal}[3]{\par\noindent{\small\upshape\raggedright\textsf{Lean (partial):} \leanlist{#2}\unskip
  \nobreak\hfil\penalty50\hskip1em\hbox{}\nobreak\hfill[\leanfile{#1}]\parfillskip=0pt\par}%
  \par\noindent{\small\upshape\textsf{Not formalized:} #3}\par\smallskip\noindent\ignorespaces}
\newcommand{\notformalized}{\par\noindent{\small\upshape\textsf{Not formalized;} a complete proof is given
  below.}\par\smallskip\noindent\ignorespaces}
\newcommand{\cited}[1]{\par\noindent{\small\upshape\textsf{Cited:} #1}\par\smallskip\noindent\ignorespaces}

% ---------------------------------------------------------------------------
% Notation
% ---------------------------------------------------------------------------
\newcommand{\N}{\mathbb{N}}
\newcommand{\Z}{\mathbb{Z}}
\newcommand{\Q}{\mathbb{Q}}
\newcommand{\R}{\mathbb{R}}
\newcommand{\C}{\mathbb{C}}
\newcommand{\F}{\mathbb{F}}
\newcommand{\T}{\mathsf{T}}
\newcommand{\clog}[1]{\lceil \log_2 #1 \rceil}
\newcommand{\abs}[1]{\lvert #1 \rvert}
\newcommand{\card}[1]{\lvert #1 \rvert}
\newcommand{\eps}{\varepsilon}
\newcommand{\prev}{\mathrm{prev}}
\newcommand{\cost}{\mathrm{cost}}
\newcommand{\desc}{\mathrm{desc}}
\newcommand{\gn}{\mathrm{gn}}
\newcommand{\Sh}{\mathrm{Sh}}
\newcommand{\cl}{\mathrm{cl}}
\newcommand{\Aut}{\mathrm{Aut}}
\newcommand{\Crux}{\mathrm{Crux}}
\newcommand{\TCP}{\mathrm{TCP}}
\DeclareMathOperator{\rank}{rank}
\DeclareMathOperator{\dist}{dist}
\DeclareMathOperator{\poly}{poly}
\DeclareMathOperator{\lev}{lev}

% ---------------------------------------------------------------------------
% Paper-specific notation
% ---------------------------------------------------------------------------
\newcommand{\Cay}{\operatorname{Cay}}
\newcommand{\GP}{\operatorname{GP}}
\newcommand{\ClassS}{\mathcal{S}}
\newcommand{\sigb}{\mathrm{b}}
\newcommand{\sigf}{\mathrm{f}}
\newcommand{\rhom}{\rho_{\mathrm{m}}}
\newcommand{\Ex}{\mathrm{X}_{15}}
\newcommand{\Gto}{\mathrm{G}_{21}}
\newcommand{\DNP}{\mathrm{DNP}}


\title{Directed Reciprocity and Deep-Nesting Bounds\\ for Matrix Transposition on Tapes}
\author{Hans Brende \and Gemini 4 Argon (Google DeepMind)}
\date{October 2026}

\begin{document}

\maketitle

\begin{abstract}
Can a deterministic multitape Turing machine transpose a $d \times d$ matrix with $\max(1, \clog{d})$-bit entries in time
linear in the input length? By the fourth paper of this series, excluding oblivious linear-time transposers splits into two
combinatorial problems: a half-rate bound for storage codes on transitive-tournament-free recovery relations (the coding
crux), and cheap depth reduction of head-walk dependency graphs (the trajectory cut problem). This paper advances three
frontiers left open by the seventh, eighth and ninth papers of the series, and every new numbered theorem, proposition,
lemma and corollary is kernel-verified in Lean~4.

First, for nested elimination on directed recovery relations, we give a purely positional characterization of the
nestedness hypothesis of the seventh paper, show that every reciprocal pair forces distinct record lists, and prove a
sharp dichotomy for chordless digon cycles with at most one active outside parent per vertex: a piece certificate exists
over some number of records if and only if the outside arcs do not form a unidirectional cycle of stride~$3$
(Theorem~DC). In particular, the $15$-vertex obstruction $\Ex$ has no piece certificate on its unique minimum nonempty
core over any number of records, and the Petersen graph $\GP(5,2)$ has no global elimination certificate over any number
of records, whereas $\GP(7,2)$ and $\GP(9,2)$ require minimum record counts $4$ and~$3$, respectively (Proposition~GP),
resolving two open formalization problems of the seventh paper. Using minimal chordal digon paths and hub-augmented
stride-$3$ pieces, we then prove that the entire class $\ClassS$ of Theorem~S of the fourth paper---transitive-tournament-free
relations of maximum in-degree~$3$ in which every vertex of in-degree~$3$ has only reciprocal children---admits a
certified piece decomposition and satisfies the half-rate bound $\card{C}^2 \le \card{\alpha}^{\card{V}}$ over every
finite alphabet and for all codes, linear or nonlinear (Theorem~S$^*$), resolving the all-code Theorem~S target of the
seventh paper, together with a directed barrier inequality (Corollary~B$^*$) and a $21$-vertex strongly connected
triangle-free witness $\Gto \in \ClassS$ on which earlier peeling rules stall (Proposition~G21).

Second, for scalar linear codes on triangle-free graphs of maximum degree~$4$ defined as the kernel of a matrix $M$ with
unit diagonal over an arbitrary field, we remove the heavy-separation hypothesis~(S2) of the eighth paper's Theorem~SH:
adjacent degree-$4$ exceptional vertices and arbitrary matchings of exceptional vertices are allowed, provided no two
exceptional vertices share a common neighbor (Theorem~SH$^*$), yielding a matching bound controlled by the heavy-matching
invariant $\rhom(G,H) \ge \rho_3(G,H)$ (Corollary~M$^*$).

Third, on the deep-nesting walk $U(s_0, m, 1)$, we formalize the recursive stroke structure and $O(m)$ turn-potential
evaluator (Proposition~SL), prove an $m$-free mod-$q$ straddle lower bound and exact parity-conditioned straddle count
$\Pi_\kappa^2 / 4$ (Theorem~V$^*$), establish a super-Fibonacci unwrapping theorem showing that level-uniform single-head
cancellation admits at most $2\lfloor \log_2 q \rfloor + 2$ consecutive levels whose blind apex and band-$0,1,2$ errors are
at most $q/8$ (Theorem~LU), determine the block-cancellation drift and Mersenne $q = 2^B - 1$ binary-clock failure residue
(Proposition~OC), and prove an unconditional three-regime witness theorem with zero $\log q$ loss for block-label pairs of
bounded block length across every $K_1, K_2 \in \Z$ and every modulus $q \ge 3$ at two fixed odd shifts, or a single shift
for symmetric pairs (Theorem~A; translating its single-head regimes to actual two-head hop errors requires blindness at the
chosen shift). The linear-time transposition question itself remains open.

\end{abstract}

\begin{quote}\small
\textbf{Contributions.} Hans Brende chose the problem, fixed the verification standard, and directed the research.
Gemini 4 Argon, a large language model developed by Google DeepMind, developed the mathematics, wrote the Lean
formalization, and wrote this paper under Hans Brende's direction. \Cref{sec:reciprocity} builds on the nested
elimination engine and piece-decomposition framework of \emph{Nested Elimination and Half-Rate Bounds for Storage Codes}
by Hans Brende and Claude Opus 5.5 (Anthropic)~\cite{SeventhPaper2026} and on the class of directed recovery relations
introduced in \emph{Curved Motion, Depth Reduction, and Storage Codes for Matrix Transposition on Tapes} by the same
authors~\cite{FourthPaper2026}. \Cref{sec:crux,sec:dnp} extend the minimal-counterexample method for degree-$4$ scalar
linear codes and the deep-nesting walk and block-label analysis of \emph{Bounded Nesting and Depth Reduction for Matrix
Transposition on Tapes}, also by Hans Brende and Claude Opus 5.5 (Anthropic)~\cite{EighthPaper2026}, and situate these
results alongside the source-cover barriers of \emph{Source Covers and Barriers to Tape Transposition} by Hans Brende and
GPT-6 Astra (OpenAI)~\cite{NinthPaper2026}.
\end{quote}

\tableofcontents
\clearpage

\section{Introduction}\label{sec:intro}

Does there exist a deterministic multitape Turing machine that transposes every $d \times d$ matrix of $b(d)$-bit entries,
where $b(d) = \max(1, \clog{d})$ and $L = d^2 b(d)$ is the bit-length of the input, in time $O(L)$? Following the earlier
papers of this series~\cite{FirstPaper2026}, we write~(T) for this assertion:
\begin{quote}
(T)\quad There exists a deterministic multitape Turing machine---with a binary read-only two-way input tape, a binary
write-only one-way output tape, and a fixed finite number of two-way read/write work tapes over a fixed finite work
alphabet---that transposes $d \times d$ matrices of $b(d)$-bit entries in $O(d^2 b(d))$ steps.
\end{quote}
By a theorem of Harvey and van der Hoeven~\cite{HvdH2025}, $N$-bit integer multiplication reduces in linear time on a
multitape Turing machine to matrix transposition; hence refuting~(T) would establish an $\Omega(N \log N)$ lower bound
for deterministic off-line multitape integer multiplication. Whether~(T) holds remains open.

\subsection{Context in the series}

The fourth paper of this series~\cite{FourthPaper2026} reduced the impossibility of \emph{oblivious} linear-time
transposers to the conjunction of two combinatorial statements that make no mention of Turing machine states:
\begin{enumerate}[label=(\roman*)]
\item \emph{The trajectory cut problem} ($\TCP(k)$): in every $k$-head walk dependency DAG of length $T$, deleting at
  most $\eps T$ edges reduces the longest directed path to length $O_{\eps,k}(1)$.
\item \emph{The coding crux}: every storage code over a finite alphabet on a transitive-tournament-free (TT-free)
  directed recovery relation of bounded degree has asymptotic rate bounded away from~$1$---and, in its sharpest form
  (Conjecture~P of~\cite{SeventhPaper2026}), rate at most~$1/2$ whenever every vertex has in-degree at most~$3$.
\end{enumerate}
Subsequent papers sharpened both halves of this program while exposing concrete obstructions to naive extensions:
\begin{itemize}
\item On the \emph{directed coding} side, the fourth paper~\cite{FourthPaper2026} proved the half-rate bound for
  \emph{scalar linear} codes on the mixed class $\ClassS$ of TT-free relations with maximum in-degree~$3$ in which every
  vertex of in-degree~$3$ has only reciprocal children (its Theorem~S). The seventh paper~\cite{SeventhPaper2026}
  introduced the nested elimination engine over general polymatroids (removing linearity altogether), proved the
  all-code half-rate bound on all undirected subcubic triangle-free graphs (its Theorem~U3), and constructed the
  $15$-vertex TT-free orientation $\Ex$ of maximum in-degree~$3$ on which minimum-core peeling fails. Within item~2
  of~\cite[Section~9.2]{SeventhPaper2026}, alongside Conjecture~P (\cite[Conjecture~9.1]{SeventhPaper2026}, for general
  in-degree at most~$3$), it singled out as an intermediate target whether the fourth paper's Theorem~S lifts from scalar
  linear codes to \emph{all} codes on the entire class~$\ClassS$, and noted two unformalized impossibility
  claims~\cite[Remarks~7.3 and~7.5; Section~9.2, item~6]{SeventhPaper2026}: that $\Ex$ admits no piece certificate on its
  unique minimum nonempty core over any number of records, and that the Petersen graph admits no global elimination
  certificate over any number of records.
\item On the \emph{degree-$4$ linear coding} side, while at degree~$5$ scalar binary storage codes already exceed
  rate~$1/2$ on the Clebsch graph (rate~$5/8$, Lean-verified in~\cite[Corollary~6.11]{FourthPaper2026}), whether any
  storage code on a triangle-free graph of maximum degree~$4$ can exceed rate~$1/2$ remains open
  (\cite[Section~9.2, item~4]{SeventhPaper2026}; \cite[Section~10, Problem~10.4]{EighthPaper2026}), and tree-local
  arguments cannot handle arbitrary clusters of degree-$4$ vertices~\cite[Proposition~6.28]{FourthPaper2026}. The eighth
  paper~\cite{EighthPaper2026} isolated the exceptional set $H$ of degree-$4$ rows in a unit-diagonal check matrix
  $M = I + B$ on a triangle-free graph of maximum degree~$4$ and proved $\dim\ker M \le \card{V}/2$ over every field
  under two hypotheses on $H$: (S1)~no two vertices of $H$ share a common neighbor, and (S2)~no two vertices of $H$ are
  adjacent. Whether~(S2) is genuinely needed or an artifact of the induction was left
  open~\cite[Section~8.2, Remark~8.8; Section~10, Problem~10.4]{EighthPaper2026}.
\item On the \emph{head-motion and deep-nesting} side, the eighth paper~\cite{EighthPaper2026} proved $\TCP(k)$ for all
  schedules of bounded nesting load (its Theorem~F) and isolated the $2$-head \emph{deep-nesting pair} $\DNP(s_0, m)$,
  running two shifted copies of the recursive nested walk $U(s_0, m, 1)$, as the sharpest open case of $\TCP(2)$. On a
  single copy of $U(s_0, m, 1)$, the eighth paper's turn counter~\cite[Section~5.2, Theorem~5.2 and Corollary~5.5]{EighthPaper2026}
  (built on Lemma~Zq of~\cite[Section~5.5, Lemma~5.15]{FourthPaper2026}) cuts depth to $2^{m + O(1)}/\eps$; the eighth
  paper introduced block-cancellation labels $B(a, K)$ to probe whether two heads can cancel turn-counter drift across
  scales, stated Conjecture~A$''$~\cite[Section~5.6, Conjecture~5.24]{EighthPaper2026}---that at some odd shift, every
  mod-$q$ label of cost at most $n_{\mathrm{h}}/8$ on the $r$-fold cover has at most $C\log_2 q + C' + \eta(m+1)$ doubly
  $\theta_0$-served levels ($\eta < 1$)---and at the end of~\cite[Section~5.6]{EighthPaper2026} identified two open steps
  for proving Conjecture~A$''$ even on block-label pairs: (1)~\emph{``a bound on $\abs{F}$ on the straddling hops, or a
  version of Theorem~V that counts the window directly''}, and (2)~\emph{``at one fixed shift for all coarse classes,
  equidistribution of a class's apex centres relative to the other head's canceller valleys of level about
  $k + \log_2(m - k)$, a mixed-radix carry problem''}. Meanwhile, the ninth
  paper~\cite{NinthPaper2026} developed an independent bypass of $\TCP$ via complete source covers and return-profile
  barriers on restricted return geometries, leaving open the unrestricted adaptive regime.
\end{itemize}

\subsection{Summary of main results}

This paper resolves the intermediate directed coding target of~\cite[Section~9.2, item~2]{SeventhPaper2026} (while
Conjecture~P, its Conjecture~9.1, remains open), removes hypothesis~(S2) from Theorem~SH of~\cite{EighthPaper2026}, and
proves and formalizes structural stroke, straddle, and super-Fibonacci laws together with three-regime failure witnesses for
bounded-range block-label pairs on the deep-nesting walk. Every new numbered result in the paper is verified in Lean~4 using
only the three standard foundational axioms (\cref{sec:prelim-lean,app:verify}).

\paragraph{1.\ Directed reciprocity, digon-cycle dichotomy, and the all-code Theorem~S$^*$ (\cref{sec:reciprocity}).}
We begin by recasting the stateful recursive predicate \lean{HalfRate.RelCharged} of~\cite{SeventhPaper2026} in purely
positional terms (\cref{prop:PN}): an elimination list $l = [v_0, \dots, v_{k-1}]$ with records $\{P_r\}_{r \in I}$,
supplying-record map $b$, and optional outside-parent map $\mathrm{par}$ is relationally nested from empty initial state
if and only if every step covers its active parents from the revealed prefix $S_j$, its supplying record $P_{b(v_j)}$,
and its optional outside parent, and whenever $j < j'$ share the same outside parent $w$, $P_{b(v_j)} \subseteq S_{j'}
\cup P_{b(v_{j'})}$ (for disjoint records, either $b(v_j) = b(v_{j'})$ or all of $P_{b(v_j)}$ is revealed before step
$j'$). As an immediate invariant on any piece certificate, no two vertices joined by a reciprocal pair (a \emph{digon})
can belong to the same own record $P_{a(u)}$, so every connected digon component on at least two vertices requires
$m \ge 2$ records.

We then classify piece certificates on chordless digon cycles $C_L$ ($L \ge 4$) in TT-free relations where each cycle
vertex has at most one active parent outside $U \cup C_L$ (\cref{prop:CS,thm:DC}): an $m$-record piece certificate
exists on $C_L$ for \emph{some} $m \ge 1$ if and only if the outside parents do not form a \emph{stride-$3$ hub triple},
that is, $3 \mid L$ together with outside vertices $w_0, w_1, w_2 \notin U \cup C_L$ such that $w_{i \bmod 3}$ is an
active parent of cycle vertex $i$ for every $0 \le i < L$. When the stride-$3$ hub obstruction is absent, $m = 2$
records suffice for $L = 4$ and $m = L + 1$ records suffice via two-sided arc expansion for every $L \ge 5$ (while on
pure stride-$s$ hub cores $C_{sk}$, $m = 2$ suffices for even $s \ge 4$ and $m = 3$ suffices for every $s \ge 4$); when
the stride-$3$ hub triple is present, a positional locking invariant rules out piece certificates across \emph{all}
$m \ge 1$. Specializing to the unique minimum nonempty core $\{0,\dots,5\}$ of $\Ex$ (\cref{cor:x15}) and to generalized
Petersen graphs (\cref{prop:GP}) resolves both parts of~\cite[Remarks~7.3 and~7.5; Section~9.2, Problem~6]{SeventhPaper2026}
in Lean~4 and reveals
a non-monotone global-certificate record hierarchy: $\GP(5,2)$ (the Petersen graph) requires $\infty$ records globally,
$\GP(7,2)$ requires exactly $4$ records, and $\GP(9,2)$ drops back to $3$ records.

Using these cycle tools, we prove \textbf{Theorem~S$^*$} (\cref{thm:Sstar}): every relation $N$ in the full class
$\ClassS$ of Theorem~S of~\cite{FourthPaper2026} admits a certified piece decomposition (\lean{HalfRate.Decomp}), so
every valid polymatroid $h$ on $V$ obeys $2 h(V) \le \sum_{v \in V} h(\{v\})$ and every
storage code $C \subseteq \alpha^V$ (linear or nonlinear, over any finite alphabet $\alpha$) satisfies
$\card{C}^2 \le \card{\alpha}^{\card{V}}$. The key mechanism overcoming embedded stride-$3$ hub cycles inside an
in-closed strongly connected component $S$ of $\ClassS$ is \emph{hub augmentation}
(\lean{exists_piece_of_scc_classS}): if $S$ has no minimal chordal digon arc and a shortest digon cycle $C_L \subseteq S$
is obstructed as a standalone piece, then $C_L$ is a stride-$3$ hub cycle with hubs $w_0, w_1, w_2 \in S \setminus C_L$;
strong connectivity forces an outgoing arc from some $u_{i_0} \in C_L$ to $S \setminus C_L$, which $\ClassS$ reciprocates
into a digon between $u_{i_0}$ and its hub $w_{r_0}$ while girth minimality forces $w_{r_0}$ to have in-degree~$2$ and no
other child on $C_L$, so adjoining $w_{r_0}$ yields an $(L+2)$-record sequential nested piece certificate on
$C_L \cup \{w_{r_0}\}$. We deduce a directed vertex-cut and isolated-outside-vertex barrier inequality for all $\ClassS$
relations (\cref{cor:Bstar}) and construct an explicit $21$-vertex strongly connected triangle-free witness $\Gto \in \ClassS$
(\cref{prop:G21}) on which low-in-degree peeling, one-way cycle peeling, and shortest-digon-cycle rules all fail
simultaneously, while Theorem~S$^*$ certifies rate at most~$1/2$.

\paragraph{2.\ Scalar linear codes at degree~$4$ without heavy separation (\cref{sec:crux}).}
In \cref{sec:crux} we prove \textbf{Theorem~SH$^*$} (\cref{thm:SHstar}): over any field $F$, if $M \in F^{V \times V}$
has unit diagonal, no recovery triangle, and satisfies $(\star 4)$ at every vertex (for instance when $M$ is supported in a
triangle-free simple graph $G = (V, E)$ of maximum degree~$4$), and the heavy set $H_M$ of vertices violating $(\star\star)$
either has size at most~$2$ or satisfies~(S1) alone---no vertex has two distinct heavy recovery neighbors (in graph form,
no two distinct vertices of $H_M$ share a common neighbor in~$G$)---then $\dim_F \ker M \le \lfloor \card{V}/2 \rfloor$
(and likewise under the same hypotheses on $M^\T$). Hypothesis~(S2) of~\cite{EighthPaper2026} is eliminated
completely: $G[H_M]$ may contain arbitrary induced matchings of adjacent degree-$4$ vertices. The proof replaces Phase~E
of~\cite{EighthPaper2026} with a unified final count (\lean{F3_star}): under~(S1) alone, a pending triple of a heavy vertex
$h$ contains at most one heavy member $g$, leaving two light members $u_1, u_2$ whose shared pending pair $\{a, b\}$ and
neutral dimension drop contradict linear independence uniformly for all $\card{W \setminus X_0} \ge 3$, while in minimal
counterexamples any degree-$4$ heavy vertex still retains at least three light neighbors under~(S1). Combining
Theorem~SH$^*$ with Tutte--Berge deficiency yields \textbf{Corollary~M$^*$} (\cref{cor:Mstar}), improving the eighth
paper's distance-$3$ packing invariant $\rho_3(G, H_M)$ to the common-neighbor-free packing invariant
$\rhom(G, H_M) \ge \rho_3(G, H_M)$.

\paragraph{3.\ Deep-nesting stroke laws, super-Fibonacci unwrapping, and bounded-range block pairs (\cref{sec:dnp}).}
In \cref{sec:dnp} we formalize the algebraic and combinatorial theory of turn-supported periodic labels modulo~$q$ on the
deep-nesting walk $U(s_0, m, 1)$ and on the deep-nesting pair $\DNP(s_0, m)$ (thereby also providing Lean~4
counterparts for the unformalized stroke, band, block and straddle lemmas of~\cite[Section~5]{EighthPaper2026}):
\begin{itemize}
\item \textbf{Proposition~SL} (\cref{prop:SL}) establishes the exact stroke decomposition bijection in $[0, \Pi_0)$,
  proves that under fair binary branch choices the stopping depth before the nested walk visits~$m$ has distribution
  $2^{-(s+1)}$ on $\{0,\dots,m-1\}$ independently of directional parity, and constructs an $O(m)$ recursive turn-potential
  evaluator \lean{evalPhiRec} proved equal to cumulative turn jumps over a full period $P = 4 \Pi_0$.
\item \textbf{Theorem~V$^*$} (\cref{thm:Vstar}) proves an $m$-free mod-$q$ straddle involution bound: whenever a level-$k$
  apex hop of one head straddles a level-$\kappa$ V-turn of the other head ($\Pi_\kappa \le R_k$), the paired half-widths
  $(d, \Pi_\kappa - d)$ cancel the nested potential identically over $\Z$, forcing their reduced mod-$q$ errors to sum to
  at least $(2 A' + 2(c_{\sigf}(\kappa) - c_{\sigb}(\kappa))) \bmod q$; moreover, summing the straddle set size over
  centers of either fixed parity across a $2\Pi_\kappa$-window yields the \emph{exact} triangular identity
  $(\Pi_\kappa / 2)^2 = \Pi_\kappa^2 / 4$.
\item \textbf{Theorem~LU} (\cref{thm:LU}) proves that for any level-uniform turn-supported label modulo $q \ge 3$, if
  $s \ge 1$ consecutive levels have single-head reduced apex and valley band-$0,1,2$ errors at most $q/8$, the modular
  band relation unwraps over $\Z$ with zero wrap-around and forces the super-Fibonacci inequality
  $b_0(\ell) \ge b_0(\ell+1) + b_0(\ell+2) + 4$, whence $32(F_s - 1) \le q$ and $s \le 2\lfloor \log_2 q \rfloor + 2$.
\item \textbf{Proposition~OC} (\cref{prop:OC}) proves exact forward and backward block-cancellation bounds for $B(a, K)$,
  establishes strictly positive mean drift $+(a(K+1)+1)/2$ per $(a+1)$-level block across all $2^{a+1}$ stroke-direction
  choices, and shows that the pure binary clock $B(\infty, 0)$ fails modulo every Mersenne number $q = 2^B - 1$ ($B \ge 2$)
  with straddle residue $q - 1$ at every third/periodic scale $B \mid (m - \kappa + 2)$.
\item \textbf{Theorem~A} (\cref{thm:A}) combines these ingredients for block-label pairs $(B(a_1, K_1),\allowbreak B(a_2, K_2))$ of
  bounded block length ($1 \le a_1, a_2$): Theorem~V$^*$ completely bypasses open step~(1) of~\cite[Section~5.6]{EighthPaper2026}
  for level-uniform labels (paired straddle sums need no bound on $\abs{F}$ and no window assumption on~$q$), while
  \lean{exists_oddShift_coarse_straddle} establishes a one-direction $\tfrac{1}{16}$-fraction form of step~(2) (weaker than
  the all-coarse-classes hypothesis of~\cite[Proposition~5.23]{EighthPaper2026}, which is obstructed by prefix lockstep of
  nested V-apices, \lean{apexOffset_diff_indep_prefix}). Whenever $m < s_0$ and
  $4(\lfloor \log_2(m+1)\rfloor + 1 + \max(a_1, a_2)) \le m + 1$, Theorem~A produces two fixed odd shifts $\xi_1, \xi_2$
  (a single odd shift $\xi$ for symmetric pairs $a_1 = a_2, K_1 = K_2$) giving three-regime failure witnesses on
  $\Omega(m)$ levels across \emph{every} $K_1, K_2 \in \Z$ and \emph{every} modulus $q \ge 3$, with \emph{zero} $\log q$
  loss: blind single-head error witnesses in Regimes~1 and~2, and two-head coarse straddle witnesses in Regime~3. We also
  record the four remaining steps separating these witnesses from Conjecture~A$''$ (\cref{rem:scope-A}).
\end{itemize}

\subsection{Related work}

The equivalence between locally recoverable storage capacity on a digraph $G$, index coding broadcast rate on the
complement, and the guessing number $\gn(G) = \sup_{\card{\alpha} \ge 2} \sup_{C} \log_{\card{\alpha}} \card{C}$ (over
locally recoverable codes $C \subseteq \alpha^V$,~\cite{Riis2007,SeventhPaper2026}) goes back to Birk and Kol~\cite{BirkKol2006},
Riis~\cite{Riis2007}, Bar-Yossef, Birk, Jayram and Kol~\cite{BBJK2011}, Gadouleau and Riis~\cite{GadouleauRiis2011}, and
Mazumdar~\cite{Mazumdar2014}. On undirected triangle-free graphs, Christofides and Markstr\"om~\cite{ChristofidesMarkstrom2011}
bounded the guessing number by the fractional matching number (giving rate~$1/2$ for all undirected triangle-free graphs
only on symmetric recovery neighborhoods), Cameron, Dang and Riis~\cite{CameronDangRiis2016} classified high-guessing
triangle-free graphs in low degree, Atkins, Rombach and Skerman~\cite{AtkinsRombachSkerman2017} determined guessing
numbers of odd cycles, and Mazumdar, McGregor and Vorotnikova~\cite{MMV2019} proved half-rate bounds under sparse global
average degree, whereas without a degree bound Barg and Z\'emor~\cite{BargZemor2022} and Barg, Schwartz and
Yohananov~\cite{BSY2024} constructed coset and Cayley storage codes on triangle-free graphs of growing degree with rates
exceeding~$1/2$ (approaching~$3/4$ and~$1$, respectively).
On directed graphs without undirected symmetry, undirected fractional-matching duality breaks down because of directed
cycles of pure one-way arcs; Theorem~S$^*$ (\cref{thm:Sstar}) is, to our knowledge, the first all-code half-rate theorem
covering a class of TT-free digraphs that simultaneously permits full subcubic undirected reciprocity, unbounded
directed one-way cycles, and mixed components combining both.

On the complexity side, reductions from network coding and circuit depth to Valiant's edge-removal problem go back to
Valiant~\cite{Valiant1977}; see~\cite{FirstPaper2026,FourthPaper2026,EighthPaper2026,NinthPaper2026} for a detailed
account of the connections between $\TCP(k)$, monotone and non-monotone depth reduction, matrix rigidity, and source
covers.

\subsection{Organization and formal reading guide}

\Cref{sec:prelim} fixes notation from the earlier papers of the series and explains our Lean~4 verification conventions
(\cref{sec:prelim-lean}). \Cref{sec:reciprocity} develops positional nestedness, the digon-cycle dichotomy, the
generalized Petersen record hierarchy, and Theorem~S$^*$. \Cref{sec:crux} proves Theorem~SH$^*$ and Corollary~M$^*$.
\Cref{sec:dnp} proves Proposition~SL, Theorem~V$^*$, Theorem~LU, Proposition~OC, and Theorem~A on the deep-nesting walk
and pair. \Cref{sec:construction} synthesizes what any hypothetical linear-time transposer must look like in light of
these bounds, and \cref{sec:open} collects the sharpened open problems. \Cref{app:index} indexes every Lean declaration
cited in the paper, and \Cref{app:verify} gives exact instructions for building the library, checking axioms, and
reproducing the finite numerical examples.

\section{Preliminaries}\label{sec:prelim}

We keep the notation and conventions of the fourth, seventh and eighth papers~\cite{FourthPaper2026,SeventhPaper2026,EighthPaper2026}
so that cited declarations align directly with their Lean~4 libraries. Throughout, logarithms are base~$2$, $\N = \{0,1,2,\dots\}$,
and for integers $a \le b$ we write $[a, b] = \{x \in \Z : a \le x \le b\}$ and $[a, b) = \{x \in \Z : a \le x < b\}$ (or
their subsets in $\N$ when the endpoints are natural numbers). For $x \in \Z$ and $q > 0$, we write $x \bmod q \in
\{0, 1, \dots, q-1\}$ for the non-negative Euclidean remainder (\texttt{Int.emod} in Lean).

\subsection{Turing machines, transposition, and the depth--coding split}

In the machine model of~\cite{FirstPaper2026,FourthPaper2026} (\lean{Transpose.TM}), a deterministic multitape Turing
machine has a binary read-only two-way input tape, a binary write-only one-way output tape, and $k$ two-way read/write
work tapes over a finite work alphabet $\{0, 1, \dots, g\}$. Given a $d \times d$ matrix of $b(d)$-bit entries in
row-major order, where $b(d) = \max(1, \clog{d})$ and $L = d^2 b(d)$, the machine must halt with the transpose in
row-major order on the output tape; statement~(T) is formalized as \lean{Transpose.LinearTimeTransposable}
(see~\cite[Section~2]{FirstPaper2026}, Lean-verified). A transposer is \emph{oblivious} if its head trajectories,
output-append times, and halting time depend only on the input dimension~$d$~\cite[Section~2]{FirstPaper2026}. By
Theorem~D of~\cite[Section~4.2, Theorem~4.10]{FourthPaper2026} (proved in the text of~\cite{FourthPaper2026}; not
formalized in Lean), if the $k$-head walk dependency DAG of an oblivious linear-time transposer admits $\eps T$-edge cuts
to bounded hop-depth ($\TCP(k)$) and storage codes on the resulting bounded-in-degree TT-free recovery relations have rate
bounded away from~$1$, then no such transposer exists.

\subsection{Directed recovery relations, polymatroids, and nested elimination}

Let $V$ be a finite vertex set. A \emph{recovery relation} on $V$ is a parent map $N : V \to \mathcal{P}(V)$; we write
$u \to v$ when $u \in N(v)$ ($u$ is a \emph{parent} of $v$, and $v$ is a \emph{child} of $u$). The relation $N$ is
\emph{irreflexive} (loopless) if $v \notin N(v)$ for all $v \in V$. Two distinct vertices $u, v$ form a \emph{reciprocal
pair} or \emph{digon} if $u \in N(v)$ and $v \in N(u)$, and $u \to v$ is a \emph{pure one-way arc} if $u \in N(v)$ and
$v \notin N(u)$. We say $N$ is \emph{transitive-tournament-free} (\lean{HalfRate.TTFree}) if there are no three distinct
vertices $a, b, c \in V$ with $a \in N(b)$, $b \in N(c)$ and $a \in N(c)$; notice that an irreflexive symmetric relation is
TT-free if and only if its underlying undirected simple graph is triangle-free. For a subset $S \subseteq V$, the induced
restriction $N|_S$ on $S$ sets $N|_S(v) = N(v) \cap S$ (formalized on the subtype $S$ as
\lean{DirectedRecovery.inducedNeighborhood} in~\cite{DirectedRecovery2026}, and on $V$ with empty parent sets outside $S$ as
\lean{resRel} in \cref{sec:reciprocity-classS}).

A \emph{storage code} on $(V, N)$ over a finite nonempty alphabet $\alpha$ (\lean{MotionDepthCodes.IsStorageCode}) is a
nonempty subset $C \subseteq \alpha^V$ such that for every $v \in V$, the coordinate $x_v$ of $x \in C$ is determined by
its parent coordinates $(x_u)_{u \in N(v)}$. Sampling $X \in C$ uniformly defines a Shannon polymatroid
$h : \mathcal{P}(V) \to \R$ (\lean{MotionDepthCodes.IsPolymatroid}) satisfying $h(\emptyset) = 0$, monotonicity,
submodularity, and validity $h(N(v) \cup \{v\}) = h(N(v))$ at every $v \in V$ (\lean{MotionDepthCodes.IsValid};
see~\cite[Section~4]{FourthPaper2026} and~\cite[Section~2]{SeventhPaper2026}, Lean-verified). We say $(V, N)$ satisfies the
polymatroid half-rate bound (\lean{MotionDepthCodes.RateBound} at~$1/2$) if every polymatroid $h$ valid for $N$ obeys
\begin{equation}\label{eq:half-rate}
  2 h(V) \le \sum_{v \in V} h(\{v\}),
\end{equation}
which specializes on any storage code $C \subseteq \alpha^V$ to $2 \log \card{C} \le \card{V} \log \card{\alpha}$, i.e.,
$\card{C}^2 \le \card{\alpha}^{\card{V}}$ (\lean{MotionDepthCodes.card_sq_le_of_rateBound_half}).

We recall the multi-record nested elimination engine of~\cite[Section~3]{SeventhPaper2026}. Let $U \subseteq V$ be a
\emph{context} (already contracted vertices), $H \subseteq V \setminus U$ be a \emph{piece}, and $W = V \setminus (U \cup H)$
be the set of \emph{outside} vertices. \emph{Elimination data} for $H$ over an index set $I$ of $m = \card{I}$ records
consist of subsets $P_r \subseteq H$ ($r \in I$), an elimination list $l = [v_0, \dots, v_{k-1}]$ with underlying set $H$,
an \emph{own record} map $a : V \to I$ and \emph{supplying record} map $b : V \to I$ with $v \in P_{a(v)}$ and $b(v) \ne a(v)$
for all $v \in l$, and an optional outside-parent map $\mathrm{par} : V \to \mathrm{Option}\ V$. At step $j < k$, given an
initial revealed set $S_0$, write $S_j = S_0 \cup \{v_0, \dots, v_{j-1}\}$ (\lean{HalfRate.rev}) for the revealed prefix
and $Q_j = S_j \cup P_{b(v_j)}$ for the supply. In~\cite{SeventhPaper2026}, relational nestedness
(\lean{HalfRate.RelCharged} applied to $N, U, P, b, \mathrm{par}, W, S_0, L_0, l$) is defined recursively along $l$ by
threading a state map $L : V \to \mathcal{P}(V)$ recording for each outside vertex $w \in W$ the set $Q_{j'} \cup \{v_{j'}\}$
at the most recent step $j'$ charged to~$w$:
\begin{itemize}
\item if $\mathrm{par}(v_j) = \mathrm{none}$ (\emph{free step}), $N(v_j) \subseteq U \cup Q_j$, and $L$ is unchanged;
\item if $\mathrm{par}(v_j) = \mathrm{some}\ w$ (\emph{charged step}), $w \in W$, $N(v_j) \subseteq U \cup \{w\} \cup Q_j$,
  and $L(w) \subseteq Q_j$, after which $L(w)$ is updated to $Q_j \cup \{v_j\}$.
\end{itemize}
By Lemma~E of~\cite[Section~3.1, Lemma~3.2]{SeventhPaper2026} (Lean-verified as \lean{HalfRate.nested_elimination}),
whenever $g$ is a polymatroid and the charging is nested, $2 g(H) \le \sum_{r \in I} g(P_r) + \sum_{w \in W} I_g(w; H)$.
A \emph{piece certificate} for $H$ in context $U$ (\lean{HalfRate.PieceCert}) bundles $m \in \N$ pairwise disjoint records
$P_0, \dots, P_{m-1}$ partitioning $H$ (note $a(v) \ne b(v)$ forces $m \ge 2$ whenever $H \ne \emptyset$), a list $l$ with
$l.\mathrm{toFinset} = H$, maps $a, b : V \to \{0, \dots, m-1\}$ with $v \in P_{a(v)}$ and $a(v) \ne b(v)$ for $v \in l$,
and $\mathrm{par} : V \to \mathrm{Option}\ V$ satisfying \lean{HalfRate.RelCharged} from initial state $S_0 = \emptyset$,
$L_0 \equiv \emptyset$ with outside set $W = V \setminus (U \cup H)$. A \emph{global elimination certificate} over $m$
records (\lean{HalfRate.ElimCert}) is the global special case $U = \emptyset$, $H = V$ (where $W = \emptyset$ and every step
is uncharged). A \emph{decomposition} of $N$ from context $U$ (\lean{HalfRate.Decomp}) is a sequence of disjoint pieces
$H_1, \dots, H_t$ partitioning $V \setminus U$ such that each $H_s$ has a piece certificate in context
$U \cup H_1 \cup \dots \cup H_{s-1}$. By Theorem~D of~\cite[Section~3.2]{SeventhPaper2026} (Lean-verified as
\lean{HalfRate.rateBound_of_decomp} and \lean{HalfRate.card_sq_le_of_decomp}), a decomposition from context $\emptyset$
implies~\eqref{eq:half-rate} and $\card{C}^2 \le \card{\alpha}^{\card{V}}$ for every storage code $C \subseteq \alpha^V$.

\subsection{Scalar linear codes on triangle-free graphs of maximum degree~\texorpdfstring{$4$}{4}}

Let $G = (V, E)$ be a finite simple undirected graph and $F$ be an arbitrary field. Following~\cite[Section~8]{EighthPaper2026},
a matrix $M \in F^{V \times V}$ has unit diagonal and support in $G$ if $M_{v,v} = 1$ for all $v \in V$ and $M_{v,u} \ne 0
\implies \{v, u\} \in E$ for $u \ne v$. Its subspace $\ker M = \{x \in F^V : M x = 0\}$ is a scalar linear storage code on
$G$: each coordinate $x_v = -\sum_{u \ne v} M_{v,u} x_u$ is recovered linearly from its active recovery neighbors
$N_M(v) = \{u \in V \setminus \{v\} : M_{v,u} \ne 0\}$ (\lean{MotionDepthCodes.recSet}). Following~\cite[Section~6.6]{FourthPaper2026}
and~\cite[Section~8.2]{EighthPaper2026}, a vertex $v$ satisfies $(\star\star)$ (\lean{MotionDepthCodes.Star}) if either
$\card{N_M(v)} \le 2$, or $\card{N_M(v)} = 3$ and $v$ reciprocates every outgoing recovery arc ($v \in N_M(u) \implies u \in
N_M(v)$); it satisfies $(\star 4)$ (\lean{BoundedNesting.Star4}) if either $\card{N_M(v)} \le 3$, or $\card{N_M(v)} = 4$
and $v$ reciprocates every outgoing recovery arc (automatically satisfied whenever $M$ is supported in an undirected graph
$G$ of maximum degree at most~$4$). We call $v \in V$ \emph{heavy} (or exceptional) if $v$ violates $(\star\star)$, and
write $H_M \subseteq V$ (\lean{BoundedNesting.heavyVerts}) for the set of heavy vertices. Theorem~SH
of~\cite[Section~8.2, Theorem~8.3]{EighthPaper2026} (Lean-verified as \lean{BoundedNesting.theoremSH}) proved $\dim_F \ker M
\le \lfloor \card{V}/2 \rfloor$ assuming $(\star 4)$ everywhere together with $\card{H_M} \le 2$ or \emph{both} (S1) no
vertex has two distinct heavy recovery neighbors (\lean{BoundedNesting.SepS1}), and (S2) no heavy vertex has a heavy
recovery neighbor (\lean{BoundedNesting.SepS2}).

\subsection{Walks, turn labels, and the deep-nesting walk \texorpdfstring{$U(s_0, m, 1)$}{U(s0, m, 1)}}

A $1$-head walk on $\Z$ is a map $x : \N \to \Z$ with $\abs{x(t+1) - x(t)} \le 1$; for $k$ work heads, each head $h \in
\{0, \dots, k-1\}$ moves by $\{-1, 0, +1\}$ at every step. Following~\cite[Section~2.3]{EighthPaper2026}
(\leanfile{BoundedNesting/WalkDag.lean}), the \emph{walk dependency DAG} over horizon $T$ has vertex set $\{0, \dots, T-1\}$,
with \emph{state edges} $t \to t - 1$ for $1 \le t < T$ and \emph{long-hop edges} $u \to p$ whenever some head visits the
same cell at times $p$ and $u$ with $p + 2 \le u < T$ and no intermediate visit in $(p, u)$ (\lean{BoundedNesting.allHops}).
By the label-cut principle~\cite[Section~2.3, Lemma~2.3]{EighthPaper2026} (Lean-verified as
\lean{BoundedNesting.labelCut_depth_walkDag}), for any label $\Psi : \N \to \{0, \dots, q-1\}$ ($q \ge 1$), deleting the
state edges at step descents $\Psi(t) < \Psi(t-1)$ (\lean{BoundedNesting.descentsOf}) and the violated hop edges with
$\Psi(u) \le \Psi(p)$ (\lean{BoundedNesting.violatedOf}) leaves at most $q - 1$ long-hop edges along any directed path in
the walk DAG. On a cyclic schedule whose long hops partition into closed cell chains, Lemma~Zq
of~\cite[Section~5.5, Lemma~5.15]{FourthPaper2026} (Lean-verified as \lean{MotionDepthCodes.lemma_Zq} and specialized to
turn-counting labels in \lean{BoundedNesting.turnCounter_cost}) expresses $q$ times the total number of cut descents and
violated hops identically as the sum of step increments $(\Psi(t) - \Psi(t-1)) \bmod q$ plus reduced hop errors
$(\Psi(u) - \Psi(p) - 1) \bmod q$.

For integers $0 \le m \le s_0$, define the dyadic scales $R_i = 2^{s_0 - i}$ for $0 \le i \le s_0$ (\lean{scaleR}) and the
nested \emph{stroke lengths} $\Pi_\kappa = (m - \kappa + 1) R_\kappa$ for $0 \le \kappa \le m$ (\lean{strokeLen}). As in
\cite[Definition~5.3]{FourthPaper2026} and~\cite[Section~2.5]{EighthPaper2026}, the \emph{deep-nesting walk} $U(s_0, m, 1)$ has fundamental period
$P = 4 \Pi_0 = 4(m+1)2^{s_0}$ (\lean{dnpPeriod}) consisting of four level-$0$ strokes (forward, backward, backward,
forward). Recursively, for $\kappa < m$, a level-$\kappa$ stroke of duration $\Pi_\kappa = R_\kappa + 2 \Pi_{\kappa+1}$
executes a plain segment of $R_\kappa$ unit steps, followed by a level-$(\kappa+1)$ \emph{back stroke} of duration
$\Pi_{\kappa+1}$ reversing direction across those $R_\kappa$ cells (making an \emph{apex} turn at offset $R_\kappa$),
followed by a level-$(\kappa+1)$ \emph{forward stroke} of duration $\Pi_{\kappa+1}$ restoring the original direction (making
a \emph{valley} turn at offset $R_\kappa + \Pi_{\kappa+1}$); at the base level $\kappa = m$, a stroke is simply $R_m =
2^{s_0 - m}$ plain steps. On a single copy of $U(s_0, m, 1)$ (and on $\DNP(s_0, m)$ at any shift), the turn counter
of~\cite[Section~5.2, Theorem~5.2 and Corollary~5.5]{EighthPaper2026} cuts hop-depth to $2^{m + O(1)}/\eps$ at edge density
$\eps$. The \emph{deep-nesting pair} $\DNP(s_0, m)$ runs head~$1$ on $U(s_0, m, 1)$ and head~$2$ on a time-shifted copy
$t \mapsto U(s_0, m, 1)((t + \xi) \bmod P)$ over a common tape.

\subsection{Formal status of the results}\label{sec:prelim-lean}

Every numbered theorem, proposition, lemma, and corollary introduced in \cref{sec:reciprocity,sec:crux,sec:dnp} has a
complete, kernel-checked proof in the accompanying Lean~4 library \leanfile{lean/DirectedReciprocity/} (Lean and Mathlib
version \texttt{v4.34.1}). Immediately below each numbered statement, a header line
\[
  \text{\textsf{Lean:} \leanname{declaration_1}, \leanname{declaration_2} \hfill [\leanname{Module.lean}]}
\]
names the exact Lean declarations (relative to the namespace \texttt{DirectedReciprocity}) and source file establishing
that result. Each new result is followed by an \emph{Outline of the formal proof} explaining the mathematical structure of
the machine-checked proof faithfully. As verified automatically by \leanfile{lean/PaperAxioms.lean}
(\cref{app:index,app:verify}), every cited declaration of our library---and every declaration of earlier libraries named in
the text---depends only on the three standard Lean axioms \texttt{propext}, \texttt{Classical.choice}, and
\texttt{Quot.sound}; no file uses \texttt{sorry}, \texttt{admit}, \texttt{native\_decide}, \texttt{unsafe}, or any custom
\texttt{axiom}. Finite case checks in \cref{sec:reciprocity} are discharged inside Lean's kernel via \texttt{decide},
\texttt{decide +kernel}, or \texttt{decide +revert}. The standalone Python script \leanfile{tools/reproduce_numerics.py} is
used solely to cross-check illustrative numerical remarks in \cref{rem:deg4-numerics,rem:dnp-numerics}, never inside any
proof.

\section{Positional Nesting, Digon Cycles, and Theorem~S\texorpdfstring{$^*$}{*}}\label{sec:reciprocity}

The multi-record nested elimination engine of~\cite{SeventhPaper2026} certifies the polymatroid half-rate bound
$2 h(V) \le \sum_{v \in V} h(\{v\})$ on a finite recovery relation $(V, N)$ by peeling $V$ into certified pieces
(\cref{sec:prelim}). In~\cite{SeventhPaper2026}, Theorem~N3 constructed piece decompositions for symmetric (undirected)
triangle-free graphs of maximum degree~$3$, Theorem~U3 extended this to loopless TT-free relations with subcubic
underlying graph, and Theorem~T constructed them for loopless TT-free relations in which the vertices with at least
three parents form an acyclic set; the intermediate target named in item~2 of~\cite[Section~9.2]{SeventhPaper2026} asked
whether piece decompositions extend to the full directed class $\ClassS$ of Theorem~S of~\cite[Section~6.6]{FourthPaper2026}
(the all-code version of the fourth paper's Theorem~S on its whole class, lying between Theorem~U3 and Conjecture~P,
numbered Conjecture~9.1 there, which remains open), where reciprocal pairs and pure one-way arcs coexist freely subject to
reciprocal closure at vertices of in-degree~$3$, while item~6 of~\cite[Section~9.2]{SeventhPaper2026} (see
also~\cite[Remarks~7.3 and~7.5]{SeventhPaper2026}) asked whether unbounded record counts $m$ can rescue obstructed cores
such as $\Ex$'s hexagon or supply global elimination certificates on cubic graphs such as the Petersen graph.

In this section we resolve all three questions. We first give a stateless, purely positional reformulation of relational
nestedness (\cref{prop:PN}) that turns impossibility proofs across unbounded record counts $m \ge 1$ into finite-order
locking arguments. We then establish exact piece-certificate dichotomies for stride-$s$ hub cores (\cref{prop:CS,cor:x15})
and arbitrary chordless digon cycles (\cref{thm:DC}), classify global record counts across generalized Petersen graphs
(\cref{prop:GP}), and prove \textbf{Theorem~S$^*$} (\cref{thm:Sstar}), \textbf{Corollary~B$^*$} (\cref{cor:Bstar}), and
the $21$-vertex sharpness witness $\Gto \in \ClassS$ (\cref{prop:G21}).

\subsection{Positional characterization of nested elimination}\label{sec:reciprocity-pn}

Fix a finite vertex set $V$, a recovery relation $N : V \to \mathcal{P}(V)$, a context $U \subseteq V$, an index set $I$
of records $P_r \subseteq V$ ($r \in I$), a supplying-record map $b : V \to I$, an optional outside-parent map
$\mathrm{par} : V \to \mathrm{Option}\ V$, an outside set $W \subseteq V$, and a list $l = [v_0, \dots, v_{k-1}]$ of
vertices. For $0 \le j \le k$, write $S_j = \{v_0, \dots, v_{j-1}\}$ (\lean{HalfRate.rev} from $\emptyset$) for the
prefix revealed strictly before step~$j$.

\begin{proposition}[Positional characterization of nested supply]\label{prop:PN}
\formal{PositionalNest.lean}{relCharged_of_nested, disjoint_nested_iff, relCharged_iff_nested, PieceCert.mem_supply_of_idxOf_lt, PieceCert.digon_ne}
Let $V, N, U, P, b, \mathrm{par}, W, l$ be as above.
\begin{enumerate}
\item \emph{(Exact positional equivalence.)} The stateful predicate
  \lean{HalfRate.RelCharged} applied to $(N, U, P, b, \mathrm{par}, W, \emptyset, \lambda\_.\,\emptyset, l)$ holds if and only if:
  \begin{enumerate}
  \item for every $0 \le j < k$ with $\mathrm{par}(v_j) = \mathrm{none}$, $N(v_j) \subseteq U \cup S_j \cup P_{b(v_j)}$;
  \item for every $0 \le j < k$ and $w \in V$ with $\mathrm{par}(v_j) = \mathrm{some}\ w$, $w \in W$ and
    $N(v_j) \subseteq U \cup \{w\} \cup S_j \cup P_{b(v_j)}$; and
  \item for every $0 \le j < j' < k$ and $w \in V$ with $\mathrm{par}(v_j) = \mathrm{par}(v_{j'}) = \mathrm{some}\ w$,
    \[
      P_{b(v_j)} \subseteq S_{j'} \cup P_{b(v_{j'})}.
    \]
  \end{enumerate}
\item \emph{(Disjoint-record simplification.)} Suppose the records $\{P_r\}_{r \in I}$ are pairwise disjoint, and let
  $0 \le j, j' < k$. Then $P_{b(v_j)} \subseteq S_{j'} \cup P_{b(v_{j'})}$ holds if and only if either
  $b(v_j) = b(v_{j'})$ or $P_{b(v_j)} \subseteq S_{j'}$.
\item \emph{(First-appearance supply forcing and digon proper coloring.)} Let $c$ be a piece certificate
  (\lean{HalfRate.PieceCert}) for $H \subseteq V$ in context $U$ with $U \cap H = \emptyset$, and write $\tau(u) =
  \mathrm{idxOf}(u, c.l)$ for $u \in H$. If $u, v \in H$ satisfy $v \in N(u)$ and $\tau(u) < \tau(v)$, then
  $v \in c.P_{c.b(u)}$ and $c.b(u) = c.a(v)$. Consequently, for any two distinct vertices $u, v \in H$ forming a digon
  ($u \in N(v)$ and $v \in N(u)$), their own records are distinct: $c.a(u) \ne c.a(v)$.
\end{enumerate}
\end{proposition}

\begin{outline}
For~(1), induction along $l$ with general initial revealed set $S_0$ and state map $L_0$ (\lean{relCharged_step_none},
\lean{relCharged_step_some}, \lean{relCharged_step_nested}) shows that if \lean{HalfRate.RelCharged} holds, then at step
$j'$ charged to $w$ after an earlier step $j < j'$ charged to $w$, the state $L(w)$ contains $\{v_j\} \cup S_j \cup
P_{b(v_j)}$ by monotonicity along intermediate updates, forcing $P_{b(v_j)} \subseteq S_{j'} \cup P_{b(v_{j'})}$.
Conversely (\lean{relCharged_of_nested}), because $v_j \in S_{j'}$ and $S_j \subseteq S_{j'}$ hold automatically whenever
$j < j'$, the inclusion $P_{b(v_j)} \subseteq S_{j'} \cup P_{b(v_{j'})}$ upgrades to $\{v_j\} \cup S_j \cup P_{b(v_j)}
\subseteq S_{j'} \cup P_{b(v_{j'})}$, preserving the invariant that $L(w)$ equals either $L_0(w)$ or $\{v_{j_*}\} \cup
S_{j_*} \cup P_{b(v_{j_*})}$ for the latest index $j_* < j'$ charged to~$w$. Specializing to $S_0 = \emptyset$ and
$L_0 \equiv \emptyset$ gives \lean{relCharged_iff_nested}.

For~(2) (\lean{disjoint_nested_iff}), if $b(v_j) \ne b(v_{j'})$, then $P_{b(v_j)} \cap P_{b(v_{j'})} = \emptyset$ by
pairwise disjointness, so $P_{b(v_j)} \subseteq S_{j'} \cup P_{b(v_{j'})}$ forces every $x \in P_{b(v_j)}$ into $S_{j'}$;
the converse is immediate.

For~(3), let $j_u = \tau(u) < \tau(v)$. By definition of first appearance (\lean{not_mem_rev_empty_of_lt_idxOf}),
$v \notin S_{j_u}$. Since $v \in H$ and $U \cap H = \emptyset$, we have $v \notin U$ and $v \notin V \setminus (U \cup H)$.
Applying~(1a) or~(1b) at index $j_u$ to $v \in N(u)$ leaves only $v \in c.P_{c.b(u)}$. Because $v \in c.P_{c.a(v)}$ and
the records of $c$ are pairwise disjoint, $c.b(u) = c.a(v)$ (\lean{PieceCert.mem_supply_of_idxOf_lt}). If $u \ne v$ form
a digon in $H$, then $\tau(u) \ne \tau(v)$; assuming $\tau(u) < \tau(v)$ by symmetry gives $c.b(u) = c.a(v)$, whereas
$c.a(u) \ne c.b(u)$, so $c.a(u) \ne c.a(v)$ (\lean{PieceCert.digon_ne}).
\end{outline}

\subsection{Stride-\texorpdfstring{$s$}{s} cores, chordless digon cycles, and generalized Petersen graphs}\label{sec:reciprocity-cycles}

Let $s \ge 3$ and $k \ge 1$. Define the canonical \emph{stride-$s$ hub relation} \lean{strideRel}$(s, k)$ on
$\Z_{sk} \sqcup \Z_s$ by giving each cycle vertex $i \in \Z_{sk}$ parents $\{i - 1, i + 1\} \subseteq \Z_{sk}$ together
with the outside hub $w_{i \bmod s} \in \Z_s$, and giving hubs empty parent sets; let $H_{s,k} = \Z_{sk} \subseteq
\Z_{sk} \sqcup \Z_s$ (\lean{strideCore}) be the cycle core. More generally, in any recovery relation $(V, N)$ with
context $U \subseteq V$, an injective cycle $f : \Z_L \hookrightarrow V$ with image $H = f(\Z_L)$ is a \emph{stride-$3$
cycle} in context $U$ (\lean{IsStride3Cycle}) if $3 \mid L$ and there exist three vertices $w_0, w_1, w_2 \in V \setminus
(U \cup H)$ such that $w_{i \bmod 3} \in N(f(i))$ for every $i \in \Z_L$.

\begin{proposition}[Stride-$s$ digon cycle dichotomy]\label{prop:CS}
\formal{StrideCycles.lean}{stride3_no_pieceCert, strideRel, strideCore, strideCore3_no_pieceCert, strideGe4_pieceCert, strideEven_pieceCert_two, strideCore_pieceCert_iff}
\begin{enumerate}
\item \emph{(Unbounded-$m$ stride-$3$ impossibility.)} Let $(V, N)$ be a finite recovery relation, $U, H \subseteq V$
  disjoint subsets, $k \ge 1$, and $f : \Z_{3k} \hookrightarrow V$ a bijection onto $H$ such that $f(i \pm 1) \in N(f(i))$
  for all $i \in \Z_{3k}$. If there exist $w_0, w_1, w_2 \in V \setminus (U \cup H)$ with $w_{i \bmod 3} \in N(f(i))$ for
  all $i \in \Z_{3k}$, then $H$ admits \emph{no} piece certificate (\lean{HalfRate.PieceCert} $N\ U\ H$) over any number of
  records $m \in \N$.
\item \emph{(Stride-$s$ threshold at $s = 4$.)} For all integers $s \ge 3$ and $k \ge 1$, the canonical stride-$s$ core
  $H_{s,k}$ admits a piece certificate in context $\emptyset$ under \lean{strideRel}$(s, k)$ (for some $m \in \N$) if and
  only if $s \ge 4$. Moreover, $m = 2$ co-supplied records suffice whenever $s \ge 4$ is even, and $m = 3$ co-supplied
  records suffice for every $s \ge 4$.
\end{enumerate}
\end{proposition}

\begin{outline}
For~(1) (\lean{stride3_no_pieceCert}), suppose $c$ is an $m$-record piece certificate for $H$ in context~$U$, and write
$\tau(i) = \mathrm{idxOf}(f(i), c.l)$. Because $w_{i \bmod 3} \in N(f(i))$ lies outside $U \cup H$, step $\tau(i)$ must
designate $\mathrm{par}(f(i)) = \mathrm{some}\ w_{i \bmod 3}$. We first establish the \emph{Locking Lemma}: if
$f(u) \in N(f(i))$, $j \equiv i \pmod 3$, and $\tau(i) \le \tau(j) < \tau(u)$, then $c.b(f(j)) = c.a(f(u))$. Indeed,
\cref{prop:PN}(3) gives $f(u) \in c.P_{c.b(f(i))}$ and $c.b(f(i)) = c.a(f(u))$; if $\tau(i) < \tau(j)$, then step $\tau(i)$
and step $\tau(j)$ share outside parent $w_{i \bmod 3}$, so \cref{prop:PN}(1c,2) forces either $c.b(f(i)) = c.b(f(j))$ or
$c.P_{c.b(f(i))} \subseteq S_{\tau(j)}$, and the latter is impossible because $f(u) \in c.P_{c.b(f(i))}$ has $\tau(j) <
\tau(u)$ and so is unrevealed at $\tau(j)$. As an immediate consequence, if $j_1 \equiv j_2 \pmod 3$ satisfy
$f(u_1) \in N(f(j_1))$, $f(u_2) \in N(f(j_2))$, and $\max(\tau(j_1), \tau(j_2)) < \min(\tau(u_1), \tau(u_2))$, applying
locking at $\max(\tau(j_1), \tau(j_2))$ yields $c.a(f(u_1)) = c.a(f(u_2))$.

Now let $z \in \Z_{3k}$ maximize $\tau(z)$ over all of $\Z_{3k}$, set $r_0 = z \bmod 3$, and let $x$ maximize $\tau(x)$
among $\{i \in \Z_{3k} : i \not\equiv r_0 \pmod 3\}$. Let $r_x = x \bmod 3 \ne r_0$ and let $r_3 \in \{0,1,2\} \setminus
\{r_0, r_x\}$ be the remaining residue class modulo~$3$. Since $f(y \pm 1)$ realize the two residues distinct from
$y \bmod 3$ for every $y \in \Z_{3k}$, $x$ has a cycle neighbor $u_3 \equiv r_3 \pmod 3$ and $z$ has cycle neighbors
$v_3 \equiv r_3 \pmod 3$ and $v_x \equiv r_x \pmod 3$. Because $u_3, v_3 \not\equiv r_0 \pmod 3$ and neither lies in
$r_x$, the choice of $x$ gives $\tau(u_3) < \tau(x) < \tau(z)$ and $\tau(v_3) < \tau(x) < \tau(z)$. Since
$u_3 \equiv v_3 \equiv r_3 \pmod 3$, the two-arc corollary forces $c.a(f(x)) = c.a(f(z))$. If $x = v_x$, then $f(x)$ and
$f(z)$ form a digon along the cycle, contradicting $c.a(f(x)) \ne c.a(f(z))$ (\cref{prop:PN}(3)). Otherwise $v_x \ne x$
lies in $r_x \ne r_0$, so $\tau(v_x) < \tau(x) < \tau(z)$; locking along $f(z) \in N(f(v_x))$ from $v_x$ to $x$ forces
$c.b(f(x)) = c.a(f(z)) = c.a(f(x))$, contradicting $c.a(f(x)) \ne c.b(f(x))$.

For~(2), when $s = 3$,~(1) rules out any \lean{HalfRate.PieceCert} (\lean{strideCore3_no_pieceCert}). When $s \ge 4$ is
even (\lean{strideEven_pieceCert_two}), $2 \mid s \mid sk$, so partitioning $H_{s,k}$ by parity into two records $P_0,
P_1$ and eliminating all even indices before all odd indices makes $b(i) = 1 - (i \bmod 2)$ constant on each hub fiber
$i \bmod s$ (co-supplied). For general $s \ge 4$ (\lean{strideGe4_pieceCert}), pulling back the proper $3$-coloring
\lean{HalfRate.ccol}$(s)$ of $\Z_s$ along $i \mapsto i \bmod s$ partitions $H_{s,k}$ into three records $P_0, P_1, P_2$;
on each residue class $i \equiv r \pmod s$, both neighbors $i \pm 1$ have colors $(r \pm 1) \bmod s$ in $\Z_s$, and as in
Theorem~N3 of~\cite[Theorem~4.3]{SeventhPaper2026} (\lean{HalfRate.cyclePiece}; see also~\cite[Lemma~5.7(b)]{SeventhPaper2026}),
ordering the fibers so that one neighbor color is already revealed while the other serves as the constant supplying record
$b(i)$ on fiber~$r$ yields an explicit $3$-record co-supplied certificate.
\end{outline}

Applying \cref{prop:CS}(1) to the $15$-vertex relation $\Ex$ of~\cite[Section~7.2, Remark~7.3]{SeventhPaper2026} (\lean{HalfRate.x15}),
whose unique minimum nonempty core is the induced $6$-cycle $\{0,1,2,3,4,5\}$ (see~\cite[Remark~7.3]{SeventhPaper2026},
Lean-verified as \lean{HalfRate.x15_cores}) attached to outside vertices $6, 7, 8$ with stride~$3$, immediately settles
the unbounded-$m$ core question of~\cite[Remark~7.3; Section~9.2, item~6]{SeventhPaper2026}:

\begin{corollary}[Unbounded-$m$ obstruction on $\Ex$]\label{cor:x15}
\formal{ClassSExamples.lean}{x15_hexagon_no_pieceCert}
The unique minimum nonempty core $\{0, 1, 2, 3, 4, 5\}$ of $\Ex$ admits no piece certificate
\lean{HalfRate.PieceCert} \lean{HalfRate.x15} $\emptyset\ \{0,1,2,3,4,5\}$ for any number of records $m \ge 1$.
\end{corollary}

\begin{outline}
Instantiate \lean{stride3_no_pieceCert} with $k = 2$, $f = [0,1,2,3,4,5]$, and outside hub triple $w = [6,7,8]$; all
finite membership and injectivity side conditions discharge by \texttt{decide}.
\end{outline}

We next show that on \emph{any} chordless digon cycle $C_L$ ($L \ge 4$) in a TT-free relation where each cycle vertex has
at most one active outside parent, the stride-$3$ hub configuration is the \emph{unique} obstruction to piece
certifiability across all $m \ge 1$.

\begin{theorem}[Universal chordless digon cycle dichotomy]\label{thm:DC}
\formal{DigonCycle.lean}{IsStride3Cycle, not_pieceCert_of_isStride3Cycle, exists_non_stride3_pos, digonCycle4_pieceCert, digonCycleGe5_pieceCert, digonCycle_pieceCert_of_not_stride3, digonCycle_pieceCert_iff}
Let $(V, N)$ be a finite recovery relation, $U \subseteq V$ a context, $L \ge 4$, and $f : \Z_L \hookrightarrow V$ an
injective map with image $H = f(\Z_L)$ disjoint from~$U$. Assume that:
\begin{enumerate}
\item no outside vertex $w \in V \setminus (U \cup H)$ is a common parent of two consecutive cycle vertices $f(i)$ and
  $f(i+1)$ (local TT-freeness);
\item $f(i \pm 1) \in N(f(i))$ for every $i \in \Z_L$, and conversely $f(j) \in N(f(i)) \implies j \in \{i - 1, i + 1\}$
  (chordless digon cycle); and
\item each $f(i)$ has at most one active parent in $V \setminus (U \cup H)$.
\end{enumerate}
Then $H$ admits a piece certificate (\lean{HalfRate.PieceCert} $N\ U\ H$) if and only if $f$ is not a stride-$3$ cycle in
context~$U$ (\lean{IsStride3Cycle} $N\ U\ f$ fails). When \lean{IsStride3Cycle} $N\ U\ f$ fails, there exists an explicit
piece certificate with $m = 2$ records for $L = 4$ and $m \le L + 1$ records for $L \ge 5$.
\end{theorem}

\begin{outline}
Necessity (\lean{not_pieceCert_of_isStride3Cycle}) follows directly from \cref{prop:CS}(1). For sufficiency
(\lean{digonCycle_pieceCert_of_not_stride3}), write $\mathrm{op}(i) \in \mathrm{Option}\ V$ for the unique outside parent
of $f(i)$ in $V \setminus (U \cup H)$ (or $\mathrm{none}$ if none exists).

When $L = 4$ (\lean{digonCycle4_pieceCert}), partition $H$ into $P_0 = \{f(0), f(2)\}$ and $P_1 = \{f(1), f(3)\}$, set
$b = 1 - a$, and eliminate in order $[f(0), f(2), f(1), f(3)]$. By~(1), $\mathrm{op}(i)$ cannot coincide on adjacent cycle
vertices, so any shared outside parent occurs on opposite vertices $\{0, 2\}$ or $\{1, 3\}$, where $b$ is constant.

When $L \ge 5$ (\lean{digonCycleGe5_pieceCert}), we first prove (\lean{exists_non_stride3_pos}) that failure of
\lean{IsStride3Cycle} $N\ U\ f$ yields an index $s_0 \in \Z_L$ such that $\mathrm{op}(s_0) = \mathrm{none}$ or
$\mathrm{op}(s_0) \ne \mathrm{op}(s_0 + 3)$: otherwise
every $i \in \Z_L$ has $\mathrm{op}(i) = \mathrm{some}\ w_i$ periodic of period~$3$ along $\Z_L$, forcing $3 \mid L$ (since
$3 \nmid L$ would make $\mathrm{op}$ constant on $\Z_L$, violating~(1)) and exhibiting $w_0, w_1, w_2$ witnessing
\lean{IsStride3Cycle}. Rotating $f$ so that $s_0 = 2$ gives $\mathrm{op}(2) = \mathrm{none}$ or $\mathrm{op}(2) \ne
\mathrm{op}(5 \bmod L)$. We then eliminate $H$ over $L + 1$ records $\{0, 1, \dots, L\}$ in the two-sided order
\[
  [f(1),\; f(2),\; f(0),\; f(L-1),\; f(L-2),\; \dots,\; f(4),\; f(3)]
\]
(\lean{dcPerm}), assigning own records (\lean{dcCol}) and supplying records (\lean{dcSupply}) as follows:
\begin{itemize}
\item $f(0)$ and $f(2)$ have own record~$0$; $f(1)$ has own record~$1$ and supplying record~$0$ (covering both neighbors
  $f(0), f(2)$ at step~$0$);
\item $f(3)$ has own record~$3$, and for $k \in \{4, \dots, L-1\}$ we set $a(f(k)) = 3$ if $\mathrm{op}(2) = \mathrm{some}\ w
  = \mathrm{op}(k + 1 \bmod L)$ and $a(f(k)) = k$ otherwise. By our rotation $\mathrm{op}(5 \bmod L) \ne \mathrm{op}(2)$
  (when $\mathrm{op}(2)$ is $\mathrm{some}\ w$), we always have $a(f(4)) = 4 \ne 3$, and by~(1) no two consecutive $k, k+1$
  can both merge into record~$3$;
\item at step~$1$, $f(2)$ uses supplying record~$3$ (covering $f(3)$, with $f(1)$ already revealed); for steps
  $j \in \{2, \dots, L-2\}$ eliminating $f(k)$ ($k \in \{0, L-1, \dots, 4\}$), its successor $f(k+1 \bmod L)$ is already
  revealed and we set $b(f(k)) = a(f(k - 1 \bmod L))$ covering its predecessor; at step $L - 1$, $f(3)$ has both neighbors
  $f(2), f(4)$ already revealed and uses the empty auxiliary record~$L$.
\end{itemize}
To verify nestedness (\cref{prop:PN}(1c)), suppose $0 \le j < j' < L$ share $\mathrm{op} = \mathrm{some}\ w$. Step~$0$
($f(1)$) has $P_{b(f(1))} = P_0 = \{f(0), f(2)\}$; by~(1) neither $f(2)$ ($j' = 1$) nor $f(0)$ ($j' = 2$) can share
$\mathrm{op}(1)$, and for all $j' \ge 3$ we already have $P_0 \subseteq S_3 \subseteq S_{j'}$. Step~$1$ ($f(2)$) has
$b(f(2)) = 3$; by~(1) the final vertex $f(3)$ ($j' = L - 1$) cannot share $\mathrm{op}(2)$, while for any $2 \le j' \le
L - 2$ with $\mathrm{op}(k) = \mathrm{op}(2) = \mathrm{some}\ w$, our definition of \lean{dcCol} placed $f(k - 1 \bmod L)$
into record~$3$, so $b(f(k)) = 3 = b(f(2))$ (co-supplied). Finally, for $2 \le j < j' < L$, either $b(v_j) = 3$, in which
case $\mathrm{op}(v_j) = \mathrm{op}(2) = \mathrm{some}\ w = \mathrm{op}(v_{j'})$, excluding $j' = L - 1$ by~(1) and
forcing $b(v_{j'}) = 3$ for $j' \le L - 2$, or else $b(v_j) = k - 1 \ge 4$ is a singleton $\{f(k-1)\}$ revealed
immediately at step $j + 1$, whereas $j' = j + 1$ cannot share $w$ with $j$ by~(1), giving $P_{b(v_j)} \subseteq S_{j+2}
\subseteq S_{j'}$.
\end{outline}

Before turning to $\ClassS$, we resolve the global-certificate question of~\cite[Remark~7.5; Section~9.2, item~6]{SeventhPaper2026} by
determining exact global record thresholds across the cubic generalized Petersen family $\GP(m_{\mathrm{gp}}, 2)$
($m_{\mathrm{gp}} \ge 5$, \lean{gpGraph}), on vertex set $\{0, \dots, 2 m_{\mathrm{gp}} - 1\}$ with outer rim edges
$i \leftrightarrow (i \pm 1 \bmod m_{\mathrm{gp}})$, inner star edges $m_{\mathrm{gp}} + i \leftrightarrow m_{\mathrm{gp}} +
(i \pm 2 \bmod m_{\mathrm{gp}})$, and spokes $i \leftrightarrow m_{\mathrm{gp}} + i$.

\begin{proposition}[Non-monotone global record hierarchy on $\GP(m_{\mathrm{gp}}, 2)$]\label{prop:GP}
\formal{PetersenHierarchy.lean}{nonempty_elimCert_iff_pieceCert_univ, gpGraph, gp52_eq_petersenGraph, gp52_no_elimCert, petersen_no_pieceCert_univ, gp72_no_elimCert_le3, gp72_elimCert4, gp92_elimCert3, generalizedPetersen_hierarchy}
\begin{enumerate}
\item \emph{(Global equivalence.)} For any finite recovery relation $(V, N)$ and $m \in \N$, there exists a global piece
  certificate $p$ of type \lean{HalfRate.PieceCert} $N\ \emptyset\ V$ with $p.m = m$ if and only if
  \lean{HalfRate.ElimCert} $N\ m$ is nonempty.
\item \emph{($\GP(5,2)$, Petersen graph, $n = 10$.)} $\GP(5,2)$ coincides with \lean{MotionDepthCodes.petersenGraph} and
  admits \emph{no} global elimination certificate (of type \lean{HalfRate.ElimCert} $(\mathrm{nbhd}\ \GP(5,2))\ m$, or
  equivalently global \lean{HalfRate.PieceCert}) for \emph{any} number of records $m \in \N$.
\item \emph{($\GP(7,2)$, $n = 14$.)} $\GP(7,2)$ admits no global elimination certificate with $m \le 3$ records, strictly
  separating $m = 3$ from $m = 4$, and admits an explicit $4$-record global elimination certificate.
\item \emph{($\GP(9,2)$, $n = 18$.)} $\GP(9,2)$ drops back from $4$ records to an explicit $3$-record global elimination
  certificate.
\end{enumerate}
\end{proposition}

\begin{outline}
For~(1) (\lean{nonempty_elimCert_iff_pieceCert_univ}), when $U = \emptyset$ and $H = V$, the outside set $W = V \setminus V
= \emptyset$ forces every step of a piece certificate $p$ to have $\mathrm{par}(v_j) = \mathrm{none}$; given any proper
subset $S \subsetneq V$, taking the first index $j$ of $p.l$ with $v_j \notin S$ gives $S_j \subseteq S$ and hence
$N(v_j) \subseteq S \cup P_{b(v_j)}$, yielding \lean{HalfRate.ElimCert} $N\ p.m$. Conversely, greedy extraction from
\lean{HalfRate.ElimCert} $N\ m$ via \lean{HalfRate.exists_relCharged} produces a global \lean{HalfRate.PieceCert} with $m$
records.

By \cref{prop:PN}(3) (\lean{elimCert_adj_ne}), any $m$-record global elimination certificate on an undirected graph $G$
induces a proper coloring $a : V \to \{0, \dots, m-1\}$ into independent sets; moreover, in an elimination certificate, no
proper subset $S \subsetneq V$ can be \emph{stuck}, meaning that every $v \notin S$ has two neighbors outside $S$ of
distinct $a$-colors (since a single supplying record $P_{b(v)}$ can belong to only one color class). For~(2)
(\lean{gp52_no_elimCert}), replacing $a(u)$ by the minimum vertex index in its own record maps any arbitrary-$m$ global
certificate on $\GP(5,2)$ to a canonical idempotent proper partition of $\{0, \dots, 9\}$, and an exhaustive kernel-checked
Boolean certificate over all canonical proper partitions exhibits a stuck proper subset $S \subsetneq \{0,\dots,9\}$ in
every case (\texttt{decide +kernel}). For~(3) (\lean{gp72_no_elimCert_le3}, \lean{gp72_elimCert4}), any $m \le 3$
certificate embeds into a normalized proper $3$-coloring of $\GP(7,2)$, every one of which admits a kernel-verified stuck
subset, while an explicit $4$-coloring and $14$-step schedule verify \lean{HalfRate.ElimCert} $(\mathrm{nbhd}\ \GP(7,2))\ 4$
by \texttt{decide}. For~(4) (\lean{gp92_elimCert3}), an explicit $3$-coloring and $18$-step schedule verify
\lean{HalfRate.ElimCert} $(\mathrm{nbhd}\ \GP(9,2))\ 3$ by \texttt{decide}.
\end{outline}

\subsection{Piece decomposition and barrier bound for \texorpdfstring{$\ClassS$}{Class S}}\label{sec:reciprocity-classS}

Recall from~\cite[Section~6.6]{FourthPaper2026} that a finite recovery relation $(V, N)$ belongs to $\ClassS$
(\lean{IsClassS}) if $N$ is irreflexive, TT-free (\lean{HalfRate.TTFree}), has maximum in-degree at most~$3$
($\card{N(v)} \le 3$ for all $v \in V$), and satisfies \emph{reciprocal closure at full in-degree}:
\[
  \forall v \in V,\quad \card{N(v)} = 3 \implies \forall u \in V,\; v \in N(u) \implies u \in N(v).
\]
Equivalently, vertices of in-degree~$3$ may have only reciprocal incident arcs (hence out-degree~$3$), whereas vertices of
in-degree at most~$2$ may send arbitrarily many pure one-way outgoing arcs. More generally, $(V, N)$ satisfies relative
$\ClassS$ on $R \subseteq V$ (\lean{ClassSOn}) if the internal restriction $N|_R$ (\lean{resRel}) belongs to $\ClassS$ and
$N$ is TT-free on~$R$. Note that \lean{ClassSOn} is hereditary under passing to arbitrary subsets $R' \subseteq R$.

\begin{theorem}[Piece decomposition and half-rate bound for $\ClassS$]\label{thm:Sstar}
\formal{ClassS.lean}{IsClassS, ClassSOn, exists_inClosed_scc, liftPieceCertInClosed, exists_piece_of_scc_classS, classSOn_decomp, classS_decomp, rateBound_classS, card_sq_le_classS}
Let $(V, N)$ be a finite recovery relation.
\begin{enumerate}
\item For every subset $R \subseteq V$ satisfying \lean{ClassSOn} $N\ R$, $(V, N)$ admits a certified piece decomposition
  (\lean{HalfRate.Decomp} $N\ (V \setminus R)\ \mathrm{pcs}$) from context $V \setminus R$.
\item In particular, if $N \in \ClassS$, then $(V, N)$ admits a certified piece decomposition from $\emptyset$, satisfies
  \lean{MotionDepthCodes.RateBound} $N\ (1/2)$ ($2 h(V) \le \sum_{v \in V} h(\{v\})$ for every valid polymatroid~$h$), and
  every storage code $C \subseteq \alpha^V$ (linear or nonlinear, over any nonempty finite alphabet $\alpha$) satisfies
  \[
    \card{C}^2 \le \card{\alpha}^{\card{V}}.
  \]
\end{enumerate}
\end{theorem}

\begin{outline}
By strong induction on $\card{R}$ (\lean{classSOn_decomp}), if $R$ is nonempty, condensation of the reachability relation
on $R$ yields a nonempty \emph{in-closed} strongly connected component $S \subseteq R$ (\lean{exists_inClosed_scc}):
every parent in $R$ of a vertex $v \in S$ already lies in~$S$, and every $u, v \in S$ reach one another via directed paths
in $M := N|_S$. Because $S$ is in-closed inside $R$, any piece certificate for $H \subseteq S$ in empty context under $M$
lifts immediately to a piece certificate for $H$ in context $V \setminus R$ under $N$ (\lean{liftPieceCertInClosed}), after
which induction applies to $R \setminus H$.

It therefore suffices to find a nonempty certified piece $H \subseteq S$ in context $\emptyset$ for a strongly connected
relation $M \in \ClassS$ supported on~$S$ (\lean{exists_piece_of_scc_classS}). Let $\mathcal{D} = {}$\lean{digonGraph}$(M)$ be
the undirected simple graph of reciprocal pairs on $S$; since $M$ is TT-free, $\mathcal{D}$ is triangle-free
(\lean{digonGraph_cliqueFree3}). Write $W_S = \{v \in S : \card{M(v)} = 2\}$.
\begin{enumerate}
\item \emph{Low in-degree vertex.} If some $v \in S$ has $\card{M(v)} \le 1$, then $H = \{v\}$ is certified by two records
  $\{v\}$ and $\emptyset$ (\lean{HalfRate.singlePiece}), uncharged if $\card{M(v)} = 0$ or charged to its sole parent if
  $\card{M(v)} = 1$.
\item \emph{Directed cycle in $W_S$.} Assume henceforth $\card{M(v)} \in \{2, 3\}$ on~$S$. If $W_S$ is not $M$-acyclic,
  let $C$ be a shortest directed cycle of $M$ inside $W_S$. Every vertex of $C$ has in-degree~$2$ in $M$, hence one cycle
  parent and at most one outside parent, and minimality plus TT-freeness make $C$ chordless with no shared outside parent
  on consecutive vertices; if some $i$ has $f(i-1) \in M(f(i))$ and $f(i) \in M(f(i-1))$, minimality forces $\card{C} = 2$
  (a digon of in-degree-$2$ vertices), while otherwise $C$ is a pure one-way chordless directed cycle of length $\ge 3$,
  both certified by $m = 2$ records (\lean{HalfRate.i5Piece}, as in Theorem~T of~\cite{SeventhPaper2026}).
\item \emph{Acyclic digon graph.} Assume $W_S$ is $M$-acyclic and $\mathcal{D}[S]$ is a forest. Fix a proper $2$-coloring
  $\kappa : S \to \{0, 1\}$ of $\mathcal{D}[S]$. Neither color class $S_k = \{v \in S : \kappa(v) = k\}$ can contain a
  directed cycle of $M$: any directed cycle inside $S_k$ cannot lie wholly in $W_S$ (since $W_S$ is $M$-acyclic), so it
  must visit some $v \in S_k$ with $\card{M(v)} = 3$, whose predecessor $u \in S_k$ on the cycle satisfies $v \in M(u)$ and
  hence $u \in M(v)$ by $\ClassS$, forming a monochromatic digon inside $S_k$. Thus each $S_k$ is $M$-acyclic, so printing
  $S_0$ and $S_1$ in topological order with $b = 1 - \kappa$ gives a global $2$-record certificate on $H = S$.
\item \emph{Shortest digon cycle and hub augmentation.} Finally, suppose $\mathcal{D}[S]$ contains cycles, and let $g \ge 4$
  be the girth of $\mathcal{D}[S]$.
  \begin{itemize}
  \item \emph{Subcase 4a (minimal chordal arc).} Suppose some $g$-cycle $c'$ of $\mathcal{D}[S]$ carries a non-cycle $M$-arc
    between vertices of cycle distance $d$; reversing $c'$ via \lean{cycRev} if necessary, we may assume $c'(i_0) \in
    M(c'(i_0 + d))$ with $2 \le d \le \lfloor g/2 \rfloor$, and choose $(d, c', i_0)$ minimizing~$d$ ($d = 2$ is ruled out by
    \lean{HalfRate.TTFree}, so $3 \le d$). Consider the $(d+1)$-vertex arc $H = \{p_0, \dots, p_d\}$ where $p_s = c'(i_0 + s)$.
    Minimality of $g$ makes $p_0 \in M(p_d)$ a \emph{pure one-way} arc ($p_d \notin M(p_0)$), so $\ClassS$ forces
    $\card{M(p_0)} \le 2$ ($p_1$ plus at most one outside parent), while every $p_s$ ($1 \le s \le d$) already has two
    internal parents in $H$ ($p_{s-1}$ and $p_{s+1}$ for $s < d$, or $p_{d-1}$ and $p_0$ for $s = d$) and hence at most one
    outside parent. Minimality of $d$ rules out all other internal $M$-arcs on $H$, and eliminating $[p_0, p_1, \dots, p_d]$
    in forward order with singleton records $P_s = \{p_s\}$ ($0 \le s \le d$), empty record $P_{d+1} = \emptyset$, and
    $b(p_s) = s + 1$ yields an explicit $(d+2)$-record nested certificate on~$H$ (\lean{minChordArc_pieceCert}).
  \item \emph{Subcase 4b (chordless non-stride-$3$ cycle).} Suppose no $g$-cycle of $\mathcal{D}[S]$ has any non-cycle
    $M$-chord. Pick any $g$-cycle $C = \{c(0), \dots, c(g-1)\}$. If $C$ is not a stride-$3$ cycle in context $\emptyset$
    under $M$, \cref{thm:DC} immediately certifies $H = C$.
  \item \emph{Subcase 4c (hub augmentation of a stride-$3$ cycle).} Suppose $C$ is a stride-$3$ cycle ($3 \mid g$, hence
    $g \ge 6$ by triangle-freeness) with hubs $w_0, w_1, w_2 \in S \setminus C$. Because $S$ is strongly connected and
    $w_0 \in S \setminus C$, some vertex $c(i_0) \in C$ sends an outgoing $M$-arc to some $y \in S \setminus C$, i.e.,
    $c(i_0) \in M(y)$. Since every vertex of $C$ has $\card{M(c(i))} = 3$ (two cycle neighbors plus $w_{i \bmod 3}$),
    reciprocal closure in $\ClassS$ forces $y \in M(c(i_0))$, so $y = w_{r_0}$ where $r_0 = i_0 \bmod 3$. Now $w_{r_0}$
    forms a digon with $c(i_0)$ and cannot form a digon with any other $c(k)$ ($k \ne i_0$), as two digons from $w_{r_0}$
    to $C$ would create a $\mathcal{D}$-cycle of length at most $\lfloor g/2 \rfloor + 2 < g$ (or a $4$-cycle when $g = 6$).
    Because $w_{r_0} \in M(c(i_0 + 3))$ and $i_0 + 3 \not\equiv i_0 \pmod g$, $c(i_0 + 3) \notin M(w_{r_0})$, so $\ClassS$
    forces $\card{M(w_{r_0})} = 2$: $w_{r_0}$ has parent $c(i_0)$ on $C$ and at most one parent outside $H := \{w_{r_0}\}
    \cup C$ (which may even be another hub $w_r$). Rotating $C$ so $i_0 = 0$ and $r_0 = 0$, the $(g+1)$-vertex augmented
    piece $H = \{w_0\} \cup C$ admits an explicit $(g+2)$-record sequential nested certificate
    (\lean{augmentedStride3_pieceCert}): list $H$ in cyclic order $[v_0, \dots, v_g] = [w_0, c(0), c(1), \dots, c(g-1)]$,
    assign own-record indices $(g, 1, 0, 3, 4, \dots, g-1, 0)$ with empty auxiliary record $P_{g+1} = \emptyset$, and set
    sequential supply $b(v_j) = j + 1$ for $0 \le j \le g$, verifying nestedness across shared outside parents via
    \cref{prop:PN}(1).
  \end{itemize}
\end{enumerate}
Applying \lean{HalfRate.rateBound_of_decomp} and \lean{MotionDepthCodes.card_sq_le_of_rateBound_half} completes~(2).
\end{outline}

For any vertex cut $S \subseteq V$ of $(V, N)$, write $I_N(S) = \{v \in V \setminus S : N(v) \subseteq S\}$
(\lean{isolatedOutDir}) for the set of vertices outside $S$ all of whose active recovery parents lie inside~$S$. Applying
\cref{thm:Sstar}(1) in context $U = S \cup I_N(S)$ yields the directed barrier bound for $\ClassS$:

\begin{corollary}[Directed barrier inequality for $\ClassS$]\label{cor:Bstar}
\formal{ClassS.lean}{isolatedOutDir, barrier_ineq_classS, log_card_le_barrier_classS}
Let $(V, N)$ be a finite recovery relation in $\ClassS$, and let $S \subseteq V$.
\begin{enumerate}
\item Every polymatroid $h$ valid for $N$ satisfies
  \[
    2 h(V) \le 2 \sum_{v \in S} h(\{v\}) + \sum_{v \in V \setminus (S \cup I_N(S))} h(\{v\}).
  \]
\item Every storage code $C \subseteq \alpha^V$ over a nonempty finite alphabet $\alpha$ satisfies
  \[
    2 \log \card{C} \le \bigl(\card{V} + \card{S} - \card{I_N(S)}\bigr) \log \card{\alpha}.
  \]
\end{enumerate}
\end{corollary}

\begin{outline}
Let $U = S \cup I_N(S)$. By \cref{thm:Sstar}(1), $V \setminus U$ admits a piece decomposition (\lean{HalfRate.Decomp} $N\ U\ \mathrm{pcs}$),
so \lean{HalfRate.decomp_ineq} gives $2(h(V) - h(U)) \le \sum_{v \notin U} (h(U \cup \{v\}) - h(U)) \le \sum_{v \notin U}
h(\{v\})$. Since every $v \in I_N(S)$ has $N(v) \subseteq S$, validity gives $h(U) = h(S) \le \sum_{v \in S} h(\{v\})$,
proving~(1) (\lean{barrier_ineq_classS}); evaluating on the uniform code-entropy polymatroid of $C$ proves~(2)
(\lean{log_card_le_barrier_classS}).
\end{outline}

Finally, we exhibit a strongly connected, completely triangle-free digraph in $\ClassS$ on which the low-in-degree,
$W_S$-cycle, and shortest-digon-cycle rules all stall simultaneously, whereas hub augmentation (Subcase~4c of
\cref{thm:Sstar}) decomposes it.

\begin{proposition}[Simultaneous stalling of earlier rules on $\Gto \in \ClassS$]\label{prop:G21}
\formal{ClassSExamples.lean}{g21Rel, g21_isClassS, g21_no_any_triangle, g21_min_indeg_and_acyclic_W, g21_stronglyConnected, g21_hex_no_pieceCert, g21_shortest_digon_cycles, rateBound_g21}
Let $\Gto$ (\lean{g21Rel}) be the recovery relation on $\{0, \dots, 20\}$ formed by three disjoint digon $6$-cycles
$A = \{0,\dots,5\}$, $B = \{6,\dots,11\}$, $C = \{12,\dots,17\}$ and three hubs $18, 19, 20$, where each $i \in \{0,\dots,17\}$
has cycle parents $i \pm 1$ (modulo~$6$ within its hexagon) plus hub parent $18 + (i \bmod 3)$, and the three hubs have
parent sets $N(18) = \{0, 15\}$, $N(19) = \{1, 7\}$, $N(20) = \{8, 14\}$. Then:
\begin{enumerate}
\item $\Gto \in \ClassS$, $\Gto$ is strongly connected, and the underlying undirected support of $\Gto$ is triangle-free;
\item every vertex of $\Gto$ has in-degree~$2$ or~$3$, and the in-degree-$2$ set $W = \{18, 19, 20\}$ has no internal arcs
  (so low-in-degree and $W$-cycle rules do not apply);
\item every cycle of length at most~$6$ in \lean{digonGraph}$(\Gto)$ has length~$6$ and vertex set $A$, $B$, or $C$ (each
  forming a mutual $6$-cycle in $\Gto$), and none of $A$, $B$, $C$ admits a piece certificate (\lean{HalfRate.PieceCert}
  $\Gto\ \emptyset\ H$) for any $m \in \N$;
\item nevertheless, $\Gto$ admits a piece decomposition and satisfies \lean{MotionDepthCodes.RateBound} $\Gto\ (1/2)$.
\end{enumerate}
\end{proposition}

\begin{outline}
Membership $\Gto \in \ClassS$ (\lean{g21_isClassS}), complete triangle-freeness (\lean{g21_no_any_triangle}), minimum
in-degree~$2$ and independence of $W = \{18,19,20\}$ (\lean{g21_min_indeg_and_acyclic_W}), and strong connectivity via an
explicit $32$-step closed walk visiting all $21$ vertices (\lean{g21_stronglyConnected}) are verified in the kernel.
Unbounded-$m$ impossibility on $A$, $B$, and $C$ (\lean{g21_hex_no_pieceCert}) follows by instantiating \cref{prop:CS}(1)
with hub triple $[18, 19, 20]$. To classify all digon cycles of length $3 \le g \le 6$ (\lean{g21_shortest_digon_cycles}),
a kernel-checked $4$-step BFS avoiding $w$ verifies that for every hub $w \in \{18,19,20\}$, its two digon neighbors lie in
distinct hexagons and have no $\mathcal{D}$-walk of length $\le 4$ avoiding~$w$, so no cycle of length $g \le 6$ can visit
any hub; projecting any remaining cycle into $\Z_6$ modulo~$6$ rules out $g \in \{3, 4, 5\}$ by parity and injectivity,
leaving $A$, $B$, and $C$. Finally, \cref{thm:Sstar} (peeling, for instance, a $7$-vertex hub-augmented piece such as
$\{18\} \cup A$ via Subcase~4c) yields \lean{MotionDepthCodes.RateBound} $\Gto\ (1/2)$ (\lean{rateBound_g21}).
\end{outline}

\section{Scalar Linear Codes at Degree~4 Without Heavy Separation}\label{sec:crux}

On $2$-head tape computations whose dependency DAG is cut to bounded depth, vertices with up to four active parents arise
naturally whenever two heads revisit cells across coupled scales~\cite[Section~8]{EighthPaper2026}. Over an arbitrary field
$F$, let $M \in F^{V \times V}$ have unit diagonal ($M_{v,v} = 1$) and no recovery triangle: there are no three distinct
vertices $a, b, c \in V$ with $b \in N_M(a)$, $c \in N_M(b)$ and $c \in N_M(a)$, where $N_M(v) = \{u \ne v : M_{v,u} \ne 0\}$
(\lean{MotionDepthCodes.ThmSHyp}). Recall from \cref{sec:prelim} that $v \in V$ is \emph{light} if it satisfies $(\star\star)$
(\lean{MotionDepthCodes.Star}: either $\card{N_M(v)} \le 2$, or $\card{N_M(v)} = 3$ and $v \in N_M(u) \implies u \in N_M(v)$)
and \emph{heavy} otherwise, with heavy set $H_M = \{v \in V : \neg(\star\star)(v)\}$ (\lean{BoundedNesting.heavyVerts}), and
that $(\star 4)$ (\lean{BoundedNesting.Star4}) requires $\card{N_M(v)} \le 3$, or $\card{N_M(v)} = 4$ with reciprocal
outgoing recovery arcs.

In~\cite[Section~8.2, Theorem~8.3]{EighthPaper2026}, the half-rate bound $\dim_F \ker M \le \lfloor \card{V}/2 \rfloor$ was
proved under $(\star 4)$ assuming either $\card{H_M} \le 2$ or \emph{both}~(S1) (\lean{BoundedNesting.SepS1}: no $v \in V$
has two distinct heavy recovery neighbors $p, q \in N_M(v) \cap H_M$) and~(S2) (\lean{BoundedNesting.SepS2}: no $h \in H_M$
has any heavy recovery neighbor in $N_M(h) \cap H_M$). On an undirected triangle-free support graph $G = (V, E)$ of maximum
degree~$4$, combining~(S1) and~(S2) forced heavy vertices to lie at graph distance at least~$3$ pairwise
(\lean{BoundedNesting.Far3}), ruling out adjacent heavy pairs. Notice, however, that on a triangle-free graph $G$, two
adjacent vertices $u, v \in H_M$ automatically share no common neighbor in~$G$. It is therefore natural to ask---as
discussed in~\cite[Section~8.2, Remark~8.8; Section~10, Problem~10.4]{EighthPaper2026}---whether~(S2) is merely an artifact
of the proof.

We prove that~(S2) can be dropped altogether, both in the directed matrix setting and in the undirected graph setting, and
deduce a strictly stronger matching-deficiency bound.

\subsection{Theorem~SH\texorpdfstring{$^*$}{*}: eliminating \texorpdfstring{(S2)}{(S2)}}\label{sec:crux-shstar}

For a subset $W \subseteq V$, write $\mathrm{rs}_W(v) = N_M(v) \cap W$ (\lean{BoundedNesting.rs}), let $(\star\star)_W$,
$(\star 4)_W$, and $H_{M,W}$ (\lean{BoundedNesting.StarOn}, \lean{BoundedNesting.Star4On}, \lean{BoundedNesting.heavySet})
be the restrictions of $(\star\star)$, $(\star 4)$, and $H_M$ to $W$, and write $\mathrm{S1}_W$ (\lean{BoundedNesting.S1On})
for the condition that no $v \in W$ has two distinct recovery neighbors in $\mathrm{rs}_W(v) \cap H_{M,W}$. Define
\lean{SHStarOn}$(M, W)$ to hold when every $v \in W$ satisfies $(\star 4)_W$ and either $\card{H_{M,W}} \le 2$ or
$\mathrm{S1}_W$ holds. For an undirected graph $G = (V, E)$, write \lean{NoCommonNbhd}$(G, u, v)$ if there is no $w \in V$
adjacent in $G$ to both $u$ and~$v$.

\begin{theorem}[Half-rate kernel bound under (S1) alone]\label{thm:SHstar}
\formal{TheoremSHStar.lean}{SHStarOn, F3_star, odd_state_star, MinCE.false_SHStar, finrank_subCode_le_half_SHStar, theoremSHStar, theoremSHStar_transpose, NoCommonNbhd, theoremSHStar_graph}
Let $F$ be a field, $V$ a finite vertex set, and $M \in F^{V \times V}$.
\begin{enumerate}
\item \emph{(Hereditary subcode form.)} Suppose $M$ satisfies \lean{MotionDepthCodes.ThmSHyp} (unit diagonal and no recovery
  triangle). For every $W \subseteq V$ satisfying \lean{SHStarOn}$(M, W)$, the restricted code $\mathrm{subCode}_M(W) =
  \{x \in F^V : \mathrm{supp}(x) \subseteq W \text{ and } (Mx)|_W = 0\}$ satisfies
  \[
    \dim_F \mathrm{subCode}_M(W) \le \lfloor \card{W} / 2 \rfloor.
  \]
\item \emph{(Global matrix and transpose forms.)} Suppose $M$ satisfies \lean{MotionDepthCodes.ThmSHyp}. If every $v \in V$
  satisfies $(\star 4)$ and either $\card{H_M} \le 2$ or \lean{BoundedNesting.SepS1}$(M)$ holds, then
  \[
    \dim_F \ker M \le \lfloor \card{V} / 2 \rfloor;
  \]
  the same bound holds if $M^\T$ satisfies $(\star 4)$ everywhere and either $\card{H_{M^\T}} \le 2$ or
  \lean{BoundedNesting.SepS1}$(M^\T)$ holds.
\item \emph{(Graph form.)} Let $G = (V, E)$ be a triangle-free simple graph of maximum degree at most~$4$, and suppose
  $M_{v,v} = 1$ for all $v \in V$ and $M_{v,u} = 0$ whenever $v \ne u$ and $\{v, u\} \notin E$. If either $\card{H_M} \le 2$
  or every pair of distinct $u, v \in H_M$ satisfies \lean{NoCommonNbhd}$(G, u, v)$, then
  \[
    \rank_F M \ge \lceil \card{V} / 2 \rceil \qquad\text{and}\qquad \dim_F \ker M \le \lfloor \card{V} / 2 \rfloor.
  \]
\end{enumerate}
\end{theorem}

\begin{outline}
First observe that \lean{SHStarOn}$(M, W)$ is hereditary under passing to subsets $W' \subseteq W$ (\lean{SHStarOn.mono}),
since $(\star 4)_{W'}$ and membership in $H_{M,W'}$ are monotone in~$W$. Thus~(1) (\lean{finrank_subCode_le_half_SHStar})
reduces by strong induction on $\card{W}$ to ruling out a minimal counterexample (\lean{BoundedNesting.MinCE} $M\ W$) on
which \lean{SHStarOn}$(M, W)$ holds (\lean{MinCE.false_SHStar}). Recall from~\cite[Section~8.2]{EighthPaper2026} that a
minimal counterexample $W$ has odd cardinality $\card{W} = 2k + 1$, $\mathrm{cdim}_W(\emptyset) = k + 1$, every proper
subset $W' \subsetneq W$ satisfies $\dim_F \mathrm{subCode}_M(W') \le \lfloor \card{W'}/2 \rfloor$, and every $v \in W$ has
$\card{\mathrm{rs}_W(v)} \ge 2$. For a state $X \subseteq W$ from start vertex $v_0 \in X$, write
$\mathrm{pend}_W(X, u) = \mathrm{rs}_W(u) \setminus X$ and $\mathrm{cdim}_W(X) = \dim_F(\mathrm{subCode}_M(W) \cap
\mathrm{zeroOn}(X))$.

Where did~\cite{EighthPaper2026} use~(S2)? Precisely in two places: Phase~E (when an odd state $Z$ expands along a heavy
vertex $h \in Z$ with $\card{\mathrm{pend}_W(Z, h)} = 3$) and Case~3 of the starting split (when choosing a light neighbor
$v_0$ of a heavy vertex $h$). Notice that Phases~F1 and~F2 of~\cite{EighthPaper2026} (\lean{BoundedNesting.F1},
\lean{BoundedNesting.F2}) already operate under \lean{BoundedNesting.F2Mode}, which requires only either at most one
remaining non-closed heavy vertex or $\mathrm{S1}_W$.

\emph{Step~1: Unified Phase~E (\lean{F3_star}).} Suppose $Z \subseteq W$ has odd size, $h \in Z$, $P = \mathrm{pend}_W(Z, h)$
has $\card{P} = 3$, $X_0 = Z \cup P$ is an even good state with neutral dimension drop $\mathrm{cdim}_W(Z) = 2 +
\mathrm{cdim}_W(X_0)$, and \emph{at most one} member of $P$ violates $(\star\star)_W$. (When $\mathrm{S1}_W$ holds, $P
\subseteq \mathrm{rs}_W(h)$ cannot contain two distinct heavy vertices by $\mathrm{S1}_W$ applied at $h$ itself; when
$\card{H_{M,W}} \le 2$, $h \in Z \cap H_{M,W}$ is already consumed, leaving at most one heavy vertex outside~$Z$.) By
recovery-triangle-freeness on $h$, $P$ is independent in $\mathrm{rs}_W$, and immediate closure (\lean{BoundedNesting.imm_closure})
ensures $\mathrm{pend}_W(X_0, t) \ne \emptyset$ for all $t \in P$.
\begin{itemize}
\item If $\card{W \setminus X_0} = 1$, say $W \setminus X_0 = \{w\}$, then every $t \in P$ has $\mathrm{pend}_W(X_0, t) =
  \{w\}$; taking a test vector $x \in \mathrm{subCode}_M(W)$ vanishing on $Z \cup \{t_1\}$ with $x_{t_2} \ne 0$
  (\lean{BoundedNesting.test_vec}), row $t_1$ forces $x_w = 0$ and row $t_2$ then forces $x_{t_2} = 0$, contradiction.
\item Suppose $\card{W \setminus X_0} \ge 3$ (recall $\card{W} - \card{X_0}$ is odd). Every light $t \in P$ has
  $\card{\mathrm{pend}_W(X_0, t)} \le 2$ (and $\ne 1$ by Phase~F1), hence $\card{\mathrm{pend}_W(X_0, t)} = 2$. If all three
  members of $P$ are light, Phase~F2 equates their $2$-element pending sets to a common pair $\{a, b\}$, so $P \cup
  \bigcup_{t \in P} \mathrm{rs}_W(t) \subseteq P \cup \{h, a, b\}$ has size at most~$6 = 2 \card{P}$, violating Hall's
  count~\cite{Hall1935} on minimal counterexamples (\lean{BoundedNesting.MinCE.hall_count}).
\item Otherwise $P = \{u_1, u_2, g\}$ where $g$ is the unique heavy member of $P$ and $u_1, u_2$ are light. Both $u_1$ and
  $u_2$ have $2$-element pending sets outside $X_0$, which coincide by Phase~F2: $\mathrm{pend}_W(X_0, u_1) =
  \mathrm{pend}_W(X_0, u_2) = \{a, b\}$ ($a \ne b$). Making the neutral move from $X_0$ along $u_1$ gives the exact drop
  $\mathrm{cdim}_W(X_0) = 1 + \mathrm{cdim}_W(X_0 \cup \{a, b\})$. If every $y \in \mathrm{subCode}_M(W) \cap
  \mathrm{zeroOn}(X_0)$ had $y_a = 0$, row $u_1$ ($M_{u_1, a} y_a + M_{u_1, b} y_b = 0$ with $M_{u_1, b} \ne 0$) would also
  force $y_b = 0$, contradicting $\mathrm{cdim}_W(X_0) > \mathrm{cdim}_W(X_0 \cup \{a, b\})$. Thus there exists $y \in
  \mathrm{subCode}_M(W) \cap \mathrm{zeroOn}(X_0)$ with $y_a \ne 0$ satisfying $M_{u_i, a} y_a + M_{u_i, b} y_b = 0$ for
  $i = 1, 2$. Meanwhile, \lean{BoundedNesting.test_vec} on $(u_2, g)$ gives $x \in \mathrm{subCode}_M(W)$ vanishing on
  $Z \cup \{u_1\}$ with $x_{u_2} \ne 0$, so $M_{u_1, a} x_a + M_{u_1, b} x_b = 0$ while $x_{u_2} + M_{u_2, a} x_a +
  M_{u_2, b} x_b = 0$, contradicting proportional two-row elimination (\lean{BoundedNesting.two_rows_zero}).
\end{itemize}

\emph{Step~2: Odd-state lemma and minimal counterexample exclusion (\lean{odd_state_star}, \lean{MinCE.false_SHStar}).}
Strong induction on the remaining non-closed heavy count $t = \card{\mathrm{NC}_W(X)}$ and $\card{W \setminus X}$ proves
\lean{BoundedNesting.OSAt}$(M, W, v_0, t)$ for all $t$: an odd state $X \subsetneq W$ has no pending singleton, preserves
parity under pending pairs, and on a pending triple $h \in X$ drops $\card{\mathrm{NC}_W}$ and hits contradiction by
\lean{F3_star}. Finally, on $W$ itself: if some $v_0 \in W$ has $\card{\mathrm{rs}_W(v_0)} = 2$ (Case~1), $X_0 = \{v_0\}
\cup \mathrm{rs}_W(v_0)$ is an odd state of size~$3$; if some heavy $h \in W$ has $\card{\mathrm{rs}_W(h)} = 3$ (Case~2),
expanding $\{h\}$ invokes \lean{F3_star} directly; otherwise (Case~3), every light vertex has degree~$3$, every heavy
vertex has degree~$4$, and $(\star\star)_W$ plus $(\star 4)_W$ make $\mathrm{rs}_W$ a symmetric simple graph on $W$. Since
$\card{W}$ is odd and the number of odd-degree (light) vertices in any simple graph is even, $\card{H_{M,W}}$ is odd
(hence $\ge 1$, and $= 1$ whenever $\card{H_{M,W}} \le 2$, so $\mathrm{S1}_W$ holds unconditionally). Picking $h \in
H_{M,W}$, $\mathrm{S1}_W$ at $h$ allows at most one heavy neighbor among the four elements of $\mathrm{rs}_W(h)$, leaving
at least three light neighbors $v_0 \in \mathrm{rs}_W(h)$. Starting at $v_0$, $\mathrm{S1}_W$ at $v_0$ guarantees that the
other two neighbors $s_1, s_2 \in \mathrm{rs}_W(v_0) \setminus \{h\}$ are light, and the exact Case~3 closure of
\cite{EighthPaper2026} contradicts \lean{BoundedNesting.F1} or \lean{odd_state_star}.

For~(2) and~(3), global instantiation $W = V$ gives \lean{theoremSHStar}; transpose invariance $\dim_F \ker M^\T =
\dim_F \ker M$ gives \lean{theoremSHStar_transpose}; and on a triangle-free graph $G$ of maximum degree~$4$,
\lean{NoCommonNbhd} on $H_M$ implies \lean{BoundedNesting.SepS1}, so~(2) bounds $\dim_F \ker M \le \lfloor \card{V}/2 \rfloor$
and rank-nullity $\rank_F M + \dim_F \ker M = \card{V}$ gives \lean{theoremSHStar_graph}.
\end{outline}

\subsection{Corollary~M\texorpdfstring{$^*$}{*}: heavy-matching deficiency bound}\label{sec:crux-matching}

For a simple graph $G = (V, E)$ and subset $H \subseteq V$, define the \emph{common-neighbor-free packing number}
(\lean{rhoMatch}) by
\begin{equation}\label{eq:rhom}
  \rhom(G, H) := \max\bigl\{\card{D} : D \subseteq H \text{ and } \forall u, v \in D,\; u \ne v \implies
  \lean{NoCommonNbhd}(G, u, v)\bigr\}.
\end{equation}
On a triangle-free graph $G$, two vertices $u \ne v$ satisfy \lean{NoCommonNbhd}$(G, u, v)$ if and only if
$\dist_G(u, v) \ne 2$, i.e., either $\{u, v\} \in E$ or $\dist_G(u, v) \ge 3$. Consequently, an admissible subset $D
\subseteq H$ is any disjoint union of edges and isolated vertices in $G[H]$ whose connected components lie at pairwise
$G$-distance at least~$3$. In particular, every distance-$3$ packing (\lean{BoundedNesting.Far3}) is admissible, so
$\rhom(G, H) \ge \rho_3(G, H)$ always holds, and whenever $H$ spans an induced matching of $r$ edges in $G$ at pairwise
distance at least~$3$, $\rhom(G, H) = 2r = 2 \rho_3(G, H)$.

\begin{corollary}[Barrier and matching bounds with $\rhom$]\label{cor:Mstar}
\formal{TheoremSHStar.lean}{finrank_ker_le_barrier_SHStar, finrank_ker_add_card_le_matching_SHStar, rhoMatch, exists_rhoMatch, rho3_le_rhoMatch, finrank_ker_le_matching_rhoMatch, finrank_ker_le_matching_rhoMatch_transpose}
Let $F$ be a field and $V$ be a finite vertex set.
\begin{enumerate}
\item \emph{(Barrier form.)} Suppose $M \in F^{V \times V}$ satisfies \lean{MotionDepthCodes.ThmSHyp}, $U \subseteq V$, and
  $\mathcal{P}$ is a partition of $V \setminus U$ into pairwise disjoint blocks such that $M_{u,v} = 0$ between distinct
  blocks $B \ne B'$ of $\mathcal{P}$ and \lean{SHStarOn}$(M, B)$ holds for every $B \in \mathcal{P}$. Then
  \[
    \dim_F \ker M \le \card{U} + \sum_{B \in \mathcal{P}} \lfloor \card{B} / 2 \rfloor.
  \]
\item \emph{(Matching-packing form.)} Let $G = (V, E)$ be a triangle-free simple graph of maximum degree at most~$4$ with
  matching number $\nu(G)$, and suppose $M \in F^{V \times V}$ has unit diagonal and off-diagonal support in~$G$. Then
  $\rho_3(G, H_M) \le \rhom(G, H_M)$ and
  \[
    \dim_F \ker M + \max\bigl(\rhom(G, H_M),\, \min(\card{H_M}, 2)\bigr) \le \nu(G) + \card{H_M},
  \]
  and the same inequality holds with $H_{M^\T}$ in place of $H_M$.
\end{enumerate}
\end{corollary}

\begin{outline}
For~(1) (\lean{finrank_ker_le_barrier_SHStar}), intersecting $\ker M$ with $\mathrm{zeroOn}(U)$ loses at most $\card{U}$
dimensions and embeds into $\bigoplus_{B \in \mathcal{P}} \mathrm{subCode}_M(B)$ because $M$ has no cross-block entries;
applying \cref{thm:SHstar}(1) to each block $B \in \mathcal{P}$ gives $\sum_{B \in \mathcal{P}} \lfloor \card{B}/2 \rfloor$.

For~(2) (\lean{finrank_ker_add_card_le_matching_SHStar} and \lean{finrank_ker_le_matching_rhoMatch}), let $D \subseteq H_M$
satisfy either $\card{D} \le 2$ or pairwise \lean{NoCommonNbhd} on $(G, u, v)$, and set $U = H_M \setminus D$. By the
Tutte--Berge formula~\cite{Berge1958} (Lean-verified in \texttt{lean-machine} as
\lean{SimpleGraph.exists_card_add_sum_card_div_two_eq_matchingNumber}), there is a Tutte barrier $U' \subseteq V$ such that
$\card{U'} + \sum_{B \in \mathrm{comp}(G - U')} \lfloor \card{B}/2 \rfloor = \nu(G)$. On each punctured component $B' =
B \setminus U$ of $V \setminus (U \cup U')$, every heavy vertex of $M$ in $B'$ belongs to $D$, so \lean{SHStarOn} holds on
$(M, B')$. Applying~(1) with cut $U \cup U'$ gives
\[
  \dim_F \ker M \le \card{U \cup U'} + \sum_{B} \lfloor \card{B \setminus U}/2 \rfloor
  \le \card{U} + \nu(G) = \card{H_M} - \card{D} + \nu(G).
\]
Choosing $D$ to attain $\rhom(G, H_M)$ (\lean{exists_rhoMatch}) or $\min(\card{H_M}, 2)$ completes the proof, and
transpose invariance of $\dim_F \ker M$ gives \lean{finrank_ker_le_matching_rhoMatch_transpose}.
\end{outline}

\begin{remark}[Linear-size separation between $\rhom$ and $\rho_3$]\label{rem:deg4-numerics}
For each integer $k \ge 2$, construct a connected bipartite (hence triangle-free) graph $G_{8k}$ on $8k$ vertices with
bipartition $A \sqcup B$, where $A = \bigsqcup_{r=0}^{k-1} \{a_{r,0}, a_{r,1}, a_{r,2}, a_{r,3}\}$ and
$B = \bigsqcup_{r=0}^{k-1} \{b_{r,0}, b_{r,1}, b_{r,2}, b_{r,3}\}$, by including, for each $0 \le r < k$, all seven edges
incident to $a_{r,0}$ or $b_{r,0}$ inside block~$r$ ($\{a_{r,0}, b_{r,j}\}$ and $\{a_{r,j}, b_{r,0}\}$ for $0 \le j \le 3$)
together with a $2$-regular bipartite subgraph on the remaining $6k$ vertices $\{a_{r,j}, b_{r,j} : 0 \le r < k,\; 1 \le j \le 3\}$
having edges $\{a_{r,j}, b_{r,j}\}$ and $\{a_{r,j}, b_{(r+1)\bmod k,\,(j\bmod 3)+1}\}$ (a single $6k$-cycle when $3 \nmid k$,
and three disjoint $2k$-cycles when $3 \mid k$). Then $G_{8k}$ is $4$-regular on $H = \{a_{r,0}, b_{r,0} : 0 \le r < k\}$
($\card{H} = 2k$), $3$-regular on $V \setminus H$, and has a perfect matching $\nu(G_{8k}) = 4k$ (verified numerically for
$2 \le k \le 6$ by \leanfile{tools/reproduce_numerics.py}). Because $G_{8k}$ is bipartite and the neighborhoods
$N(a_{r,0}) = \{b_{r,0},\dots,b_{r,3}\}$ and $N(b_{r,0}) = \{a_{r,0},\dots,a_{r,3}\}$ are pairwise disjoint across $r$, no two
distinct vertices of $H$ share a common neighbor in $G_{8k}$, whereas each pair $\{a_{r,0}, b_{r,0}\}$ is an edge:
\begin{itemize}
\item the eighth paper's distance-$3$ packing invariant satisfies $\rho_3(G_{8k}, H) = k$, yielding only the upper bound
  $\dim_F \ker M \le 4k + 2k - k = 5k = \frac{5}{8}\card{V}$;
\item by contrast, $\rhom(G_{8k}, H) = 2k = \card{H}$, so \cref{thm:SHstar} and \cref{cor:Mstar} certify the sharp half-rate
  bound $\dim_F \ker M \le 4k = \frac{1}{2}\card{V}$ over every field~$F$; over $\F_2$, the all-ones unit-diagonal matrix
  supported on $G_{8k}$ has heavy set $H$ and attains $\dim_{\F_2} \ker M = 4k$ (verified for $2 \le k \le 6$ by
  \leanfile{tools/reproduce_numerics.py}).
\end{itemize}
\end{remark}

\section{Deep-Nesting Stroke Laws, Super-Fibonacci Unwrapping, and Block Pairs}\label{sec:dnp}

While \cref{sec:reciprocity,sec:crux} push half-rate coding barriers across $\ClassS$ and degree-$4$ graphs under~(S1),
completing the conditional impossibility program of~\cite{FourthPaper2026,EighthPaper2026} for $2$-head tapes also
requires understanding when $2$-head walk dependency DAGs admit $\eps T$-edge cuts to bounded depth ($\TCP(2)$). The
canonical obstruction testbed recalled in~\cite[Section~2.5]{EighthPaper2026} (after~\cite[Definitions~5.3 and~5.9]{FourthPaper2026})
is the \emph{deep-nesting pair} $\DNP(s_0, m)$, running two copies of the recursive turn-and-revisit walk $U(s_0, m, 1)$ at
relative time shift $\xi \in \{0, \dots, P-1\}$ over a shared tape (\cref{sec:prelim}). On a single copy of $U(s_0, m, 1)$,
a turn-supported label $\Phi : \Z \to \Z_q$ cuts depth to $2^{m + O(1)}$ by keeping $\Phi(t+1) - \Phi(t) \equiv 1 \pmod q$
on plain steps and adjusting $\Phi$ at turns so that same-cell revisit hops satisfy $\Phi(u) - \Phi(p) \equiv 1 \pmod q$.
On two shifted copies, however, a revisit hop $(p, u)$ of head~$1$ is \emph{interrupted} (non-blind) whenever head~$2$
executes a turn in $(p, u]$, adding head~$2$'s net potential increment $F$ across $(p, u]$ to the mod-$q$ consistency
relation.

In~\cite[Section~5]{EighthPaper2026}, the path-sum law, band decomposition, block-label family $B(a, K)$, and straddle
involution were developed on paper without Lean formalization, and Conjecture~A$''$~\cite[Section~5.6, Conjecture~5.24]{EighthPaper2026}---asserting
that at some odd shift $\xi$, every mod-$q$ label of cost at most $n_{\mathrm{h}}/8$ on the $r$-fold cover has at most
$C\log_2 q + C' + \eta(m+1)$ doubly $\theta_0$-served levels ($\eta < 1$)---left two explicit steps open at the end
of~\cite[Section~5.6]{EighthPaper2026} even for pairs of block labels:
\begin{enumerate}
\item \emph{``a bound on $\abs{F}$ on the straddling hops, or a version of Theorem~V that counts the window directly''}
  (since $\abs{F}$ in~\cite[Theorem~5.20]{EighthPaper2026} can grow linearly in~$m$ and wrap modulo~$q$ unless $q$ is of
  order $(m - \kappa)Q$); and
\item \emph{``at one fixed shift for all coarse classes, equidistribution of a class's apex centres relative to the other
  head's canceller valleys of level about $k + \log_2(m - k)$, a mixed-radix carry problem''}.
\end{enumerate}

In this section we formalize the turn-potential, V-straddle, and super-Fibonacci apparatus in Lean~4: \cref{thm:Vstar}
bypasses step~(1) altogether for level-uniform labels by summing straddling pairs directly modulo~$q$ with zero wrap-around
loss, \lean{exists_oddShift_coarse_straddle} establishes a weaker one-direction $\tfrac{1}{16}$-fraction form of step~(2)
over odd shifts, and \cref{thm:A} combines these with \cref{thm:LU} to produce three-regime failure witnesses for all
block-label pairs $(B(a_1, K_1), B(a_2, K_2))$ of bounded block length at two fixed odd shifts (a single shift for
symmetric pairs).

\subsection{Exact stroke law and \texorpdfstring{$O(m)$}{O(m)} recursive turn potential}\label{sec:dnp-strokes}

Fix integers $0 \le m \le s_0$, dyadic scales $R_i = 2^{s_0 - i}$ (\lean{scaleR}), stroke lengths $\Pi_\kappa =
(m - \kappa + 1) R_\kappa$ (\lean{strokeLen}), and fundamental period $P = 4 \Pi_0 = 4(m+1)2^{s_0}$ (\lean{dnpPeriod}) of
$U(s_0, m, 1)$~\cite[Definition~5.3]{FourthPaper2026},~\cite[Section~2.5]{EighthPaper2026}. A \emph{level-uniform turn-supported potential} on
$U(s_0, m, 1)$ is specified by integer level parameters $A, E : \N \to \Z$ (or equivalently back-stroke and forward-stroke
path-sum increments $c_{\sigb}(\kappa) = 1 + A_\kappa$ and $c_{\sigf}(\kappa) = 2 + A_\kappa + E_\kappa$, \lean{pathSumBack}
and \lean{pathSumFwd}) together with an outer tip jump $j_{\mathrm{tip}} \in \Z$ at the period reversals $\Pi_0$ and
$3\Pi_0$: at every internal apex of level $\kappa$ ($1 \le \kappa \le m$), the turn jump adds $c_{\sigb}(\kappa)$; at every
internal valley of level $\kappa$, the turn jump adds $c_{\sigf}(\kappa) - c_{\sigb}(\kappa) - M_{\sigf}(\kappa)$, where
$M_{\sigf}(\kappa) = \sum_{l = \kappa + 1}^m c_{\sigf}(l)$ (\lean{strokeMass}) is the cumulative turn jump inside one
level-$\kappa$ stroke, arranged so that entering the level-$\kappa$ forward stroke starts at cumulative offset
$c_{\sigf}(\kappa)$ and completing the entire level-$\kappa$ V contributes net turn jump
$c_{\sigf}(\kappa) + M_{\sigf}(\kappa) = M_{\sigf}(\kappa - 1)$.

\begin{proposition}[Exact stroke law and recursive turn potential]\label{prop:SL}
\formal{DNPStrokes.lean}{scaleR, strokeLen, dnpPeriod, decompDepth_valid, reconstruct_decomp, decomp_reconstruct, card_stroke_stopDepth_branchWord, card_stroke_stopDepth, card_stroke_stopDepth_branchWord_parity, card_stroke_stopDepth_parity, strokePhiRec, strokePhiRec_back, strokePhiRec_fwd, evalPhiRec, evalPhiRec_eq_sum_turnJumps}
Let $0 \le \kappa \le m \le s_0$.
\begin{enumerate}
\item \emph{(Stroke decomposition bijection.)} Recursive stroke decomposition (\lean{decompDepth}) and reconstruction
  (\lean{reconstructDepth}) establish an explicit bijection between offsets $o \in \{0, \dots, \Pi_\kappa - 1\}$ of a
  level-$\kappa$ stroke and valid triples $(s, w, \mathrm{rem})$ (\lean{ValidStrokeCode}) consisting of a stopping depth
  $s \in \{0, \dots, m - \kappa\}$, a nested branch word $w \in \{\sigb, \sigf\}^s$ ($\mathrm{false} = \sigb$,
  $\mathrm{true} = \sigf$), and a plain-segment remainder $\mathrm{rem} \in \{0, \dots, R_{\kappa + s} - 1\}$.
\item \emph{(Uniform stopping depth, fair-coin branches, and parity independence.)} For each $0 \le s \le m - \kappa$ and
  each word $w \in \{\sigb, \sigf\}^s$, there are exactly $R_{\kappa + s} = R_\kappa / 2^s$ offsets $o \in [0, \Pi_\kappa)$
  with code $(s, w, \cdot)$, and summing over $w \in \{\sigb, \sigf\}^s$ gives \emph{exactly} $R_\kappa = \Pi_\kappa /
  (m - \kappa + 1)$ offsets at stopping depth~$s$. Moreover, when $m < s_0$ (so every $R_i$ for $i \le m$ is even),
  $o \equiv \mathrm{rem} \pmod 2$, and for each parity $p \in \{0, 1\}$, stopping depth $s$, and word $w \in \{\sigb,
  \sigf\}^s$, the joint count is $R_{\kappa + s}/2$ and the marginal count at $(s, p)$ is $R_\kappa / 2$.
\item \emph{($O(m)$ recursive turn potential and V-stroke decomposition.)} For any $c_{\sigb}, c_{\sigf} : \N \to \Z$ and
  $j_{\mathrm{tip}} \in \Z$, let $Y_i(o) = {}$\lean{strokePhiRec}$(s_0, m, c_{\sigb}, c_{\sigf}, i, o)$ be the $O(m - i)$
  recursive stroke potential at level $i \le m$. The $O(m)$-step period evaluator $\Phi(t) =
  {}$\lean{evalPhiRec}$(s_0, m, c_{\sigb}, c_{\sigf}, j_{\mathrm{tip}}, t)$ equals the exact cumulative sum of turn jumps
  $\sum_{0 < \tau \le t} J(\tau)$ (\lean{periodTurnJump}) for all $0 \le t < P$. Inside any enclosing level-$(\kappa - 1)$
  stroke ($1 \le \kappa \le m$), its internal level-$\kappa$ V at apex offset $R_{\kappa - 1}$ satisfies, for all
  $0 \le o < \Pi_\kappa$,
  \begin{equation}\label{eq:v-decomp}
    Y_{\kappa - 1}(R_{\kappa - 1} + o) = c_{\sigb}(\kappa) + Y_\kappa(o), \qquad
    Y_{\kappa - 1}(R_{\kappa - 1} + \Pi_\kappa + o) = c_{\sigf}(\kappa) + Y_\kappa(o).
  \end{equation}
\end{enumerate}
\end{proposition}

\begin{outline}
Part~(1) (\lean{decompDepth_valid}, \lean{reconstruct_decomp}, \lean{decomp_reconstruct}) inducts on remaining depth
$d = m - \kappa$ using the stroke-length identity $\Pi_\kappa = R_\kappa + 2 \Pi_{\kappa+1}$ (\lean{strokeLen_succ}),
partitioning $[0, \Pi_\kappa)$ into the initial plain segment $[0, R_\kappa)$ ($s = 0$), the back sub-stroke $[R_\kappa,
R_\kappa + \Pi_{\kappa+1})$ ($b_1 = \sigb$), and the forward sub-stroke $[R_\kappa + \Pi_{\kappa+1}, \Pi_\kappa)$
($b_1 = \sigf$). Part~(2) (\lean{card_stroke_stopDepth_branchWord}, \lean{card_stroke_stopDepth},
\lean{card_stroke_stopDepth_branchWord_parity}, and \lean{card_stroke_stopDepth_parity}) counts fibers via
\texttt{Finset.card\_bij}; when $m < s_0$, $2 \mid R_m$ divides every summand $R_i$ and $\Pi_i$ added by
\lean{reconstructDepth}, so $o \bmod 2 = \mathrm{rem} \bmod 2$ and $[0, R_{\kappa+s})$ splits into $R_{\kappa+s}/2$ even
and $R_{\kappa+s}/2$ odd remainders. For~(3), unfolding one step of \lean{strokePhiRec} yields~\eqref{eq:v-decomp}
(\lean{strokePhiRec_back}, \lean{strokePhiRec_fwd}), induction on $m - \kappa$ equates \lean{strokePhiRec} with the sum of
\lean{strokeTurnJump} over $(0, o]$ (\lean{strokePhiRec_eq_sum_turnJumps}) by cancelling $-M_{\sigf}(\kappa)$ at the valley
turn against the completed back-stroke jump sum $M_{\sigf}(\kappa)$, and summing across the four level-$0$ strokes of
$[0, P)$ establishes \lean{evalPhiRec_eq_sum_turnJumps}.
\end{outline}

\subsection{Unconditional mod-\texorpdfstring{$q$}{q} V-straddle bound and odd-center straddle sums}\label{sec:dnp-straddle}

Consider a coarse revisit hop $(c_0 - d, c_0 + d)$ of half-width $d \in \{1, \dots, R_k - 1\}$ centered at an apex mapped
to time $c_0 \in \Z$ on the other head's timeline, and suppose the other head executes a level-$\kappa$ V on $[\alpha_V,
\alpha_V + 2\Pi_\kappa)$ with $\Pi_\kappa \le R_k$. We say $d$ \emph{straddles} this V (\lean{straddleSet}) if its left
endpoint lands in the back stroke and its right endpoint lands in the forward stroke:
\begin{equation}\label{eq:straddle-set}
  d \in \mathrm{Strad}(c_0, \alpha_V, \Pi_\kappa, R_k) \iff
  1 \le d \le R_k - 1,\quad \alpha_V \le c_0 - d < \alpha_V + \Pi_\kappa \le c_0 + d < \alpha_V + 2\Pi_\kappa.
\end{equation}
Notice the reflection symmetry $d \mapsto \Pi_\kappa - d$ (\lean{mem_straddleSet_sub_iff}): when $\Pi_\kappa \le R_k$,
$d \in \mathrm{Strad}(c_0, \alpha_V, \Pi_\kappa, R_k)$ if and only if $\Pi_\kappa - d \in \mathrm{Strad}(c_0, \alpha_V,
\Pi_\kappa, R_k)$, because the back-stroke offset $o_p(d) = (c_0 - d) - \alpha_V$ and forward-stroke offset $o_u(d) =
(c_0 + d) - (\alpha_V + \Pi_\kappa)$ swap under $d \leftrightarrow \Pi_\kappa - d$:
\[
  o_p(\Pi_\kappa - d) = (c_0 - \Pi_\kappa + d) - \alpha_V = o_u(d), \qquad
  o_u(\Pi_\kappa - d) = (c_0 + \Pi_\kappa - d) - (\alpha_V + \Pi_\kappa) = o_p(d).
\]

\begin{theorem}[Unconditional mod-$q$ V-straddle inequality and parity-invariant straddle count]\label{thm:Vstar}
\formal{DNPStraddle.lean}{straddleSet, mem_straddleSet_sub_iff, HasLevelUniformVDecomp, strokePhiRec_hasLevelUniformVDecomp, straddle_otherHeadInc_sum, theoremVStar_pair, theoremVStar_sum, sum_straddleSet_card_parity}
Let $\Phi : \Z \to \Z$ admit a level-uniform V decomposition on $[\alpha_V, \alpha_V + 2\Pi_\kappa)$ with base constant
$C_V \in \Z$ and parameters $(c_{\sigb}, c_{\sigf}, Y)$, i.e., $\Phi(\alpha_V + o) = C_V + c_{\sigb} + Y(o)$ and
$\Phi(\alpha_V + \Pi_\kappa + o) = C_V + c_{\sigf} + Y(o)$ for $0 \le o < \Pi_\kappa$ (\lean{HasLevelUniformVDecomp};
instantiated for \lean{strokePhiRec} via~\eqref{eq:v-decomp} by \lean{strokePhiRec_hasLevelUniformVDecomp}), and write
$F(d) = \Phi(c_0 + d) - \Phi(c_0 - d)$ (\lean{otherHeadInc}) for the other head's net potential increment across
$(c_0 - d, c_0 + d)$.
\begin{enumerate}
\item \emph{(Exact $\Z$-cancellation.)} For every $d \in \mathrm{Strad}(c_0, \alpha_V, \Pi_\kappa, R_k)$,
  \[
    F(d) + F(\Pi_\kappa - d) = 2(c_{\sigf} - c_{\sigb}) \quad\text{in }\Z.
  \]
\item \emph{(Unconditional mod-$q$ pair and set bounds.)} Let $q > 0$ and $A' \in \Z$, and define the reduced mod-$q$
  interrupted hop error $e(d) = (A' + F(d)) \bmod q \in \{0, \dots, q - 1\}$. For every $d \in S := \mathrm{Strad}(c_0,
  \alpha_V, \Pi_\kappa, R_k)$,
  \begin{equation}\label{eq:vstar-pair}
    \bigl(2 A' + 2(c_{\sigf} - c_{\sigb})\bigr) \bmod q \;\le\; e(d) + e(\Pi_\kappa - d),
  \end{equation}
  and whenever $\Pi_\kappa \le R_k$, summing over $d \in S$ yields
  \begin{equation}\label{eq:vstar-sum}
    \card{S} \cdot \bigl(\bigl(2 A' + 2(c_{\sigf} - c_{\sigb})\bigr) \bmod q\bigr) \;\le\; 2 \sum_{d \in S} e(d),
  \end{equation}
  with \emph{zero dependence on $m$ or on the integer magnitude of $Y$ or $F(d)$}.
\item \emph{(Exact parity-conditioned straddle count.)} Suppose $\Pi_\kappa = 2L$ is even and $2L \le R_k$. For \emph{either}
  fixed parity $p \in \{0, 1\}$, summing $\card{\mathrm{Strad}(c_0, \alpha_V, 2L, R_k)}$ over all centers $c_0 \in
  [\alpha_V, \alpha_V + 4L)$ with $c_0 \bmod 2 = p$ equals \emph{exactly} $L^2 = \Pi_\kappa^2 / 4$.
\end{enumerate}
\end{theorem}

\begin{outline}
For~(1) (\lean{straddle_otherHeadInc_sum}), the V decomposition gives $F(d) = (c_{\sigf} - c_{\sigb}) + Y(o_u(d)) -
Y(o_p(d))$ and $F(\Pi_\kappa - d) = (c_{\sigf} - c_{\sigb}) + Y(o_p(d)) - Y(o_u(d))$; adding cancels $Y(o_u(d)) - Y(o_p(d))$
identically in~$\Z$.

For~(2) (\lean{theoremVStar_pair}, \lean{theoremVStar_sum}), $e(d) + e(\Pi_\kappa - d) \ge 0$ is congruent modulo $q$ to
$(A' + F(d)) + (A' + F(\Pi_\kappa - d)) = 2A' + 2(c_{\sigf} - c_{\sigb})$; since for any integer $S_{\mathrm{sum}} \ge 0$ and
$q > 0$ we always have $S_{\mathrm{sum}} \bmod q \le S_{\mathrm{sum}}$,~\eqref{eq:vstar-pair} holds without any wrap-around
loss. Because $d \mapsto \Pi_\kappa - d$ is an involution of $S$ when $\Pi_\kappa \le R_k$, $\sum_{d \in S} (e(d) +
e(\Pi_\kappa - d)) = 2 \sum_{d \in S} e(d)$, giving~\eqref{eq:vstar-sum}.

For~(3) (\lean{sum_straddleSet_card_parity}), swap the order of summation between $c_0 \in [\alpha_V, \alpha_V + 4L)$
($c_0 \bmod 2 = p$) and $d \in \{1, \dots, 2L - 1\}$. For each fixed $d \in \{1, \dots, 2L - 1\}$, solving the inequalities
in~\eqref{eq:straddle-set} for $c_0$ shows that $\{c_0 : d \in \mathrm{Strad}(c_0, \alpha_V, 2L, R_k)\}$ is the half-open
interval $[\alpha_V + \max(d, 2L - d),\, \alpha_V + \max(d, 2L - d) + 2 \min(d, 2L - d))$, whose length $2 \min(d, 2L - d)$
is even. Every half-open integer interval $[a, a + 2M)$ of even length contains \emph{exactly} $M$ even integers and
\emph{exactly} $M$ odd integers (\lean{card_Ico_even_len_mod_two}), so the parity-$p$ fiber has exact size $\min(d, 2L - d)$,
and $\sum_{d=1}^{2L-1} \min(d, 2L - d) = L^2$ (\lean{sum_min_sub_two_mul}).
\end{outline}

\subsection{Super-Fibonacci unwrapping for level-uniform labels}\label{sec:dnp-fib}

On a single head of $U(s_0, m, 1)$ carrying a level-uniform potential $(A, E)$ modulo $q \ge 3$, consider the blind
(uninterrupted) revisit hops around the apices and valleys of level $k$:
\begin{itemize}
\item every blind apex hop at level $k$ has reduced mod-$q$ hop error $a(k) := A_k \bmod q \in [0, q)$
  (\lean{blindApexErr}, verified on \lean{strokePhiRec} by \lean{strokePhiRec_apex_hop});
\item for $1 \le k \le l \le m$, every blind valley hop at level $k$ whose half-width stops in the plain segment of nested
  level $l$ (\emph{valley band $l - k$}) has reduced mod-$q$ hop error (\lean{strokePhiRec_valley_band_hop})
  \begin{equation}\label{eq:band-err}
    b_{l - k}(k) = \lean{blindBandErr}(A, E, q, k, l) :=
    \Bigl(E_k - \sum_{i = k + 1}^l (A_i + E_i + 2)\Bigr) \bmod q \in [0, q).
  \end{equation}
\end{itemize}
Whenever a level-$k$ class is $\theta$-served ($\theta \le 1/256$) and at least half of its apex hop family or valley
sub-band $0, 1, 2$ is blind at the chosen shift (by~\cite[Proposition~5.17 and Corollary~5.18]{EighthPaper2026}, fine
classes are mostly blind on average over shifts, whereas at a fixed shift half-blindness is an explicit hypothesis),
Markov's inequality bounds each corresponding reduced blind error by $q / 8$ (\lean{blind_error_le_div_eight}).
Unwrapping~\eqref{eq:band-err} across bands $0, 1, 2$ forces \emph{super-Fibonacci} growth of the blind errors from fine to
coarse scales, unconditionally ruling out $s \ge 2\lfloor\log_2 q\rfloor + 3$ consecutive levels of blind errors at most
$q/8$ (and hence bounding consecutive half-blind $\theta$-served levels conditionally on half-blindness at the chosen shift).

\begin{theorem}[Super-Fibonacci unwrapping for level-uniform labels]\label{thm:LU}
\formal{DNPFibonacci.lean}{strokePhiRec_apex_hop, strokePhiRec_valley_band_hop, blindApexErr, blindBandErr, unwrap_blindBand_step, superFib_step, superFib_lower_bound, log_bound_of_fib_le, blind_error_le_div_eight, theoremLU_consecutive}
Let $A, E : \N \to \Z$ and $q \ge 3$.
\begin{enumerate}
\item \emph{(Exact $\Z$-unwrapping and super-Fibonacci step.)} For any $k \le l$, if
  \lean{blindBandErr}$(A, E, q, k, l+1)$, \lean{blindApexErr}$(A, q, l+1)$, and
  \lean{blindBandErr}$(A, E, q, l+1, l+1)$ are each at most $\lfloor q / 8 \rfloor$ (i.e., $8 e \le q$), then the mod-$q$
  relation lifts to exact equality in~$\Z$:
  \begin{align*}
    &\lean{blindBandErr}(A, E, q, k, l) - \lean{blindBandErr}(A, E, q, k, l+1) \\
    &\qquad = \lean{blindApexErr}(A, q, l+1) + \lean{blindBandErr}(A, E, q, l+1, l+1) + 2.
  \end{align*}
  Consequently, if $b_1(k), b_2(k), a(k+1), b_0(k+1), a(k+2), b_0(k+2)$ are each at most $q/8$, adding the $(k, k)$ and
  $(k, k+1)$ identities telescopes $b_1(k)$ and yields the super-Fibonacci inequality
  \begin{equation}\label{eq:super-fib}
    b_0(k+1) + b_0(k+2) + 4 \;\le\; b_0(k).
  \end{equation}
\item \emph{(Fibonacci lower bound and $2\lfloor\log_2 q\rfloor + 2$ run limit.)} Suppose $s \ge 1$ consecutive levels
  $k_0, \dots, k_0 + s - 1$ satisfy $8 a(k_0 + i) \le q$ and $8 b_0(k_0 + i) \le q$ for $0 \le i < s$, $8 b_1(k_0 + i) \le q$
  for $i + 1 < s$, and $8 b_2(k_0 + i) \le q$ for $i + 2 < s$. Writing $F_s = \texttt{Nat.fib}(s)$ ($F_0 = 0, F_1 = F_2 = 1,
  F_3 = 2, \dots$), we have
  \[
    4(F_s - 1) \le b_0(k_0), \qquad 32(F_s - 1) \le q, \qquad\text{and}\qquad s \le 2 \lfloor \log_2 q \rfloor + 2.
  \]
\end{enumerate}
\end{theorem}

\begin{outline}
For~(1) (\lean{unwrap_blindBand_step}), write $\mathrm{ex}$, $\mathrm{ey}$, $\mathrm{ea}$, and $\mathrm{ee}$ for
\lean{blindBandErr} at $(A, E, q, k, l)$, \lean{blindBandErr} at $(A, E, q, k, l+1)$, \lean{blindApexErr} at $(A, q, l+1)$,
and \lean{blindBandErr} at $(A, E, q, l+1, l+1)$, and let $D = (\mathrm{ex} - \mathrm{ey}) - (\mathrm{ea} + \mathrm{ee} + 2)$.
By~\eqref{eq:band-err}, $q \mid D$. Nonnegativity gives $D \le \mathrm{ex} - 2 < q$. For the lower bound, if $q \ge 8$, the
hypotheses $8 \cdot \mathrm{ey} \le q$, $8 \cdot \mathrm{ea} \le q$, and $8 \cdot \mathrm{ee} \le q$ give $\mathrm{ey} +
\mathrm{ea} + \mathrm{ee} + 2 \le 3q/8 + 2 < q$, while if $3 \le q < 8$, $8 x \le q$ forces $\mathrm{ey} = \mathrm{ea} =
\mathrm{ee} = 0$ and hence $D = \mathrm{ex} - 2 \ge -2 > -q$. Thus $\abs{D} < q$ and $q \mid D$, forcing $D = 0$. Applying
the step lemma \lean{unwrap_blindBand_step} at $l = k$ (using $8 b_1(k) \le q$) and at $l = k + 1$ (using $8 b_2(k) \le q$) sums to
$b_0(k) - b_2(k) = a(k+1) + b_0(k+1) + a(k+2) + b_0(k+2) + 4$; dropping the nonnegative terms $b_2(k), a(k+1), a(k+2) \ge 0$
yields~\eqref{eq:super-fib} (\lean{superFib_step}).

For~(2) (\lean{superFib_lower_bound}, \lean{log_bound_of_fib_le}, \lean{theoremLU_consecutive}), strong induction on $s$
shows $b_0(k_0) \ge 4(F_s - 1)$ since $4(F_{s-1} - 1) + 4(F_{s-2} - 1) + 4 = 4(F_s - 1)$. Combining with $8 b_0(k_0) \le q$
gives $32(F_s - 1) \le q$. Using $F_{2k+3} \ge 2^k + 1$ (\lean{two_pow_add_one_le_fib}) for $s \ge 3$ ($k = \lfloor (s-3)/2
\rfloor$) gives $2^{k+5} \le 32(F_s - 1) \le q$, so $k + 5 \le \lfloor \log_2 q \rfloor$ and $s \le 2k + 4 \le 2\lfloor
\log_2 q \rfloor - 6 < 2\lfloor \log_2 q \rfloor + 2$.
\end{outline}

\subsection{Block-cancellation drift, Mersenne clocks, and bounded-range block pairs}\label{sec:dnp-block}

To evade \cref{thm:LU} over more than $O(\log q)$ levels on a single walk,~\cite[Section~5.3]{EighthPaper2026} introduced
the periodic \emph{block-cancellation family} $B(a, K)$ ($a \ge 1$, $K \in \Z$): partition $\{0, \dots, m\}$ from the bottom
level~$m$ upward into blocks of length $a + 1$, designating $\kappa$ as a \emph{canceller} if $(m - \kappa) \bmod (a + 1)
= a$ (\lean{isCanceller}) and as a \emph{served} level with $j = (m - \kappa) \bmod (a + 1) \in \{0, \dots, a-1\}$ served
levels below it otherwise. Setting $Q(a, K) = (2^a - 1)(2K + 2)$ (\lean{blockQ}), define (\lean{blockA}, \lean{blockE}):
\[
  A_\kappa = \begin{cases} 0 & \text{if }\kappa\text{ is a canceller},\\ K & \text{if }\kappa\text{ is served},\end{cases}
  \qquad
  E_\kappa = \begin{cases} -Q(a, K) - 2 & \text{if }\kappa\text{ is a canceller},\\
    K + (2^j - 1)(2K + 2) & \text{if }\kappa\text{ is served.}\end{cases}
\]
On served levels (\lean{block_served_mass}), $A_\kappa + E_\kappa + 2 = 2^j(2K + 2)$, $c_{\sigb}(\kappa) = 1 + K$, and
$c_{\sigf}(\kappa) = 2^j(2K + 2)$; on cancellers (\lean{block_canceller_mass}), $c_{\sigb}(\kappa) = 1$, $c_{\sigf}(\kappa)
= -Q(a, K)$, and $c_{\sigf}(\kappa) - c_{\sigb}(\kappa) = -Q(a, K) - 1$, resetting the cumulative sum $\sum (A_i + E_i + 2)$
across each $(a+1)$-level block to~$0$. Moreover, reducing $K \mapsto K \bmod q$ leaves every mod-$q$ blind error invariant
(\lean{block_blind_errors_mod_K}).

We first record the exact path-sum drift laws of $B(a, K)$ and the Mersenne-modulus behavior of the pure binary clock
$B(\infty, 0)$ ($c_{\sigb} = 1$, $c_{\sigf}(\kappa) = 2^{m - \kappa + 1}$, $A \equiv 0$, which makes every single-head hop
error $0$ in $\Z$ before modular reduction, and modulo a Mersenne number $q = 2^B - 1$ makes single-head hop errors vanish
periodically in~$B$, yet incurs straddle residue $q - 1$):

\begin{proposition}[One-sided block drift and Mersenne straddle failure]\label{prop:OC}
\formal{DNPConjectureApp.lean}{blockQ, blockA, blockE, block_served_mass, block_canceller_mass, block_blind_errors_mod_K, propOC_fwd_bounds, propOC_back_bounds, propOC_sum_all_paths, propOC_mersenne_straddle}
Let $a \ge 1$ and $K \ge 0$, and for branch choices $(d_{\mathrm{canc}}, d_{\mathrm{srv}}) \in \{\sigb, \sigf\} \times
\{\sigb, \sigf\}^a$ across one complete $(a+1)$-level block define the block path sum \lean{blockPathSum}$(a, K,
d_{\mathrm{canc}}, d_{\mathrm{srv}})$.
\begin{enumerate}
\item \emph{(Forward vs.\ back canceller bounds.)} When $d_{\mathrm{canc}} = \sigf$, the block path sum lies in
  $[-Q(a, K) + a(1 + K),\, 0]$ and equals $0$ when $d_{\mathrm{srv}} \equiv \sigf$. When $d_{\mathrm{canc}} = \sigb$, the
  block path sum lies in $[1 + a(1 + K),\, Q(a, K) + 1]$.
\item \emph{(Positive mean block drift.)} For every $K \in \Z$, summing \lean{blockPathSum}$(a, K, p_1, p_2)$ over all
  $2^{a+1}$ branch choices $p \in \{\sigb, \sigf\} \times \{\sigb, \sigf\}^a$ equals $2^a\bigl(a(K + 1) + 1\bigr)$, which
  is strictly positive for all $K \ge 0$ (mean drift $+\tfrac{1}{2}(a(K+1) + 1)$ per $(a+1)$-level block under the fair-coin
  law of \cref{prop:SL}(2)).
\item \emph{(Mersenne binary-clock straddle residue.)} For any $B \ge 2$ and Mersenne modulus $q = 2^B - 1$, the pure
  binary clock ($c_{\sigb} = 1$, $c_{\sigf}(\kappa) = 2^{m - \kappa + 1}$, $A' = 0$) satisfies, at every level $\kappa \le m$
  with $B \mid (m - \kappa + 2)$,
  \[
    \bigl(2 A' + 2(c_{\sigf}(\kappa) - c_{\sigb}(\kappa))\bigr) \bmod q = q - 1.
  \]
\end{enumerate}
\end{proposition}

\begin{outline}
Parts~(1) and~(2) (\lean{propOC_fwd_bounds}, \lean{propOC_back_bounds}, and \lean{propOC_sum_all_paths}, proved in Lean for
all $a \ge 0$) use the geometric identity $\sum_{j=0}^{a-1} 2^j = 2^a - 1$ (\lean{sum_fin_two_pow}) and coordinatewise
involution $f_j \mapsto \neg f_j$ on $\{\sigb, \sigf\}^a$: summing $d_{\mathrm{canc}} \in \{\sigb, \sigf\}$ gives
$1 - Q(a, K)$, while each served coordinate $j$ averages $\tfrac{1}{2}(2^j(2K + 2) + 1 + K)$, summing to $Q(a, K) +
a(1 + K)$ and leaving net mean $\tfrac{1}{2}(1 + a(1 + K))$. For~(3) (\lean{propOC_mersenne_straddle}), $B \mid (m -
\kappa + 2)$ implies $(2^B - 1) \mid (2^{m - \kappa + 2} - 1)$, so $2(2^{m - \kappa + 1} - 1) = (2^{m - \kappa + 2} - 1) -
1 \equiv -1 \equiv q - 1 \pmod q$.
\end{outline}

We now assemble \cref{prop:SL}, \cref{thm:Vstar}, and \cref{thm:LU} with prefix lockstep and simultaneous reverse-Markov
shift selection to establish three-regime failure witnesses on $\DNP(s_0, m)$ for all block-label pairs
$(B(a_1, K_1), B(a_2, K_2))$ of bounded block length. For $1 \le k \le \kappa \le m$, write
\lean{periodVApices}$(s_0, m, \kappa) \subseteq [0, P)$ for the set of $2^{\kappa + 1}$ level-$\kappa$ V
start times in one period, \lean{oddRange}$(P) = \{\xi \in [0, P) : \xi \bmod 2 = 1\}$, \lean{cancellerBelow}$(m, a, k)$
for the first canceller level $\kappa \in [k + \lfloor \log_2(m+1) \rfloor + 1,\, k + \lfloor \log_2(m+1) \rfloor + 1 + a]$
of $B(a, \cdot)$ satisfying $\Pi_\kappa \le R_k$, and \lean{classStraddleCount}$(s_0, m, k, \kappa, \mathrm{dir}, \xi)$ for
the total number of straddling half-widths $d$ across all $2^{k+1}$ apices of level $k$ on the active head ($\mathrm{dir} =
\mathrm{true}$ for head~$1$ at $(c + \xi) \bmod P$, $\mathrm{dir} = \mathrm{false}$ for head~$2$ at $(c + P - \xi) \bmod P$)
against all level-$\kappa$ V's of the other head.

\begin{theorem}[Three-regime witnesses for bounded-range block pairs on $\DNP$]\label{thm:A}
\formal{DNPConjectureApp.lean}{apexOffset_diff_indep_prefix, periodVApices_spans_disjoint, card_periodVApices, even_of_mem_periodVApices, sum_classStraddleCount_odd, reverse_markov_simultaneous, classStraddleCount, cancellerBelow_props, exists_oddShift_coarse_straddle, block_straddle_residue_ge_half_q, block_overflow_band_witness, theoremA_three_regimes, theoremA_symmetric_blockPair}
Let $a_1, a_2 \ge 1$, $m < s_0$, and $4(\lfloor \log_2(m+1) \rfloor + 1 + \max(a_1, a_2)) \le m + 1$. Set $\Omega =
{}$\lean{oddRange}$(P)$ and $K_{\mathrm{coarse}} = \{1, \dots, \lfloor (m+1)/4 \rfloor\}$. There exist \emph{two fixed odd
shifts} $\xi_1, \xi_2 \in \Omega$---chosen once and for all from $(s_0, m, a_1, a_2)$, \emph{independent of $K_1, K_2$,
and $q$} (with a \emph{single} odd shift $\xi \in \Omega$, chosen solely from $(s_0, m, a_1)$, sufficing whenever $a_1 = a_2$
and $K_1 = K_2$)---such that for \emph{every} $K_1, K_2 \in \Z$ and \emph{every} modulus $q \ge 3$, setting $\tilde{Q}_h =
Q(a_h, K_h \bmod q)$, at least one of the following three failure regimes occurs:
\begin{enumerate}
\item \emph{(Small-$q$ super-Fibonacci regime, $3 \le q < 32$.)} For every $k_0 \in \N$, the window $[k_0, k_0 + 11)$ of
  $11$ consecutive levels on head~$1$ contains at least one level whose reduced blind apex error or reduced blind band-$0,
  1, 2$ error strictly exceeds $q / 8$ (realized as walk valley-band hop errors whenever $1 \le k_0$ and $k_0 + 10 \le m$).
\item \emph{(Large-$\tilde{Q}$ blind fine-overflow regime, $q \ge 32$ and $q < 4\tilde{Q}_h + 4$ for some $h \in \{1, 2\}$.)}
  On head~$h$, there is a served residue $\mathrm{rem} < a_h$ such that \emph{every} level $k < m$ with $(m - k) \bmod
  (a_h + 1) = \mathrm{rem}$ has reduced blind apex error $\ge q/16$, or reduced blind band-$0$ or band-$1$ error $\ge q/8$.
\item \emph{(Small-$\tilde{Q}$ coarse V-straddle regime, $q \ge 32$ and $4\tilde{Q}_1 + 4 \le q$, $4\tilde{Q}_2 + 4 \le q$.)}
  \begin{itemize}
  \item If $\tilde{Q}_1 \le \tilde{Q}_2$, then at odd shift $\xi_1$, at least $\tfrac{1}{16}\card{K_{\mathrm{coarse}}}$
    coarse levels $k \in K_{\mathrm{coarse}}$ of head~$1$ simultaneously satisfy, against $\kappa =
    {}$\lean{cancellerBelow}$(m, a_2, k)$ of head~$2$,
    \begin{align*}
      \lean{classStraddleCount}(s_0, m, k, \kappa, \mathrm{true}, \xi_1) &\;\ge\; \frac{2^{k+2} R_k}{2^{a_2 + 7}}, \\
      \bigl(2 A^{(1)}_k + 2(c_{\sigf}^{(2)}(\kappa) - c_{\sigb}^{(2)}(\kappa))\bigr) \bmod q &\;\ge\; \frac{q}{2}.
    \end{align*}
  \item Symmetrically, if $\tilde{Q}_2 < \tilde{Q}_1$, then at odd shift $\xi_2$, at least
    $\tfrac{1}{16}\card{K_{\mathrm{coarse}}}$ coarse levels $k \in K_{\mathrm{coarse}}$ of head~$2$ satisfy the analogous
    bounds with $a_1$ and $(B(a_2, K_2), B(a_1, K_1))$.
  \end{itemize}
\end{enumerate}
\end{theorem}

In both subcases of Regime~3, combining the straddle count and residue lower bounds of \cref{thm:A}(3) with the per-V
straddle inequality~\eqref{eq:vstar-sum} of \cref{thm:Vstar}(2) shows that at every such coarse level $k \in K_{\mathrm{coarse}}$,
the straddling apex hops accumulate total reduced mod-$q$ hop error at least $\Omega(2^{-a} q \cdot 2^k R_k)$---a constant
fraction $2^{-O(a)}$ of the entire class capacity $q \cdot 2^{k+1} R_k$.

\begin{outline}
We highlight the four formal components in \leanfile{DNPConjectureApp.lean} (which also establishes prefix lockstep
$\alpha_\kappa(\eps) - \alpha_k(\eps)$ independent of $(\eps_1, \dots, \eps_{k-1})$, \lean{apexOffset_diff_indep_prefix}):
\begin{enumerate}
\item \emph{Exact odd-shift straddle sum (\lean{sum_classStraddleCount_odd}).} The $2^{\kappa+1}$ level-$\kappa$ V-spans
  $[\alpha_V, \alpha_V + 2\Pi_\kappa)$ in $[0, P)$ are pairwise disjoint (\lean{periodVApices_spans_disjoint},
  \lean{card_periodVApices}) and have even start times when $m < s_0$ (\lean{even_of_mem_periodVApices}). For each coarse
  apex $c$ of level $k$, as $\xi$ ranges over the $P/2$ odd shifts \lean{oddRange}$(P)$, $(c + \xi) \bmod P$ (and likewise
  $(c + P - \xi) \bmod P$) bijects onto the $P/2$ elements of $[0, P)$ of parity opposite to~$c$. Summing over $\alpha_V$
  and applying \cref{thm:Vstar}(3) on each $[\alpha_V, \alpha_V + 2\Pi_\kappa)$ yields the exact identity
  \[
    \sum_{\xi \in \lean{oddRange}(P)} \lean{classStraddleCount}(s_0, m, k, \kappa, \mathrm{dir}, \xi)
    = 2^{k+1} \cdot 2^{\kappa+1} \cdot (\Pi_\kappa / 2)^2.
  \]
\item \emph{Simultaneous reverse-Markov selection (\lean{reverse_markov_simultaneous} and \lean{exists_oddShift_coarse_straddle}).}
  For every $k \in K_{\mathrm{coarse}}$ and $\kappa = \mathrm{cancellerBelow}(m, a, k)$ (\lean{cancellerBelow},
  \lean{cancellerBelow_props}), we have $\Pi_\kappa \le R_k \le 2^{a+2} \Pi_\kappa$ and $m - \kappa + 1 \ge \tfrac{1}{2}(m+1)$,
  so letting $X(\xi, k)$ denote \lean{classStraddleCount} at $(s_0, m, k, \kappa, \mathrm{dir}, \xi)$ gives
  $X(\xi, k) \le M_k := 2^{k+1}\Pi_\kappa$ (because at most one V-span contains any shifted center,
  \lean{totalStraddleAtCenter_le}), while its mean over $\Omega = \lean{oddRange}(P)$ ($\card{\Omega} = 2(m+1)2^{s_0}$) satisfies $2^{\kappa+1}(\Pi_\kappa/2)^2 / \card{\Omega} =
  \tfrac{m - \kappa + 1}{4(m+1)} \Pi_\kappa \ge \tfrac{1}{8}\Pi_\kappa$, giving $\mathbb{E}_{\xi \in \Omega}[X(\xi, k)] \ge
  \tfrac{1}{8} M_k$. Applying reverse Markov per level $k$ (\lean{reverse_markov_single}) gives $\Pr_{\xi \in \Omega}[X(\xi, k)
  \ge \tfrac{1}{16} M_k] \ge \tfrac{1}{16}$, and double-counting over $K_{\mathrm{coarse}}$ (\lean{reverse_markov_simultaneous})
  exhibits a \emph{single} odd shift $\xi \in \Omega$ at which at least $\tfrac{1}{16}\card{K_{\mathrm{coarse}}}$ coarse
  levels satisfy $X(\xi, k) \ge \tfrac{1}{16} M_k \ge 2^{k+2} R_k / 2^{a+7}$.
\item \emph{Mod-$q$ straddle residue without wrap-around (\lean{block_straddle_residue_ge_half_q}).} Suppose $q \ge 32$,
  $\tilde{Q}_1 \le \tilde{Q}_2$, and $4\tilde{Q}_2 + 4 \le q$. Let $\tilde{K}_h = K_h \bmod q \in [0, q)$. At any level $k$
  of $B(a_1, \tilde{K}_1)$, $A^{(1)}_k \in \{0, \tilde{K}_1\}$, while $2\tilde{K}_1 + 2 \le \tilde{Q}_1 \le \tilde{Q}_2$
  (since $a_1 \ge 1$). At canceller $\kappa$ of $B(a_2, \tilde{K}_2)$, $c_{\sigf}^{(2)}(\kappa) - c_{\sigb}^{(2)}(\kappa) =
  -\tilde{Q}_2 - 1$. Therefore the integer $Z := 2 A^{(1)}_k + 2(-\tilde{Q}_2 - 1)$ satisfies
  \[
    -\frac{q}{2} \;\le\; -2\tilde{Q}_2 - 2 \;\le\; Z \;\le\; 2\tilde{K}_1 - 2\tilde{Q}_2 - 2 \;\le\; -4 \;<\; 0.
  \]
  Because $Z \in [-q/2, -1]$, its Euclidean remainder $Z \bmod q = Z + q$ lies in $[q/2, q - 4]$, hence $\ge q/2$.
\item \emph{Fine-scale doubling overflow (\lean{block_overflow_band_witness}).} Suppose $q \ge 32$ and $4\tilde{Q}_h + 4 > q$.
  Set $\tilde{K} = K_h \bmod q \in [0, q)$. If $16\tilde{K} \ge q$, every served level ($\mathrm{rem} = 0 < a_h$) has
  $a(k) = \tilde{K} \ge q/16$. If $16\tilde{K} < q$, consider the doubling sequence
  $T_r = (2^r - 1)(2\tilde{K} + 2)$ for $0 \le r \le a_h$: we have $T_0 = 0$, $T_{r+1} = 2 T_r + (2\tilde{K} + 2)$, and
  $4 T_{a_h} + 4 = 4\tilde{Q}_h + 4 > q$. Since $2\tilde{K} + 2 \le q/8 + 2 < q/4$ for $q \ge 32$, the first index $r \le
  a_h$ where $\tilde{K} + T_r \ge q/8$ cannot jump past $7q/8$, landing strictly inside $[q/8, 7q/8] \subseteq [q/8, q - 1]$.
  When $r < a_h$, $\tilde{K} + T_r = E_k \bmod q = b_0(k)$ on residue $\mathrm{rem} = r$; when $r = a_h$, $\tilde{K} +
  T_{a_h} = (E_k - (A_{k+1} + E_{k+1} + 2)) \bmod q = b_1(k)$ across the canceller $k + 1$ on residue $\mathrm{rem} = 0$.
\end{enumerate}
Combining Regimes~1 (from \cref{thm:LU}(2) with $s = 11 > 2\lfloor\log_2 31\rfloor + 2 = 10$),~2, and~3 proves
\lean{theoremA_three_regimes} and \lean{theoremA_symmetric_blockPair}.
\end{outline}

\begin{remark}[Scope relative to Conjecture~A$''$]\label{rem:scope-A}
\Cref{thm:A} gives three-regime failure witnesses for block-label pairs $(B(a_1, K_1), B(a_2, K_2))$ of bounded block
length across all $K_1, K_2 \in \Z$ and all $q \ge 3$ with \emph{zero logarithmic loss in $q$}, at two fixed odd shifts
(one shift for symmetric pairs $a_1 = a_2, K_1 = K_2$). Deriving Conjecture~A$''$~\cite[Section~5.6, Conjecture~5.24]{EighthPaper2026}
for block pairs---and ultimately for all low-cost labels on the $r$-fold cover---requires four further steps:
\begin{enumerate}[label=(\alph*)]
\item \emph{One common odd shift for asymmetric pairs.} Conjecture~A$''$ fixes one odd shift $\xi$ defining the pair
  $x_2 = x_1(\cdot + \xi)$ before quantifying over labels and~$q$. In Regime~3 of \cref{thm:A}, $\xi_1$ is used when
  $\tilde{Q}_1 \le \tilde{Q}_2$ and $\xi_2$ when $\tilde{Q}_2 < \tilde{Q}_1$, which depends on $(K_1, K_2, q)$; applying
  first-moment reverse Markov separately to $\mathrm{dir} = \mathrm{true}$ and $\mathrm{dir} = \mathrm{false}$ gives two
  subsets of $\Omega$ of density at least $\tfrac{1}{16}$ that need not intersect without a joint second-moment
  (carry-decorrelation) bound across \lean{periodVApices}.
\item \emph{Half-blindness at the fixed shift.} Regimes~1 and~2 bound \emph{blind} single-head errors (\lean{blindApexErr}
  and \lean{blindBandErr}, where $F = 0$), which coincide with actual two-head hop errors only on hops whose window contains
  no turn of the other head. In~\cite[Proposition~5.17 and Corollary~5.18]{EighthPaper2026}, blindness of fine classes is
  proved on average over shifts rather than at the specific shift chosen for Regime~3.
\item \emph{From hop-error witnesses to doubly $\theta_0$-served level counts.} Conjecture~A$''$ bounds the number of doubly
  $\theta_0$-served levels by $C\log_2 q + C' + \eta(m+1)$ ($\eta < 1$). Passing from \cref{thm:A}'s witnesses to that
  level-count bound requires combining~(b) on fine levels with the explicit class-fraction weights of the apex and
  valley-band hop families in Regimes~1--3 to determine $\theta_0(a_0)$ and $\eta < 1$.
\item \emph{Unbounded block lengths and general low-cost labels.} Beyond bounded block lengths $a_1, a_2 \le a_0$, when
  $\log_2 q \ll m \ll q$ and $a \asymp \tfrac{1}{2}\log_2 q$, the single-canceller straddle density $2^{-a} \asymp q^{-1/2}$
  drops below constant $\theta_0$; and full Conjecture~A$''$ covers \emph{every} label of cost at most $n_{\mathrm{h}}/8$ on
  every $r$-fold cover (not only turn-supported or level-uniform block potentials, and subject to period closure modulo
  $rP$).
\end{enumerate}
\end{remark}

\begin{remark}[Finite numerical illustration]\label{rem:dnp-numerics}
The standalone script \leanfile{tools/reproduce_numerics.py} evaluates exact stroke counts, odd-center straddle sums,
super-Fibonacci band unwrapping, exhaustive odd-shift straddle sums (e.g., on $s_0 = 7, m = 6$), and all three regimes of
\cref{thm:A} on finite parameters (e.g., $s_0 = 24, m = 23, a_1 = a_2 = 1$, where $4(\lfloor\log_2 24\rfloor + 1 + 1) =
24 = m + 1$ and $K_{\mathrm{coarse}} = \{1, \dots, 6\}$), confirming that the explicit odd-shift reverse-Markov witness
and mod-$q$ straddle residues match the formal bounds across all tested $(K_1, K_2, q)$.
\end{remark}

\section{What a linear-time transposer would need}\label{sec:construction}

Our aim across this series is to decide~(T), whether by constructing a linear-time transposer or by proving that none
exists. Impossibility theorems and candidate constructions inform one another directly: every conditional or restricted
half-rate bound and label-cut barrier rules out an entire mechanism on which a linear-time machine cannot rely, narrowing
the design space to the exact structural regimes where impossibility is not yet certified. This section collects what the
results of \cref{sec:reciprocity,sec:crux,sec:dnp} require of any linear-time transposer (oblivious unless stated
otherwise) and identifies the specific mechanisms that fail. It introduces no new numbered theorems: every claim is one of
the formalized results above, a cited result of an earlier paper in the series, or explicitly marked as a heuristic.

\subsection{The depth--coding interface}\label{sec:con-interface}

For an oblivious transposer~$M$ with $k$ work tapes, value width $\beta$, and running time $T \le C(L + 1)$ at input
length $L = d^2 b(d)$ ($b(d) = \max(1, \clog{d})$), Theorem~D of the fourth paper~\cite[Section~4.2, Theorem~4.10]{FourthPaper2026}
(proved in the text of~\cite{FourthPaper2026}) presents a dichotomy for every $\eps \in (0, 1]$ and $D \ge 0$:
\begin{itemize}[leftmargin=2em]
\item \emph{Either} the dependency DAG $H_M$ admits no genie of cost at most $\eps L$ and depth at most~$D$ (defeating
  cheap depth reduction);
\item \emph{or} the revealed word of such a genie forms an index code of at most $\eps L$ bits for the bounded-depth cone
  side information $N_I$ on $[0, L)$, and its largest fibre is a binary storage code $C \subseteq \{0,1\}^L$ of rate at
  least $1 - \eps$ on $N_I$~\cite[Remark~4.9(b)]{FourthPaper2026}.
\end{itemize}
By the fifth paper's paired-distance analysis~\cite[Theorem~3.3 and Corollaries~3.7 and~3.8]{DirectedRecovery2026}, for
every fixed $\delta > 0$ and odd radius $r$, deleting a set $Z$ of at most $\delta L + O(L / \log d)$ coordinates
(high-reuse positions together with short paired-incidence paths, measured by the short-pair count $b_E(r)$ in the
bipartite incidence graph of $N_I$) leaves a bounded-degree relation $N_I - Z$ with $b_E(r) = 0$; at $r = 3$, $N_I - Z$
is loopless and TT-free~\cite[Lemma~2.1]{SeventhPaper2026} (\lean{HalfRate.ttFree_iff_shortPaired_eq_empty}). Moreover, by
paid coordinate deletion (\cite[Corollary~6.9]{SeventhPaper2026}, Lean-verified as
\lean{HalfRate.card_sq_le_of_paidDeletion}), if $N_I - Z$ satisfies the polymatroid half-rate bound~\eqref{eq:half-rate},
then $\card{C}^2 \le 2^{L + \card{Z}}$, bounding the rate of $C$ by $\tfrac{1}{2} + \card{Z}/(2L) \le \tfrac{1}{2} +
\tfrac{\delta}{2} + o(1)$ and ruling out rate $1 - \eps$ whenever $\eps < \tfrac{1}{2} - \tfrac{\delta}{2}$.

\subsection{Necessary conditions on the cone codes}\label{sec:con-codes}

Suppose an oblivious linear-time transposer admits a bounded-depth genie of cost $\eps L$ ($\eps < \tfrac{1}{2}$),
yielding a binary storage code $C$ of rate at least $1 - \eps$ on $N_I$, and let $N' = N_I - Z$ be the loopless TT-free
sub-relation obtained after paid deletion of $\card{Z} \le \delta L + o(L)$ coordinates. In addition to the conditions
established in~\cite{FourthPaper2026,DirectedRecovery2026,SeventhPaper2026,EighthPaper2026}, our coding theorems impose
two sharp structural requirements on $N'$:

\begin{enumerate}[label=(N\arabic*),leftmargin=3em]
\item\label{it:con-N1} \emph{Positive density of unreciprocated full-in-degree vertices ($\neg\ClassS$).} For \emph{every}
  subset $W$ of coordinates of $N'$ on which the induced relation $N'[W]$ belongs to $\ClassS$, we must have
  \[
    \card{W} \;\le\; (2\eps + \delta)L + o(L).
  \]
  Indeed, applying \cref{thm:Sstar}(2) (\lean{rateBound_classS}) to $N'[W]$ and paid deletion to the remaining
  $L - \card{W}$ coordinates of $N_I$ gives $\card{C}^2 \le 2^{L + (L - \card{W})} = 2^{2L - \card{W}}$, whereas
  $\card{C} \ge 2^{(1 - \eps)L}$ forces $2(1 - \eps)L \le 2L - \card{W}$, i.e., $\card{W} \le 2\eps L$ (and at most
  $(2\eps + \delta)L + o(L)$ inside $N'$). Likewise, by \cref{cor:Bstar}, on any $\ClassS$ sub-relation no vertex cut $S$
  can create an isolated-parent surplus $\card{I_{N'[W]}(S)} - \card{S}$ without correspondingly shrinking code capacity.

  \emph{Design consequence:} When $N'$ has maximum in-degree~$3$, a coding gain above rate~$\tfrac{1}{2}$ (whether linear
  or nonlinear, over any alphabet) cannot come from subcubic undirected exchange, pure one-way directed cycles of
  in-degree-$2$ vertices, chordal one-way shortcuts across digon cycles, or even stride-$3$ hub cycles such as those of
  $\Gto$ (\cref{prop:G21}) whose hubs have in-degree~$2$. Rather, a linear fraction of coordinates in strongly connected
  components of $N'$ must have \emph{in-degree~$3$ together with at least one unreciprocated outgoing arc} ($v \in N'(u)$
  with $u \notin N'(v)$), or else have in-degree at least~$4$.

\item\label{it:con-N2} \emph{Distance-$2$ heavy clusters for degree-$4$ scalar linear codes.} Suppose the machine's data
  path is linear over a field $F$ and a fibre code on a triangle-free symmetric support $G = (V, E)$ of maximum degree~$4$
  is realized as $\ker M$ for a unit-diagonal check matrix $M \in F^{V \times V}$. By \cref{thm:SHstar} and
  \cref{cor:Mstar} (\lean{finrank_ker_le_matching_rhoMatch}),
  \[
    \dim_F \ker M \;\le\; \nu(G) + \card{H_M} - \rhom(G, H_M) \;\le\; \frac{\card{V}}{2} + \card{H_M} - \rhom(G, H_M),
  \]
  and the same bound holds with $H_{M^\T}$ in place of $H_M$. To achieve $\dim_F \ker M \ge (1 - \eps)\card{V}$, the
  code must satisfy
  \[
    \card{H_M} - \rhom(G, H_M) \;\ge\; \Bigl(\frac{1}{2} - \eps\Bigr)\card{V}, \qquad
    \card{H_{M^\T}} - \rhom(G, H_{M^\T}) \;\ge\; \Bigl(\frac{1}{2} - \eps\Bigr)\card{V}.
  \]
  In~\cite{EighthPaper2026}, where $\rho_3(G, H_M)$ stood in place of $\rhom(G, H_M)$, adjacent pairs of degree-$4$ rows
  remained a plausible source of linear excess (as on the bipartite family $G_{8k}$ of \cref{rem:deg4-numerics}, where
  $\card{H} - \rho_3(G_{8k}, H) = k = \tfrac{1}{8}\card{V}$). \Cref{cor:Mstar} closes that loophole completely: isolated
  heavy edges ($u \sim v$ with no shared neighbor, since $G$ is triangle-free) contribute full weight~$2$ to
  $\rhom(G, H_M)$ and yield \emph{zero} rate gain above~$\tfrac{1}{2}$. Any degree-$4$ scalar linear gain must concentrate
  a linear fraction of heavy vertices at pairwise distance \emph{exactly~$2$} in~$G$, sharing common neighbors.
\end{enumerate}

\subsection{Necessary conditions on head schedules and cut labels}\label{sec:con-walks}

By Theorems~F and~M of~\cite{EighthPaper2026}, any oblivious linear-time transposer that defeats cheap depth reduction must
employ unbounded owner-load on at least two work heads, or on one work head coupled with bounded-load heads across
unboundedly many separated scales; the $2$-head deep-nesting pair $\DNP(s_0, m)$ is the minimal prototype of the two-head
regime. Our results in \cref{sec:dnp} sharpen both sides of the $\DNP(s_0, m)$ frontier:

\begin{itemize}[leftmargin=2em]
\item \emph{Why simple turn-label cuts fail to prove $\TCP(2)$ on $\DNP(s_0, m)$.} An impossibility proof via $\TCP(2)$
  must exhibit, for every $\eps > 0$, a bounded modulus $q = O_\eps(1)$ and a low-cost label $\Psi : [0, T) \to \Z_q$.
  \Cref{thm:LU} proves that no level-uniform turn potential can keep blind apex and band-$0,1,2$ errors at most $q/8$ across
  more than $2\lfloor\log_2 q\rfloor + 2$ consecutive levels without periodic negative resets; \cref{prop:OC}(2)--(3) proves
  that periodic $(a+1)$-block resets carry strictly positive mean path drift $+\tfrac{1}{2}(a(K+1)+1)$ across random nested
  branches while pure binary clocks modulo Mersenne numbers $q = 2^B - 1$ suffer $q - 1$ straddle error whenever
  $B \mid (m - \kappa + 2)$; and \cref{thm:A} exhibits three-regime blind-error and coarse-straddle witnesses on $\Omega(m)$
  levels at two fixed odd shifts $\xi_1, \xi_2$ (a single odd shift for symmetric pairs) for \emph{every} bounded-range block
  pair $(B(a_1, K_1), B(a_2, K_2))$ across \emph{all} $K_1, K_2 \in \Z$ and \emph{all} $q \ge 3$. Thus any cut refuting
  depth-robustness of $\DNP(s_0, m)$ at misaligned odd shifts where fine bands are half-blind must exploit mechanisms
  outside \cref{thm:A} (\cref{rem:scope-A}), such as block lengths $a$ growing with $\log_2 q$, non-level-uniform jumps
  varying across sibling V's at the same scale, or cuts not supported at turns.
\item \emph{What a depth-robust $2$-head schedule must do (heuristic reading).} Conversely, if a $2$-head oblivious
  transposer seeks to evade cheap depth cuts through deep nesting, aligned shifts ($\xi = 1$) fail because both heads
  share a single doubled block potential~\cite[Proposition~5.14]{EighthPaper2026}, whereas misaligned odd shifts create the
  simultaneous multi-scale V-straddle interference quantified by \cref{thm:Vstar} and \cref{thm:A}: coarse apex hops of
  scale $R_k$ straddle canceller valleys of scale $\Pi_\kappa \asymp R_k$ on the other head with exact parity-invariant
  weight $\Pi_\kappa^2 / 4$ per V (\cref{thm:Vstar}(3)). Even if such a schedule achieves depth-robustness on its deep
  core, Theorem~F of~\cite{EighthPaper2026} restricts that core to $o(L)$ heavily visited scratch cells unless the machine
  also overcomes the Scratch Question (\cite[Section~10, Problem~10.5]{EighthPaper2026})---or, for adaptive machines,
  escapes the complete source-cover and return-profile barriers of~\cite{NinthPaper2026}.
\end{itemize}

\section{Open problems}\label{sec:open}

\subsection{Updated map of half-rate and deep-nesting bounds}

For oblivious machines, Theorem~D~\cite[Section~4.2, Theorem~4.10]{FourthPaper2026} reduces $\neg$(T) to depth reduction
($\TCP(k)$) together with asymptotic coding bounds on paired-distance-separated recovery relations
(\cite[Question~9.1]{DirectedRecovery2026}). The table below updates the status of exact radius-$3$ (loopless TT-free)
half-rate bounds from~\cite[Section~9.1]{SeventhPaper2026} and~\cite[Section~8]{EighthPaper2026}:

\begin{center}
\small
\begin{tabular}{@{}>{\raggedright\arraybackslash}p{0.31\textwidth}>{\raggedright\arraybackslash}p{0.33\textwidth}>{\raggedright\arraybackslash}p{0.30\textwidth}@{}}
\toprule
\textbf{regime} & \textbf{proved} & \textbf{open nearby} \\
\midrule
directed in-degree $\le 3$, full reciprocal closure at in-degree~$3$ ($\ClassS$) &
  all codes, all alphabets: piece decomposition and $\card{C}^2 \le \card{\alpha}^{\card{V}}$ (\textbf{Theorem~S$^*$}, \cref{thm:Sstar}) &
  general TT-free in-degree $\le 3$ (\cref{prob:conjP}; Conjecture~P of~\cite{SeventhPaper2026}) \\
\addlinespace
chordless digon cycles with $\le 1$ outside parent per vertex &
  piece certificate exists for some $m \in \N$ iff not a stride-$3$ hub cycle (\textbf{Theorem~DC}, \cref{thm:DC}); $\Ex$ core and $\GP(5,2)$ ruled out for all $m \in \N$ (\cref{cor:x15,prop:GP}) &
  chorded digon cycles with one-way chords (Conjecture~K of~\cite{SeventhPaper2026}); $O(1)$ record bound on $\ClassS$ (\cref{prob:records}) \\
\addlinespace
triangle-free degree $\le 4$, scalar linear ($\ker M$, unit diagonal) &
  $\dim_F \ker M \le \lfloor \card{V}/2 \rfloor$ under~(S1) alone (heavy matchings allowed; \textbf{Theorem~SH$^*$}, \cref{thm:SHstar}), plus $\rhom$ barrier (\textbf{Corollary~M$^*$}, \cref{cor:Mstar}) &
  heavy vertices at distance~$2$ sharing common neighbors; all codes at degree~$4$ (\cref{prob:deg4-dist2}) \\
\addlinespace
$\DNP(s_0, m)$, turn-supported mod-$q$ labels &
  blind-error run $\le 2\lfloor\log_2 q\rfloor + 2$ (\textbf{Theorem~LU}, \cref{thm:LU}); three-regime failure witnesses for bounded-range block pairs at two odd shifts, one for symmetric pairs (\textbf{Theorem~A}, \cref{thm:A}) &
  common asymmetric shift, fixed-shift blindness, served-level conversion, $a \asymp \tfrac{1}{2}\log_2 q$, and general $\TCP(2)$ (\cref{prob:App-full}) \\
\bottomrule
\end{tabular}
\end{center}

\subsection{Open problems}

\begin{enumerate}[label=\arabic*.,leftmargin=2em]
\item \emph{Decide (T).} Construct a deterministic multitape Turing machine that transposes $d \times d$ matrices of
  $b(d)$-bit entries ($b(d) = \max(1, \clog{d})$) in time $O(d^2 b(d))$, or prove that none exists (checked against
  \lean{Transpose.LinearTimeTransposable} in~\cite{FirstPaper2026}). On adaptive machines, extending the complete
  source-cover and return-profile barriers of~\cite{NinthPaper2026} beyond restricted return geometries remains the central
  non-oblivious barrier question.

\item \emph{Conjecture~P beyond $\ClassS$, and Conjecture~K.} \Cref{thm:Sstar} resolves the intermediate target named in
  item~2 of~\cite[Section~9.2]{SeventhPaper2026} (the all-code version of the fourth paper's Theorem~S on its whole class).
  Conjecture~P (\cite[Conjecture~9.1]{SeventhPaper2026}) remains open; the gap to the full directed degree-$3$ case lies
  in vertices of in-degree~$3$ that send unreciprocated outgoing arcs:
  \begin{problem}[Directed in-degree three without reciprocal closure]\label{prob:conjP}
  Does every finite loopless TT-free recovery relation with $\card{N(v)} \le 3$ for all $v \in V$ admit a certified piece
  decomposition (\lean{HalfRate.Decomp}, Conjecture~P of~\cite[Section~9.2]{SeventhPaper2026}), or at least satisfy the
  polymatroid half-rate bound $2 h(V) \le \sum_{v \in V} h(\{v\})$? As a concrete special case (Conjecture~K
  of~\cite[Section~9.2]{SeventhPaper2026}), does every loopless TT-free chorded digon cycle with at most one incoming and at
  most one outgoing one-way chord per vertex admit a global elimination certificate?
  \end{problem}
  By~\cite[Proposition~8.1]{SeventhPaper2026} (\lean{HalfRate.dn3_rate_eventually_le}), an affirmative answer to
  \cref{prob:conjP} implies an affirmative answer to the paired-distance coding question~\cite[Question~9.1]{DirectedRecovery2026}
  for all recovery relations of in-degree at most~$3$, with no restriction on out-degree. On $\Ex \notin \ClassS$, vertex~$6$
  has in-degree~$3$ and sends a pure one-way arc $6 \to 3$, obstructing its minimum core $\{0,\dots,5\}$ across all $m \in \N$
  (\cref{cor:x15}), yet $\Ex$ still admits a $3$-record global certificate~\cite[Remark~7.3]{SeventhPaper2026}.

\item \emph{Record complexity of piece decompositions in $\ClassS$.} In \cref{prop:GP} we showed that global elimination
  certificates exhibit a strict record hierarchy ($m = \infty$ on $\GP(5,2)$, $m = 4$ on $\GP(7,2)$, and $m = 3$ on
  $\GP(9,2)$). For piece decompositions in $\ClassS$, pure stride-$s$ cores ($s \ge 4$) use $m \le 3$ records
  (\cref{prop:CS}(2)), whereas in \cref{thm:DC} and \cref{thm:Sstar} our certificates on non-stride-$3$ chordless digon
  $L$-cycles ($L \ge 5$), minimal chordal arcs of span $d \ge 3$ (Subcase~4a), and hub-augmented stride-$3$ $g$-cycles
  (Subcase~4c) use up to $L + 1$, $d + 2$, and $g + 2$ records, respectively.
  \begin{problem}[Bounded-record decompositions for $\ClassS$]\label{prob:records}
  Is there an absolute constant $m_0$ (for instance $m_0 = 3$ or $m_0 = 4$) such that every finite relation $N \in \ClassS$
  admits a certified piece decomposition in which every piece certificate uses at most $m_0$ records?
  \end{problem}

\item \emph{Triangle-free storage codes at degree~$4$ with distance-$2$ heavy clusters.} \Cref{thm:SHstar} and
  \cref{cor:Mstar} show that for scalar linear codes $\ker M$ on triangle-free graphs of maximum degree~$4$, heavy vertices
  need not be non-adjacent: only shared neighbors between distinct heavy vertices ($\dist_G(u, v) = 2$, violating~(S1)) can
  break the minimal-counterexample run.
  \begin{problem}[Distance-$2$ heavy clusters and all-code degree-$4$ bounds]\label{prob:deg4-dist2}
  (a) Over an arbitrary field~$F$, does $\dim_F \ker M \le \lfloor \card{V}/2 \rfloor$ hold for every unit-diagonal matrix
  $M$ supported on a finite triangle-free graph $G$ of maximum degree~$4$, when $H_M$ contains vertices at distance~$2$
  in~$G$---or at least up to $o(\card{V})$ along locally bipartite families of growing odd girth?
  (b) Does the polymatroid half-rate bound $2 h(V) \le \sum_{v \in V} h(\{v\})$ hold on every triangle-free undirected graph
  of maximum degree~$4$?
  \end{problem}
  Note that Barg and Z\'emor's high-rate binary linear codes on triangle-free $4$-regular Cayley graphs~\cite{BargZemor2022}
  employ multi-bit vector symbols per vertex rather than scalar unit-diagonal checks, and at degree~$5$ scalar binary codes
  already reach rate~$5/8$ on the Clebsch graph~\cite[Corollary~6.11]{FourthPaper2026}.

\item \emph{Closing the block-pair gaps to Conjecture~A$''$, and $\TCP(2)$ on the deep-nesting pair.} \Cref{thm:A} gives
  three-regime failure witnesses for block-label pairs $(B(a_1, K_1), B(a_2, K_2))$ of bounded block length on
  $\DNP(s_0, m)$ across all $K_1, K_2 \in \Z$ and $q \ge 3$ at two fixed odd shifts (one shift for symmetric pairs). On the
  $r$-fold periodic cover of $\DNP(s_0, m)$ (containing $n_{\mathrm{h}} = 2 r P = 8 r (m+1) 2^{s_0}$ long hops across the
  two heads), the four steps of \cref{rem:scope-A} remain to establish Conjecture~A$''$~\cite[Section~5.6, Conjecture~5.24]{EighthPaper2026}
  and decide $\TCP(2)$:
  \begin{problem}[From Theorem~A to Conjecture~A$''$ and walk-DAG depth of $\DNP$]\label{prob:App-full}
  (a) Prove the existence of a \emph{single} odd shift $\xi \in \Omega$ that works simultaneously for both coarse straddle
  directions ($\mathrm{dir} = \mathrm{true}$ and $\mathrm{dir} = \mathrm{false}$) and satisfies half-blindness on fine
  sub-bands, and deduce from \cref{thm:A} an explicit upper bound $C\log_2 q + C' + \eta(m+1)$ ($\eta < 1$) on doubly
  $\theta_0$-served levels for all bounded-range block pairs.
  (b) Extend the bound when $q$ is exponential in~$m$ and $\max(a_1, a_2) \asymp \tfrac{1}{2}\log_2 q$, and from block pairs
  to \emph{every} mod-$q$ label of cost at most $n_{\mathrm{h}}/8$ on the $r$-fold cover (full Conjecture~A$''$).
  (c) Are the walk DAGs of $\DNP(s_0, m)$ depth-robust as $m \to \infty$ at some sequence of odd shifts, refuting $\TCP(2)$
  (and hence $4$-stack depth reduction $\mathrm{SR}(4)$; compare $\mathrm{SR}(3)$ and the Scratch Question
  in~\cite[Section~10, Problems~10.2 and~10.5]{EighthPaper2026})?
  \end{problem}
\end{enumerate}


\appendix
\section{Index of Lean declarations}\label{app:index}

This index lists every declaration of this paper's library cited in the text, relative to the namespace
\texttt{DirectedReciprocity}. The declarations named in status lines are grouped by the module named there; those named only
in the text are listed at the end. All of them depend only on the axioms \texttt{propext}, \texttt{Classical.choice} and
\texttt{Quot.sound} (\cref{app:verify}). Declarations of earlier libraries of the series that are named in the text,
in the namespaces \texttt{Transpose}, \texttt{MotionDepthCodes}, \texttt{DirectedRecovery}, \texttt{HalfRate},
\texttt{BoundedNesting} and \texttt{SimpleGraph} (\texttt{lean-machine}), are defined and verified in those
libraries~\cite{FirstPaper2026,FourthPaper2026,DirectedRecovery2026,SeventhPaper2026,EighthPaper2026}; all except the
standalone \texttt{Transpose} package of~\cite{FirstPaper2026} are also re-checked directly in \leanfile{lean/PaperAxioms.lean}.

\paragraph{\texttt{DirectedReciprocity.ClassS}}
{\small\raggedright \texttt{ClassSOn}, \texttt{IsClassS}, \texttt{barrier\_ineq\_classS}, \texttt{card\_sq\_le\_classS}, \texttt{classSOn\_decomp}, \texttt{classS\_decomp}, \texttt{exists\_inClosed\_scc}, \texttt{exists\_piece\_of\_scc\_classS}, \texttt{isolatedOutDir}, \texttt{liftPieceCertInClosed}, \texttt{log\_card\_le\_barrier\_classS}, \texttt{rateBound\_classS}.\par}

\paragraph{\texttt{DirectedReciprocity.ClassSExamples}}
{\small\raggedright \texttt{g21Rel}, \texttt{g21\_hex\_no\_pieceCert}, \texttt{g21\_isClassS}, \texttt{g21\_min\_indeg\_and\_acyclic\_W}, \texttt{g21\_no\_any\_triangle}, \texttt{g21\_shortest\_digon\_cycles}, \texttt{g21\_stronglyConnected}, \texttt{rateBound\_g21}, \texttt{x15\_hexagon\_no\_pieceCert}.\par}

\paragraph{\texttt{DirectedReciprocity.DNPConjectureApp}}
{\small\raggedright \texttt{apexOffset\_diff\_indep\_prefix}, \texttt{blockA}, \texttt{blockE}, \texttt{blockQ}, \texttt{block\_blind\_errors\_mod\_K}, \texttt{block\_canceller\_mass}, \texttt{block\_overflow\_band\_witness}, \texttt{block\_served\_mass}, \texttt{block\_straddle\_residue\_ge\_half\_q}, \texttt{cancellerBelow\_props}, \texttt{card\_periodVApices}, \texttt{classStraddleCount}, \texttt{even\_of\_mem\_periodVApices}, \texttt{exists\_oddShift\_coarse\_straddle}, \texttt{periodVApices\_spans\_disjoint}, \texttt{propOC\_back\_bounds}, \texttt{propOC\_fwd\_bounds}, \texttt{propOC\_mersenne\_straddle}, \texttt{propOC\_sum\_all\_paths}, \texttt{reverse\_markov\_simultaneous}, \texttt{sum\_classStraddleCount\_odd}, \texttt{theoremA\_symmetric\_blockPair}, \texttt{theoremA\_three\_regimes}.\par}

\paragraph{\texttt{DirectedReciprocity.DNPFibonacci}}
{\small\raggedright \texttt{blindApexErr}, \texttt{blindBandErr}, \texttt{blind\_error\_le\_div\_eight}, \texttt{log\_bound\_of\_fib\_le}, \texttt{strokePhiRec\_apex\_hop}, \texttt{strokePhiRec\_valley\_band\_hop}, \texttt{superFib\_lower\_bound}, \texttt{superFib\_step}, \texttt{theoremLU\_consecutive}, \texttt{unwrap\_blindBand\_step}.\par}

\paragraph{\texttt{DirectedReciprocity.DNPStraddle}}
{\small\raggedright \texttt{HasLevelUniformVDecomp}, \texttt{mem\_straddleSet\_sub\_iff}, \texttt{straddleSet}, \texttt{straddle\_otherHeadInc\_sum}, \texttt{strokePhiRec\_hasLevelUniformVDecomp}, \texttt{sum\_straddleSet\_card\_parity}, \texttt{theoremVStar\_pair}, \texttt{theoremVStar\_sum}.\par}

\paragraph{\texttt{DirectedReciprocity.DNPStrokes}}
{\small\raggedright \texttt{card\_stroke\_stopDepth}, \texttt{card\_stroke\_stopDepth\_branchWord}, \texttt{card\_stroke\_stopDepth\_branchWord\_parity}, \texttt{card\_stroke\_stopDepth\_parity}, \texttt{decompDepth\_valid}, \texttt{decomp\_reconstruct}, \texttt{dnpPeriod}, \texttt{evalPhiRec}, \texttt{evalPhiRec\_eq\_sum\_turnJumps}, \texttt{reconstruct\_decomp}, \texttt{scaleR}, \texttt{strokeLen}, \texttt{strokePhiRec}, \texttt{strokePhiRec\_back}, \texttt{strokePhiRec\_fwd}.\par}

\paragraph{\texttt{DirectedReciprocity.DigonCycle}}
{\small\raggedright \texttt{IsStride3Cycle}, \texttt{digonCycle4\_pieceCert}, \texttt{digonCycleGe5\_pieceCert}, \texttt{digonCycle\_pieceCert\_iff}, \texttt{digonCycle\_pieceCert\_of\_not\_stride3}, \texttt{exists\_non\_stride3\_pos}, \texttt{not\_pieceCert\_of\_isStride3Cycle}.\par}

\paragraph{\texttt{DirectedReciprocity.PetersenHierarchy}}
{\small\raggedright \texttt{generalizedPetersen\_hierarchy}, \texttt{gp52\_eq\_petersenGraph}, \texttt{gp52\_no\_elimCert}, \texttt{gp72\_elimCert4}, \texttt{gp72\_no\_elimCert\_le3}, \texttt{gp92\_elimCert3}, \texttt{gpGraph}, \texttt{nonempty\_elimCert\_iff\_pieceCert\_univ}, \texttt{petersen\_no\_pieceCert\_univ}.\par}

\paragraph{\texttt{DirectedReciprocity.PositionalNest}}
{\small\raggedright \texttt{PieceCert.digon\_ne}, \texttt{PieceCert.mem\_supply\_of\_idxOf\_lt}, \texttt{disjoint\_nested\_iff}, \texttt{relCharged\_iff\_nested}, \texttt{relCharged\_of\_nested}.\par}

\paragraph{\texttt{DirectedReciprocity.StrideCycles}}
{\small\raggedright \texttt{stride3\_no\_pieceCert}, \texttt{strideCore}, \texttt{strideCore3\_no\_pieceCert}, \texttt{strideCore\_pieceCert\_iff}, \texttt{strideEven\_pieceCert\_two}, \texttt{strideGe4\_pieceCert}, \texttt{strideRel}.\par}

\paragraph{\texttt{DirectedReciprocity.TheoremSHStar}}
{\small\raggedright \texttt{F3\_star}, \texttt{MinCE.false\_SHStar}, \texttt{NoCommonNbhd}, \texttt{SHStarOn}, \texttt{exists\_rhoMatch}, \texttt{finrank\_ker\_add\_card\_le\_matching\_SHStar}, \texttt{finrank\_ker\_le\_barrier\_SHStar}, \texttt{finrank\_ker\_le\_matching\_rhoMatch}, \texttt{finrank\_ker\_le\_matching\_rhoMatch\_transpose}, \texttt{finrank\_subCode\_le\_half\_SHStar}, \texttt{odd\_state\_star}, \texttt{rho3\_le\_rhoMatch}, \texttt{rhoMatch}, \texttt{theoremSHStar}, \texttt{theoremSHStar\_graph}, \texttt{theoremSHStar\_transpose}.\par}

\paragraph{Also named in the text}
{\small\raggedright \texttt{SHStarOn.mono}, \texttt{ValidStrokeCode}, \texttt{augmentedStride3\_pieceCert}, \texttt{blockPathSum}, \texttt{cancellerBelow}, \texttt{card\_Ico\_even\_len\_mod\_two}, \texttt{cycRev}, \texttt{dcCol}, \texttt{dcPerm}, \texttt{dcSupply}, \texttt{decompDepth}, \texttt{digonGraph}, \texttt{digonGraph\_cliqueFree3}, \texttt{elimCert\_adj\_ne}, \texttt{isCanceller}, \texttt{minChordArc\_pieceCert}, \texttt{not\_mem\_rev\_empty\_of\_lt\_idxOf}, \texttt{oddRange}, \texttt{otherHeadInc}, \texttt{pathSumBack}, \texttt{pathSumFwd}, \texttt{periodTurnJump}, \texttt{periodVApices}, \texttt{reconstructDepth}, \texttt{relCharged\_step\_nested}, \texttt{relCharged\_step\_none}, \texttt{relCharged\_step\_some}, \texttt{resRel}, \texttt{reverse\_markov\_single}, \texttt{strokeLen\_succ}, \texttt{strokeMass}, \texttt{strokePhiRec\_eq\_sum\_turnJumps}, \texttt{strokeTurnJump}, \texttt{sum\_fin\_two\_pow}, \texttt{sum\_min\_sub\_two\_mul}, \texttt{totalStraddleAtCenter\_le}, \texttt{two\_pow\_add\_one\_le\_fib}.\par}

\section{Building and verifying the formalization}\label{app:verify}

The Lean~4 formalization lives in \leanfile{lean/} next to this paper. It uses Lean and Mathlib at version \texttt{v4.34.1}
(\leanfile{lean/lean-toolchain}, \leanfile{lean/lakefile.toml}, \leanfile{lean/lake-manifest.json}) and builds directly on
the Lean libraries of the seventh paper (\texttt{halfRate},~\cite{SeventhPaper2026}) and the eighth paper
(\texttt{boundedNesting},~\cite{EighthPaper2026}) via relative paths; through them it also imports the libraries of the
fourth, fifth, second, and first papers~\cite{FourthPaper2026,DirectedRecovery2026,SecondPaper2026,FirstPaper2026} together
with the repository's shared library \texttt{lean-machine}. Keep the repository's sibling directory layout intact so that
these relative package paths resolve.

\paragraph{Build.} On a fresh clone of the repository, first fetch the second paper's pinned external dependency as
documented in its README (\texttt{python3 verification/fetch\_openai\_metric.py}, run inside
\leanfile{../Local-Coding-Barriers-and-the-Limits-of-Tape-Geometry-for-Matrix-Transposition/}). Then, with \texttt{elan}
installed, run in \leanfile{lean/}:
\begin{center}
\texttt{lake exe cache get \&\& lake build}
\end{center}
The first command downloads precompiled Mathlib \texttt{.olean} caches; the second compiles all eleven modules of
\leanfile{lean/DirectedReciprocity/} and the root module \leanfile{lean/DirectedReciprocity.lean} at Lean's default
\texttt{maxHeartbeats} setting ($200{,}000$), with zero errors and zero warnings.

\paragraph{Axioms and declaration index.} The script \leanfile{tools/lean_index.py} scans \leanfile{paper/sections/*.tex}
and regenerates both the declaration index of \cref{app:index} (\leanfile{paper/sections/leanindex.tex}) and
\leanfile{lean/PaperAxioms.lean} (pass \texttt{--index-only} to write only the index). Running in \leanfile{lean/}:
\begin{center}
\texttt{lake env lean PaperAxioms.lean}
\end{center}
executes \texttt{\#print axioms} on every declaration of \texttt{DirectedReciprocity}---and every declaration of earlier
libraries---cited in this paper. Every declaration depends at most on the three standard foundational axioms
\texttt{propext}, \texttt{Classical.choice}, and \texttt{Quot.sound}. No file of this library uses \texttt{sorry}, \texttt{admit}, \texttt{native\_decide}, \texttt{unsafe},
or any custom \texttt{axiom}; finite graph and partition checks in \cref{sec:reciprocity} are verified inside Lean's kernel
via \texttt{decide}, \texttt{decide +kernel}, and \texttt{decide +revert}.

\paragraph{Standalone numerical cross-checks.} The portable Python~3 script \leanfile{tools/reproduce_numerics.py} uses
only the Python standard library and independently reproduces the finite numerical illustrations of
\cref{rem:deg4-numerics,rem:dnp-numerics}:
\begin{center}
\texttt{python3 tools/reproduce\_numerics.py}
\end{center}
Specifically, it constructs the bipartite degree-$4$ family $G_{8k}$ ($2 \le k \le 6$) of \cref{rem:deg4-numerics},
verifying exact matching number $\nu(G_{8k}) = 4k$, packing numbers $\rho_3(G_{8k}, H) = k$ and $\rhom(G_{8k}, H) = 2k$,
and explicit $\F_2$ kernel dimension $4k$; and it evaluates stroke-law counts (\cref{prop:SL}), parity-conditioned straddle
sums (\cref{thm:Vstar}), super-Fibonacci band unwrapping (\cref{thm:LU}), block-path mean drift (\cref{prop:OC}), and the
explicit reverse-Markov odd shift and three-regime error witnesses of \cref{thm:A} (\cref{rem:dnp-numerics}). None of the
proofs in the paper depends on this script.

\paragraph{Module guide (\texttt{lean/DirectedReciprocity/}).}
\begin{itemize}[leftmargin=2em]
\item \leanfile{PositionalNest.lean}: stateless positional characterization of \lean{HalfRate.RelCharged}, disjoint-record
  nesting simplification, first-appearance supply forcing, and proper record coloring across digons (\cref{prop:PN}).
\item \leanfile{StrideCycles.lean}: positional locking lemma on stride-$3$ hub cycles, unbounded-$m$ piece impossibility
  when $s = 3$, and $2$-record (even $s \ge 4$) and $3$-record (all $s \ge 4$) co-supplied piece certificates on canonical
  stride-$s$ cores (\cref{prop:CS}).
\item \leanfile{DigonCycle.lean}: universal chordless digon cycle dichotomy ($L \ge 4$), non-stride-$3$ witness extraction,
  and two-sided $(L+1)$-record piece expansion (\cref{thm:DC}).
\item \leanfile{PetersenHierarchy.lean}: equivalence between global \lean{HalfRate.PieceCert} and \lean{HalfRate.ElimCert},
  canonical partition reduction and unbounded-$m$ global impossibility on $\GP(5,2)$ (Petersen graph), $m \le 3$
  impossibility and $m = 4$ certificate on $\GP(7,2)$, and $m = 3$ certificate on $\GP(9,2)$ (\cref{prop:GP}).
\item \leanfile{ClassS.lean}: definition of $\ClassS$ and relative \lean{ClassSOn}, condensation to an in-closed strongly
  connected component, minimal chordal one-way arc certificates, shortest digon cycle and $(g+2)$-record sequential hub
  augmentation, Theorem~S$^*$ (\cref{thm:Sstar}), and the directed barrier bound Corollary~B$^*$ (\cref{cor:Bstar}).
\item \leanfile{ClassSExamples.lean}: unbounded-$m$ core impossibility on $\Ex$ (\cref{cor:x15}), and construction and
  complete cycle/peeling analysis of the $21$-vertex strongly connected triangle-free witness $\Gto \in \ClassS$
  (\cref{prop:G21}).
\item \leanfile{TheoremSHStar.lean}: unified final count \lean{F3_star} under~(S1) alone, odd-state lemma \lean{odd_state_star},
  minimal-counterexample exclusion removing~(S2), Theorem~SH$^*$ in matrix, transpose, and graph forms (\cref{thm:SHstar}),
  and the common-neighbor-free packing invariant \lean{rhoMatch} and Tutte--Berge deficiency bound Corollary~M$^*$
  (\cref{cor:Mstar}).
\item \leanfile{DNPStrokes.lean}: dyadic scales, stroke lengths, recursive stroke decomposition bijection, stopping-depth
  and parity-independence counts, and $O(m)$ recursive turn potential \lean{evalPhiRec} (\cref{prop:SL}).
\item \leanfile{DNPStraddle.lean}: straddle reflection involution, $\Z$-cancellation of nested stroke potentials, $m$-free
  mod-$q$ pair and set bounds, and exact parity-conditioned straddle count $\Pi_\kappa^2 / 4$ (\cref{thm:Vstar}).
\item \leanfile{DNPFibonacci.lean}: blind apex and valley-band error formulas on \lean{strokePhiRec}, exact $\Z$-unwrapping
  under $\le q/8$ bounds, super-Fibonacci step $b_0(k) \ge b_0(k+1) + b_0(k+2) + 4$, and the $2\lfloor\log_2 q\rfloor + 2$
  consecutive-level bound (\cref{thm:LU}).
\item \leanfile{DNPConjectureApp.lean}: periodic block-cancellation family $B(a, K)$, forward/backward path bounds, positive
  mean drift, Mersenne binary-clock straddle failure (\cref{prop:OC}), prefix lockstep of nested V-apices, simultaneous
  reverse-Markov odd-shift selection, fine-scale doubling overflow, and the three-regime completeness Theorem~A
  (\cref{thm:A}).
\end{itemize}


\bibliographystyle{plainurl}
\bibliography{references}

\end{document}
